\documentclass[11pt]{article}
\usepackage[margin=1in]{geometry}
\usepackage{amsmath,amssymb,amsthm}
\usepackage{array,longtable,graphicx}
\usepackage[colorlinks=true,linkcolor=blue,citecolor=blue,urlcolor=blue]{hyperref}
\hypersetup{pdftitle={Unknown-frame recovery in sparse block-regular graphs},pdfauthor={Research manuscript}}
\newtheorem{theorem}{Theorem}[section]
\newtheorem{lemma}[theorem]{Lemma}
\newtheorem{proposition}[theorem]{Proposition}
\newtheorem{corollary}[theorem]{Corollary}
\newtheorem{definition}[theorem]{Definition}
\theoremstyle{remark}
\newtheorem{remark}[theorem]{Remark}
\newcommand{\one}{\mathbf 1}
\newcommand{\Prob}{\mathbb P}
\newcommand{\E}{\mathbb E}
\newcommand{\Fp}{\mathbb F_p}
\newcommand{\F}{\mathbb F}
\newcommand{\R}{\mathbb R}
\newcommand{\Q}{\mathbb Q}
\newcommand{\Z}{\mathbb Z}
\newcommand{\rank}{\operatorname{rank}}
\newcommand{\nullity}{\operatorname{nullity}}
\newcommand{\supp}{\operatorname{supp}}
\newcommand{\Var}{\operatorname{Var}}
\newcommand{\spec}{\operatorname{spec}}
\newcommand{\TV}{\operatorname{TV}}
\newcommand{\KL}{D_{\mathrm{KL}}}
\newcommand{\im}{\operatorname{im}}
\newcommand{\nnz}{\operatorname{nnz}}
\newcommand{\diag}{\operatorname{diag}}
\newcommand{\eps}{\varepsilon}
\renewcommand{\labelitemi}{$\bullet$}
\title{Unknown-frame recovery in sparse block-regular graphs\\Moments, certified algorithms, and sharp ambiguity}
\author{Research manuscript}
\date{November 15, 2026}
\begin{document}
\maketitle
\begin{abstract}
We recover an unknown bounded-size degree frame and its planted classes
in sparse block-regular graphs, from the graph and fixed order and degree
bounds alone. Independent uniform simple internal and cross factors may
have unequal class sizes, varying total degrees, and absent internal
edges. The hypothesis is a symmetric omission certificate with weights
in $[0,1/2)$ and unit physical row sums; zero weights are permitted.
The plant is, with high probability, the unique finest bounded-class
equitable partition. A bounded raw finite-field channel moment gives
complete partition listing in unconditional expected nearly quadratic
bit work and expected nearly linear space. Its proof uses heterogeneous
root/output spaces and selected cancellation-edge certificates to handle
inactive quotient channels. A certified complete-kernel algorithm
charges every input before averaging. We also prove disjoint old and
new spectra for all fixed connected simple regular base lifts, including
bipartite bases and irrational old modes. Exact two-class ambiguity and
hidden-prior Bayes error constants cover every $a\ge0$, $b\ge1$,
$a+b\ge3$. Complete bases give sharp lift ambiguity, while cycle lifts
have optimal success tending to zero. Separate prescribed-cross
spectral theorems yield rational recovery under their stated internal
degree hypotheses. Direct cover-recovery algorithms and classical
counting, mixing and linear-algebra inputs are credited; historical
priority is not certified.
\end{abstract}
\tableofcontents
\section{Main results and relation to earlier work}\label{sec:overview}
An equitable partition records the exact number of neighbors a vertex
has in each class. Its indicator $J$ satisfies $AJ=JQ$, where $Q$ is
the degree frame. We recover the finest partition with a fixed bound
on its class count. Neither that frame, its order, nor its class sizes
are supplied. The bound is essential, since the singleton partition
is always equitable.

The model has independent uniform simple regular internal factors and
biregular cross factors on a fixed connected reversible frame. Sizes
may be unequal and total degrees may vary. Our sufficient condition is
the feasibility of the symmetric linear certificate
\[
 q_{ij}=q_{ji},\qquad 0\le q_{ij}<\tfrac12,\qquad
 \sum_jd_{ij}q_{ij}=1.
\]
The weights belong to physical neighbor slots; an internal weight
has no extra half factor. They need not be supplied to the algorithm.
Every common total degree $d\ge3$ is covered by $q_{ij}=1/d$.
The condition also covers every connected frame with at least two types
and internal degrees at least two, and further mixed-degree frames with
no internal edges.
Some frames require zero weights; those are included.

Theorem~\ref{thm:unknown-frame} proves that every bounded-class
equitable partition coarsens the plant, with failure
$O(n^{-(d_{\min}-2)})$. The raw finite-field moment in
Theorem~\ref{thm:frame-bounded-moment} and the certified kernel algorithm
give complete partition listing in unconditional expected
$\widetilde O(n^2)$ bit work, expected $\widetilde O(n)$ workspace,
and expected $O(n)$ output bits. The algorithm receives only fixed
class-count and degree bounds. It retains a partition with the largest
class count and recovers the frame from integer neighbor counts.
A deterministic polynomial truncation preserves the structural rate.

Independent lifts of every fixed connected simple regular base of
degree at least three are a special case, for every positive lift
size. All old eigenvalues also retain their exact base multiplicities
with high probability, including bipartite bases and irrational
old modes (Theorem~\ref{lift:thm:all-base}).
Complete bases attain the structural ambiguity power, with explicit
leading constant. The fixed-motif occurrence coefficients have classical
precursors in Friedman--Kohler~\cite{FKMotifs}; the argument here also
classifies the alternatives contributing to recovery error.
In contrast, unconditioned cycle lifts have optimal
hidden-fiber recovery success tending to zero even with their base
supplied.

The two-class results cover every fixed $a\ge0,b\ge1,a+b\ge3$,
including assortative, equal-degree and disassortative choices.
We compute exact leading ambiguity and hidden-prior Bayes constants.
Separate theorems allow prescribed cross factors and identify fixed
quotient eigenspaces, giving exact rational row clustering. Those
robust prescribed-cross quantifiers have their own internal-degree
hypotheses; they are not inferred from the random-cross theorem.

Here $\widetilde O$ suppresses logarithmic bit factors. All degrees,
frame orders, fields, target dimensions and class bounds are fixed.
An always-correct Las Vegas listing algorithm terminates almost surely.
Recovering the particular hidden plant is a separate statistical event.
The complete-kernel certificate is deterministic to verify.

\paragraph{Direct cover-recovery precedents.}
Fellows and Koblitz explicitly discussed recovery of the four hidden
fibers in random covers of $K_4$ in their Kid Krypto
construction~\cite{KidKrypto}. Bernstein, Cottaar and Lange's
2023 evaluation of IPCC already applies echelon reduction to the
public equations $(A+I)v=\one$ and groups identical affine linear
forms to expose those fibers~\cite{BCL}. Their reported experiments
include both exact four-fiber recovery and localized ambiguities.
This is a direct ancestor of the algebraic fiber-recovery method,
and its use of Gaussian elimination and coordinate grouping is not
claimed as originating here. Our statements concern the explicitly
proved fixed-parameter ranges, raw moments, statistical constants
and certified expected-cost composition.

\paragraph{Other recovery and spectral precedents.}
Brito, Dumitriu, Ganguly, Hoffman and Tran proved rigidity, recovery
under an adjacency outlier condition, and bootstrapping in the
two-class regular model~\cite{BDGHT}. Barucca proposed full recovery
in equitable graphs beyond that regime, with numerical
evidence~\cite{Barucca}. Karmakar and Podder proved multicommunity
uniqueness for sufficiently large positive degree entries under
comparability hypotheses~\cite{KP}; their scope is not restricted
to assortative frames. Brito, Dumitriu and Harris give clustering
for supplied nonsingular frames under degree scaling and eigenvalue
separation~\cite{BDH}. Zhu and Zhu's Lemma~9.1 already states
planted-eigenvalue simplicity in the random-cross two-class
model~\cite{ZZ}, so no first-simplicity claim is made.
A norm bound alone does not exclude an exact collision inside the
limiting bulk. Exact inherited eigenvalues already make noiseless
lift detection easy; identifying the finest hidden fibers differs
from that detection task and its noisy versions~\cite{KY}.

\paragraph{Counting, moments, and computation precedents.}
The graph-cover multinomial enumerator and its entropy interpretation
are classical~\cite{Vontobel}. Gaussian and lattice prefactors for
positive interior phases have close precedents in
Mori--Tanaka~\cite{MoriTanaka,MoriTanakaCentral}.
Huang's thesis and Huang--Vontobel's loop-calculus argument also give
uniform bounds on normalized raw cover expectations under a loop-weight dominance
condition~\cite{HuangThesis,HV}.
These results are credited rather than treating cover prefactors or
bounded raw expectations as new methods. Multitype star/edge entropy
also appears in the Gaussian-wave framework of
Dembo--McKenzie~\cite{DM}.

M\'esz\'aros proves stronger fixed-group mixing in global
permutation/matching regular ensembles~\cite{Meszaros}; vector
alphabets there encode arbitrary matrix targets.
Nguyen--Wood develop cokernel moments and modular-to-real
transfer~\cite{NW}, and Huang combines configuration counts, local
limit estimates and large deviations for sparse
invertibility~\cite{Huang}. Our proof classifies the actual boundary
phases and controls their leakage with conditional switching and
rank comparison. The sparse algebra uses Wiedemann sequences and
established preconditioners~\cite{Wiedemann,Chen}, with full per-input
certification and an explicit averaging argument.

These comparisons delimit the results proved here. They do not
certify historical priority. No noise theorem, growing-parameter
uniformity, simplicity of the new spectrum, or numerical-stability
claim is implicit.


\section{Recovery of an unknown degree frame}\label{sec:unknown-frame}
Fix supplied constants $K\ge1,D\ge3$. The unknown fixed nonnegative
integer frame $Q=(d_{ij})$ has $k\le K$ types, total degrees
$d_i=\sum_jd_{ij}\le D$, and connected off-diagonal support, with the
one-type convention. Its positive planted sizes satisfy
\begin{equation}\label{eq:unknown-balance}
 N_id_{ij}=N_jd_{ji},\qquad N_id_{ii}\text{ even},\qquad n=\sum_iN_i.
\end{equation}
Sample all present internal regular and cross biregular factors
independently and uniformly simple. All limits use admissible sizes
where those simple factors exist.

Our hypothesis is a symmetric omission certificate on every present
factor:
\begin{equation}\label{eq:omission-certificate}
 q_{ij}=q_{ji},\qquad 0\le q_{ij}<\tfrac12,\qquad
 \sum_jd_{ij}q_{ij}=1\quad\text{for every }i.
\end{equation}
An internal weight is used once per physical slot. Zero weights are
allowed. The certificate is a hypothesis on the true frame and is
not supplied to the algorithm. Each vertex has at least three
positive-weight slots, so $d_{\min}:=\min_i d_i\ge3$.
Put $c=\min_{d_{ij}>0}(1-2q_{ij})>0$.

Let $J$ be the planted indicator, so $AJ=JQ$. A candidate partition
with $r$ nonempty classes is a one-hot matrix
$Y\in\{0,1\}^{n\times r}$ with nonzero columns. It is equitable if
$AY=YR$ for an integer nonnegative quotient $R$.
A compatible planted coarsening has $Y=JT$ for an onto one-hot map
$T$, with $QT=TR$.

\begin{theorem}[Unknown-frame recovery]\label{thm:unknown-frame}
Under \eqref{eq:omission-certificate}, with failure
$O_{K,D}(n^{-(d_{\min}-2)})$, every equitable partition into at most
$K$ nonempty classes coarsens the plant. The plant is the unique
finest such partition, up to names.

Given only $A,K,D$, an always-correct Las Vegas algorithm lists all
such partitions in unconditional expected $\widetilde O_{K,D}(n^2)$
bit work and expected $\widetilde O_{K,D}(n)$ working bits, with
expected $O_{K,D}(n)$ output bits. Selecting the largest class count
recovers the plant, $k$, its sizes and $Q$ on the preceding event.
A deterministic truncation uses bounded
$\widetilde O_{K,D}(n^3+n^{d_{\min}})$ bit work and $O_K(n^2)$ space,
with the same recovery failure order.
\end{theorem}

Connectedness and balance compare adjacent class sizes by ratios of
integers at most $D$. Therefore
\begin{equation}\label{eq:unknown-rho}
 \alpha_i:=N_i/n\ge\rho:=1/(KD^{K-1}).
\end{equation}
Only finitely many frames occur at fixed $K,D$. Choose one certificate
for each qualifying frame; their finitely many positive $c$ values
have a positive minimum. This makes the stated constants uniform.

\subsection{The weighted entropy and an explicit one-hot gap}
Temporarily use independent distinguished-stub configurations, allowing
loops and repeated edges. At old type $i$ an ordered star records
its root candidate label and the candidate label at every physical
neighbor slot, satisfying the original integer equation
\[
 \sum_{s=1}^{d_i}e_{Z_{is}}=R_{X_i,*}.
\]
Its exact expected type count is \eqref{eq:fmp-exact-types}.
The vertex multinomials already assign all root labels.
At fixed $K,D$ there are polynomially many types, and uniform
zero-safe factorial estimates give exponent $n\Phi$ with only a
polynomial remainder, where
\[
 \Phi=\sum_i\alpha_iH(P_i)-\sum_{i\le j}w_{ij}H(\nu_{ij}),\qquad
 w_{ii}=\alpha_id_{ii}/2,\quad w_{ij}=\alpha_id_{ij}\ (i<j).
\]
Absent terms are omitted. Symmetrize slots aimed at the same old
type. This preserves legality, root laws $\mu_i$ and pooled pair
laws $\nu_{ij}$, while increasing star entropy.

For the symmetric laws let $I_{ij}$ be pair mutual information.
After omission of physical slot $t$ let $J_{it}$ be the retained
conditional total correlation given the root. Every omission is
lossless. Average its entropy chain rule with the slot weights
$q_{it}$; their sum is one. Pair reversal, balance and symmetric
weights cancel all root entropies, giving
\begin{equation}\label{eq:unknown-entropy}
 -\Phi=\sum_{i\le j}w_{ij}(1-2q_{ij})I_{ij}
             +\sum_i\alpha_i\sum_tq_{it}J_{it}\ge0.
\end{equation}
For the original unsymmetrized exponent this is an upper bound.
No positive lower bound on $q$ and no field inverse is used.

Set
\[
 \epsilon=\sum_i\alpha_i(1-\max_x\mu_i(x)),
\]
the normalized distance from the closest arbitrary old-type point map.
For omission $(i,t)$, keep the root law but draw the retained neighbors
independently from their unconditional old-type marginals. The
divergence of the actual retained law from this product law is
\[
 \kappa_{it}=J_{it}+\sum_{s\ne t}I(X_i;Z_{is}).
\]
Writing $L=-\Phi$ and using \eqref{eq:unknown-entropy}, weighted
summation gives
\begin{equation}\label{eq:unknown-KL}
 A_q:=\sum_i\alpha_i\sum_tq_{it}\kappa_{it}
 =J_q+\sum_{i\le j}2w_{ij}(1-q_{ij})I_{ij}\le2L/c,
\end{equation}
where $J_q=\sum_i\alpha_i\sum_tq_{it}J_{it}$.
Products of marginals are positive on every actual supported tuple,
so these divergences are valid at zero marginals.

Put $\mu_{is}=\E e_{Z_{is}}$, $V_{is}=1-\|\mu_{is}\|^2$,
$S_i=\sum_sV_{is}$, and
$B_i=\E\|R_{X_i,*}-\sum_s\mu_{is}\|^2\ge0$.
For the retained data use
\[
 f_{it}=\left\|R_{X_i,*}-\sum_{s\ne t}e_{Z_{is}}-\mu_{it}\right\|^2.
\]
Its actual expectation is $V_{it}$ and its product-law expectation is
$S_i-V_{it}+B_i$. Thus their difference is
$\Delta_{it}=S_i-2V_{it}+B_i$.
Both vectors inside the norm have nonnegative mass $d_i$, so
$0\le f_{it}\le2D^2$ under both laws. Pinsker gives
$|\Delta_{it}|\le\sqrt2D^2\sqrt{\kappa_{it}}$.
Individual differences may be negative, but
\[
 \sum_tq_{it}\Delta_{it}
 =S_i-2\sum_tq_{it}V_{it}+B_i\ge cS_i.
\]
Weighted Cauchy--Schwarz, whose total weight is
$\sum_i\alpha_i\sum_tq_{it}=1$, gives
\[
 c\sum_i\alpha_iS_i
 \le\sqrt2D^2\sqrt{A_q}\le2D^2\sqrt{L/c}.
\]
By reversal and $d_i\ge3$,
\[
 \sum_i\alpha_iS_i
 =\sum_j\alpha_jd_j(1-\|\mu_j\|^2)
 \ge3\sum_j\alpha_j(1-\max_x\mu_j(x))=3\epsilon.
\]
Consequently every one-hot collection satisfies the strict gap
\begin{equation}\label{eq:unknown-gap}
 \Phi\le-\gamma\epsilon^2,\qquad \gamma=9c^3/(4D^4)>0.
\end{equation}
This argument includes zero omission weights, unused candidate labels,
and internal degrees zero or one.

For each fixed $\eta>0$, the exact finite-alphabet type count and
\eqref{eq:unknown-gap} bound the expected number of candidates at
distance at least $\eta n$ from every point map by
$\operatorname{poly}_{K,D}(n)e^{-\gamma\eta^2n}$.
There are finitely many candidate integer quotients of order at most
$K$ and row sums at most $D$. Conditioning all factors simple costs a
bounded factor, so macroscopic alternatives are exponentially unlikely.
No lower bound on candidate class masses is required.

\subsection{Small row supports and compatibility}
Call a vertex set $S$ bad if no exterior vertex has exactly one edge
into it, counting multiplicity. In the configuration model let
$s=|S|$, let $a$ edges lie inside it and $b$ leave it, reaching
$w$ exterior vertices. With $e=a+b$,
\[
 2a+b\ge d_{\min}s,\quad b\ge2w,\quad
 e-s-w\ge(d_{\min}-2)s/2.
\]
Every present pool has at least $\rho n$ stubs. For
$s\le\rho n/(4D)$, one $e$-pair prescription has probability at most
$(2/(\rho n))^e$. Choose $S$, its exterior vertices, selected exterior
stubs, and pairs in at most
$\binom ns\binom nw2^{Dw}(2Ds)^e$ ways.
Using factorial bounds and $(s/w)^w\le e^{s/e_{\rm Euler}}$ for
$w>0$, with value one at $w=0$, yields
\begin{equation}\label{eq:unknown-support}
 \E\#\{\text{bad }S:|S|=s\}
 \le C_{K,D}^s(s/n)^{\lceil(d_{\min}-2)s/2\rceil}.
\end{equation}
For sufficiently small fixed $\eta$, the corresponding real-power
summands have consecutive ratio at most one half through $\eta n$.
Their sum from two is $O(n^{-(d_{\min}-2)})$.
A singleton is not bad in a simple graph of positive degree.
Positive-probability simplicity conditioning proves this same bound
for any nonempty small bad support in the target law.

On the complementary graph event, for every field and matrix $R$,
\begin{equation}\label{eq:unknown-no-support}
 AV=VR,\quad |\supp_{\rm row}V|\le\eta n\quad\Longrightarrow\quad V=0.
\end{equation}
An exterior vertex with one neighbor in the nonzero row support
would otherwise have a nonzero row in $AV$ and zero in $VR$.
This needs no union bound over entries or target matrices.

Take an integer-verified $Y$ within $\eta n$ rows of $JT$ and choose
$\eta<\rho/[2(D+1)]$. Changed closed neighborhoods occupy at most
$(D+1)\eta n$ vertices. Every old class contains an unchanged closed
neighborhood; evaluating the two integer identities there gives
$QT=TR$. Therefore $V=Y-JT$ satisfies $AV=VR$ and has small row
support, so \eqref{eq:unknown-no-support} forces $Y=JT$.
Together with the macro bound this proves the structural assertion.

\subsection{Exact enumeration and covered examples}
Enumerate every $r\le K$ and every nonnegative integer $r\times r$
matrix with row sums at most $D$, at most
$\sum_{r=1}^K(D+1)^{r^2}$ targets. Solve its finite-field intertwiner
system, enumerate the complete right kernel incrementally, and retain
exactly one-hot matrices with nonempty columns satisfying $AY=YR$
over the integers. Keep a candidate with the largest class count.
The true $(J,Q)$ appears. On the rigidity event every candidate is
$JT$ with onto $T$, so $r\le k$ and equality means a permutation.

Theorem~\ref{thm:frame-bounded-moment} supplies the bounded actual-law
moment. The certified per-input algorithm and expectation calculation
of Section~\ref{sec:vector-algorithm} give the stated expected work,
space and output bounds. For a deterministic cap, use binary dense
elimination and interrupt a kernel exceeding $n^{d_{\min}-1}$,
reading $d_{\min}$ from the graph. The finite-family bounded binary
moment gives interruption $O(n^{-(d_{\min}-1)})$.
Elimination plus Gray enumeration costs
$\widetilde O(n^3+n^{d_{\min}})$; the structural failure dominates
the smaller cap probability. This proves the theorem.

Common degree $d\ge3$ has certificate $q_{ij}=1/d$.
If $k\ge2$ and every internal degree $h_i=d_{ii}\ge2$, let
$b_i=\sum_{j\ne i}d_{ij}>0$, choose $\varepsilon=1/(4D)$ on all
cross factors, and put $q_{ii}=(1-\varepsilon b_i)/h_i$.
All weights are positive and strictly below one half.
Thus both earlier regular and mixed-internal-degree classes are
special cases of the single criterion.

\begin{corollary}[Unknown base fibers]\label{cor:random-lift}
Every fixed connected simple base with a certificate
\eqref{eq:omission-certificate} has its fibers and base recovered
from an independent uniform random lift, given only fixed order and
degree bounds, in expected nearly quadratic bit work.
Every positive lift size is allowed. In particular all connected
simple regular bases of degree at least three are included.
\end{corollary}
\begin{proof}
Use its adjacency matrix as $Q$, equal planted sizes, and one
independent uniform matching for each base edge. These are exactly
the simple factors in the theorem, automatically simple for every $N$.
\end{proof}

This includes mixed-degree bases. For $K_5$ minus one edge, give each
edge from a degree-three vertex weight $1/3$, and each edge between
degree-four vertices weight $1/6$. A forced-zero example is
$Q=\left(\begin{smallmatrix}0&3\\3&1\end{smallmatrix}\right)$ with
equal sizes: $q_{12}=1/3$ forces $q_{22}=0$.
For simple bases, two disjoint copies of $K_{3,3}$ joined by one
bridge have a certificate with weight $1/3$ internally and zero on
the bridge; summing the row equations over one copy's bipartition
forces that zero. Degree two and weights equal to one half are
excluded. The exact cycle obstruction appears in
Section~\ref{sec:lift-boundary}.


\section{Bounded channel moments for omission-certified frames}
\label{ch:sec:channel-proof}\label{sec:frame-moment-proof}
We prove the raw moment bound in a closed heterogeneous channel class.
The omission certificate may have forced zero weights. All channels
are injective, and positive-weight channels are isomorphisms.
The common-space intertwiner and regular-lift statements below are
specializations of this theorem.

\subsection{Model and statements}
Fix finitely many old types $i\le k$, positive integers $a_i$, and nonnegative integers $d_{ij}$ with
\[
 a_id_{ij}=a_jd_{ji},\qquad d_i=\sum_jd_{ij},\qquad D=\max_i d_i<\infty.
\]
At scale $N$, type $i$ has $n_i=a_iN$ vertices; put $n=\sum_i n_i$ and $\alpha_i=n_i/n$. Each type-$i$ vertex has $d_{ij}$ labeled physical stubs aimed at type $j$. Independently for every present cross factor $i<j$, choose a uniform bijection between its two stub sets. Independently for every present internal factor, pair its stubs uniformly. Scales are admissible when every $n_id_{ii}$ is even. Retain loops and parallel edges; each loop contributes two incidences. The type support may be connected, as in the intended application, but the proof does not require connectivity.

Assume symmetric omission weights on the present factors:
\begin{equation}\label{ch:eq:q}
 q_{ij}=q_{ji},\qquad 0\le q_{ij}<\tfrac12,\qquad
 \sum_j d_{ij}q_{ij}=1\quad\text{for every }i.
\end{equation}
Call a factor or stub \emph{active} when its weight is positive. Present means $d_{ij}>0$ and permits a zero weight. Every vertex has at least three active physical stubs, since two weights strictly below $1/2$ cannot sum to one. The active type support need not be connected.

For a fixed prime $p$, choose finite-dimensional $\Fp$ spaces $V_i,U_i$, allowing dimension zero. On each present ordered type pair choose a linear injection
\[
 L_{ij}:V_j\longrightarrow U_i,
\]
required to be an isomorphism when $q_{ij}>0$. Use this same map in every slot with that ordered type pair. Opposite-direction maps need not be inverses. Fix arbitrary linear maps $T_i:V_i\to U_i$. Let $Z_T(G_N)$ count assignments satisfying
\begin{equation}\label{ch:eq:system}
 Y_v\in V_i,\qquad
 \sum_{\text{stubs }vu\text{ at }v}L_{ij}(Y_u)=T_i(Y_v),
 \qquad \operatorname{type}(v)=i,\ \operatorname{type}(u)=j.
\end{equation}

\begin{theorem}[Bounded channel moments]\label{ch:thm:main}
For each fixed set of the preceding data,
\[
 \sup_{N\ \mathrm{admissible}}\E Z_T(G_N)<\infty.
\]
The bound can be maximized over the finitely many admissible channels and on-site maps on any fixed set of finite spaces. It also holds after direct conditioning of every factor on simplicity, at sizes where the required simple factors exist.
\end{theorem}

\begin{corollary}[Raw intertwiner moments]\label{ch:cor:intertwiner}\label{thm:frame-bounded-moment}
Let $A_N$ be the adjacency matrix of the configuration frame, with a loop counted twice on the diagonal. Under \eqref{ch:eq:q}, for every fixed prime $p$ and dimension $r$, uniformly in $R\in\operatorname{Mat}_r(\Fp)$,
\[
 \sup_N\E\#\{Y\in\Fp^{n\times r}:A_NY=YR\}<\infty.
\]
Every fixed power of this count also has bounded expectation, in both configuration and independent uniform-simple models.
\end{corollary}

The corollary is the case $V_i=U_i=\Fp^r$, $L_{ij}=I$, and $T_i(x)=xR$. The count is raw: no planted subspace is divided out and no exceptional graph is deleted.

\begin{corollary}[Regular-lift specialization]\label{ch:cor:lift}\label{thm:lift-bounded-moment}
Let $B$ be a fixed connected simple $d$-regular base with $d\ge3$,
and let $A_N$ be its independent uniform random $N$-lift.
For every fixed prime $p$ and dimension $r$,
\[
 \sup_{N\ge1}\sup_{R\in\Fp^{r\times r}}
 \E|\{Y\in\Fp^{n\times r}:A_NY=YR\}|<\infty.
\]
This includes bipartite bases and every positive lift size.
\end{corollary}
\begin{proof}
Take $d_{ij}$ to be the base adjacency entries, $a_i=1$,
$q_{ij}=1/d$ on present factors, and all channels identity.
Every factor is a perfect matching, so the configuration law is
exactly the simple random-lift law.
Apply Corollary~\ref{ch:cor:intertwiner}.
\end{proof}

For an explicit forced-zero example, take equal type sizes and
\begin{equation}\label{ch:eq:example}
 (d_{ij})=\begin{pmatrix}0&3\\3&1\end{pmatrix}.
\end{equation}
The first row forces $q_{12}=1/3$ and the second forces $q_{22}=0$. Total degrees are three and four. With zero on-site maps in a common ambient space, phases can have $H_2\subsetneq H_1$, with output spaces $K_1=H_2$, $K_2=H_1$. The inactive restricted internal channel is the inclusion $H_2\to H_1$. Thus a common-subspace induction would not be closed, even though the full type support is nonbipartite.

\subsection{Exact ordered-star count and entropy}
A legal star at type $i$ consists of a root $x\in V_i$ and one neighbor $y_s\in V_{j(s)}$ for every labeled physical slot $s$, satisfying its equation in \eqref{ch:eq:system}. Let $t_i(\omega)$ be counts of these legal profiles and let $m_{ij}(x,y)$ pool the root-neighbor table over the $d_{ij}$ slots aimed at $j$. Cross compatibility is $m_{ji}=m_{ij}^{\mathsf T}$; internal compatibility is symmetry with even diagonal. Set $W=0$ otherwise. On compatible arrays put
\begin{align}
 W_{ij}(m)&=\frac{\prod_{x\in V_i,y\in V_j}m_{ij}(x,y)!}{(n_id_{ij})!},\quad i<j,\label{ch:eq:cross}\\
 W_{ii}(m)&=\frac{\prod_x(m_{ii}(x,x)-1)!!\prod_{x<y}m_{ii}(x,y)!}{(n_id_{ii}-1)!!},
 \qquad W(m)=\prod_{i\le j\text{ present}}W_{ij}(m).
 \label{ch:eq:internal}
\end{align}
Absent factors are omitted and $(-1)!!=1$. These are exact labeled-stub matching probabilities. Consequently
\begin{equation}\label{eq:fmp-exact-types}
 \E Z_T=\sum_t\left(\prod_i\frac{n_i!}{\prod_\omega t_i(\omega)!}\right)W(m).
\end{equation}
The roots have already been assigned by these multinomials.

Write $P_i=t_i/n_i$ for star laws, $\mu_i$ for root laws, and $\nu_{ij}=m_{ij}/(n_id_{ij})$ for pooled pair laws. Compatibility implies that $\nu_{ij}$ has marginals $\mu_i,\mu_j$, including two copies of $\mu_i$ internally. Uniform factorial estimates, valid with zero entries, give log contribution $n\Phi+O(\log(n+1))$, where
\begin{equation}\label{ch:eq:entropy}
 \Phi(P)=\sum_i\alpha_iH(P_i)
 -\sum_{i<j}\alpha_id_{ij}H(\nu_{ij})
 -\sum_i\frac{\alpha_id_{ii}}2H(\nu_{ii}).
\end{equation}
There are only polynomially many empirical profiles at fixed data.

Permute slots only within the same target type and symmetrize, obtaining $\bar P_i$. This preserves legality, since those slots have the same channel, and preserves pooled pair laws. Put $\Gamma_i=H(\bar P_i)-H(P_i)\ge0$. Strict concavity gives $\Gamma_i=0$ precisely when $P_i=\bar P_i$. For the symmetric law, every slot to $j$ has law $\nu_{ij}$. Let $I_{ij}$ be its root-neighbor mutual information, and set
\[
 J_{it}=\sum_{s\ne t}H(Y_s\mid X)-H((Y_s)_{s\ne t}\mid X)\ge0.
\]
The root and retained neighbors determine an omitted neighbor on legal stars, by channel injectivity. Averaging the chain rule over omissions with weights $q_{it}$ yields
\begin{equation}\label{ch:eq:localentropy}
 H(\bar P_i)=(2-d_i)H(\mu_i)
 +\sum_jd_{ij}(1-q_{ij})H(\nu_{ij})-\sum_tq_{it}J_{it}.
\end{equation}
Here $\sum_tq_{it}=1$ is a sum over physical slots. Subtract the matching entropies in \eqref{ch:eq:entropy} and expand $H(\nu_{ij})=H(\mu_i)+H(\mu_j)-I_{ij}$. Reversibility and symmetry give coefficient
\[
 \alpha_i\left[(2-d_i)+\sum_jd_{ij}(1-2q_{ij})\right]=0
\]
at every $H(\mu_i)$. Thus, with $w_{ij}=\alpha_id_{ij}$ for $i<j$ and $w_{ii}=\alpha_id_{ii}/2$,
\begin{equation}\label{ch:eq:identity}
 \Phi=-\sum_i\alpha_i\Gamma_i
 -\sum_{i\le j\text{ present}}w_{ij}(1-2q_{ij})I_{ij}
 -\sum_i\alpha_i\sum_tq_{it}J_{it}\le0.
\end{equation}
All mutual-information coefficients are strictly positive, even on inactive factors. Equality requires every star to be symmetric, every present pair to be a product, and every retained-neighbor tuple after an \emph{active} omission to be a product conditional on the root. No condition $J_{it}=0$ is inferred at an inactive omission.

\subsection{The complete finite phase family}
\begin{lemma}[Three-coordinate reflection]\label{ch:lem:triple}
Suppose $A,B,C$ take values in a finite $\Fp$ space, their sum is fixed, and all three pair laws are products of their marginal laws. Then the three laws are uniform cosets of one common subspace.
\end{lemma}
\begin{proof}
Write the fixed sum as $b$. For any support point $a$ of $A$, pair independence and the sum give
\[
 \mu_B(z)=\mu_C(b-a-z)\qquad\text{for every }z.
\]
Conditioning instead on support points of $B,C$ gives reflected translates in every pair, so all support-difference spaces agree, say with $K$. Compare the displayed identity at $a,a'$: $\mu_C$ is invariant under every difference $a-a'$, and hence under their additive span $K$. Its support lies in one coset of $K$; invariance makes it uniform on that whole coset. The reflected translates give the other two laws. Additive subgroups of an $\Fp$ space are subspaces, including the zero subgroup.
\end{proof}

At an entropy zero fix a type $i$ and a root support point $x$. Transform the neighbor variables by their channels into $U_i$. Select three active physical slots $a,b,c$, and condition all other neighbor variables to any tuple in their product support. By the three active retained-product identities, each pair among the selected three retains its original product law after this conditioning. Their sum is fixed by the local equation. Lemma~\ref{ch:lem:triple} forces their image laws to be uniform cosets of a common $K_i\le U_i$. These are their original marginals, independent of the conditioning. Fix two active slots and successively use any other active slot as the third. Every active image law at $i$ therefore has this same support-difference subspace $K_i$.

Every old type $j$ occurs as an active neighbor somewhere: its unit row supplies an active outgoing factor, and reversibility and symmetry supply the opposite active factor. Its channel there is an isomorphism. Consequently every root law is uniform on a coset $c_j+H_j$, with $H_j\le V_j$, and
\begin{equation}\label{ch:eq:activephase}
 L_{ij}(H_j)=K_i\qquad(q_{ij}>0).
\end{equation}
No connectivity of the active support was used. At an active omission the retained variables are independent, include another active uniform-$K_i$ variable, and have image-sum support a coset of
\[
 K_i+\sum_{s\text{ inactive}}L_{i,j(s)}(H_{j(s)}).
\]
By the local equation this is a reflected translate of the omitted active image support, a coset of $K_i$. Therefore
\begin{equation}\label{ch:eq:inactivephase}
 L_{ij}(H_j)\subseteq K_i\quad\text{on every factor}.
\end{equation}
Varying the root over its coset in the local equation now gives
\begin{equation}\label{ch:eq:onsitephase}
 T_i(H_i)\subseteq K_i,\qquad
 T_i(c_i)-\sum_jd_{ij}L_{ij}(c_j)\in K_i.
\end{equation}
Multiplicity scalars are interpreted in the field; no degree is inverted.

Conversely suppose the subspaces and root cosets satisfy \eqref{ch:eq:activephase}--\eqref{ch:eq:onsitephase}. Choose the root uniformly on its coset and all neighbors except one active slot independently uniformly on their cosets. Solve for the omitted neighbor by its isomorphism. The result lies in the required coset. This samples uniformly the entire legal affine star support, of size
\[
 |H_i|\frac{\prod_s|H_{j(s)}|}{|K_i|}.
\]
The same parametrization works at every active omission, proving all active retained-product identities. A different active omission proves every root-neighbor pair is a product. Uniform legal supports are invariant under repeated-target slot permutations. Thus this is an entropy zero, uniquely determined by the cosets and subspaces: the root and any active retained tuple already have their specified product law and determine the last label.

There are finitely many phases at fixed finite spaces. On restriction put
\begin{equation}\label{ch:eq:restriction}
 V_i'=H_i,\quad U_i'=K_i,\quad L_{ij}'=L_{ij}|_{H_j},\quad T_i'=T_i|_{H_i}.
\end{equation}
Every channel remains injective and every active channel remains an isomorphism. Hence this is a closed induction class. If every $H_i=V_i$, an active isomorphism in row $i$ forces $K_i=U_i$. This is the unique full-uniform phase. Every proper phase strictly reduces $\sum_i\dim V_i$. At total root dimension zero, active isomorphisms force all output spaces to be zero and there is exactly one assignment.

\begin{corollary}[Positive-certificate phase specialization]\label{ch:cor:positive-phases}
Suppose the present type support is connected, every present
$q_{ij}$ is positive, all root and output spaces are the same
$V=\Fp^r$, and every channel is identity. Then the entropy zeros
have the following more specific forms.
On a nonbipartite support there is one common subspace $H$ with
$H_i=K_i=H$. On a bipartite support there are two subspaces
$H_{\mathcal L},H_{\mathcal R}$, with $H_i$ equal to its side's
subspace and $K_i$ equal to the opposite side's subspace.
In either case the root laws are uniform on cosets $c_i+H_i$,
\[
 T_i(H_i)\subseteq K_i,\qquad
 T_i(c_i)-\sum_jd_{ij}c_j\in K_i,
\]
and these data give precisely the unique uniform legal-star phases.
\end{corollary}
\begin{proof}
All present factors are active, so \eqref{ch:eq:activephase} says
$H_j=K_i$ for every neighboring type. The subspaces $H_i$ are
therefore constant along even walks. Connected nonbipartite support
has one even-walk component, counting an internal factor as a
self-incidence; connected bipartite support has its two parts.
The displayed conditions and converse are exactly
\eqref{ch:eq:onsitephase} and the preceding phase construction.
\end{proof}

This common-subspace or two-subspace conclusion requires positivity
on every present factor. With zero weights, the full present support
can be nonbipartite while some $H_i$ differ, as in
\eqref{ch:eq:example}. In particular the scalar sharp-moment
refinement uses this positive-certificate specialization; no exact
scalar coefficient is inferred here for general zero-weight frames.

\subsection{The heterogeneous full-uniform central estimate}
Retain a separate table $M_{is}(x,y)$ for each physical slot, even when slots have the same old target. If that target is $j$, the table has width $|V_j|$. Write $A_i=|V_i|$, $B_{is}=|V_{j(s)}|$, and $C_i=|U_i|$. Let
\[
 \mu_{is}(x,y)=\frac1{A_iB_{is}},\qquad
 \Delta=M-(n_i\mu_{is})_{i,s}.
\]
Use a fixed Euclidean norm. We prove that a fixed neighborhood $\|\Delta\|\le\epsilon n$ has expected contribution with limsup at most
\begin{equation}\label{ch:eq:index}
 I_*:=\prod_i C_i^{A_i}.
\end{equation}

\paragraph{Exact reference normalization}
Let $P_i$ be uniform on legal stars and $Q_i$ independently uniform on the root and all neighbor slots. One active isomorphism can absorb any other labels, so the legal-star count is $k_i=A_i\prod_sB_{is}/C_i$. Put $K_N=\prod_i k_i^{n_i}$, and let $P_N,Q_N$ be the independent-vertex full-table laws. Exact counting gives
\begin{equation}\label{ch:eq:normalizers}
 \E Z_T=K_N\E_{P_N}W,\qquad
 \E_{Q_N}W=\prod_{i,s}B_{is}^{-n_i}.
\end{equation}
The second identity follows because each configuration and root labeling specify exactly one consistent neighbor labeling. Active isomorphisms and \eqref{ch:eq:q} give
\begin{align*}
 \sum_i n_i\log C_i
 &=\sum_{i,j}n_id_{ij}q_{ij}\log|V_j|\\
 &=\sum_jn_j\log|V_j|\sum_i d_{ji}q_{ji}
 =\sum_jn_j\log A_j.
\end{align*}
Consequently
\begin{equation}\label{ch:eq:exactone}
 K_N\E_{Q_N}W=\prod_i(A_i/C_i)^{n_i}=1.
\end{equation}
The normalized weighted law
\[
 \nu_N(m)=K_NQ_N(M=m)W(m)
\]
is exactly the full-table law of a random configuration followed by iid uniform roots in their respective $V_i$. Indeed, for each full ordered-star array, $K_NQ_N$ is $\prod_i C_i^{-n_i}=\prod_i A_i^{-n_i}$, its uniform-root probability, and $W$ is its configuration probability. This proves the assertion before passing to tables.

\paragraph{Covariance and maximal lattices}
At old type $i$ let $\mathcal H_i$ be the real space of slot-table deviations having zero total in every table and a common root-row margin vector. Its dimension is
\[
 h_i=(A_i-1)+A_i\sum_s(B_{is}-1).
\]
The means and covariance of the table-indicator vector agree under $P_i,Q_i$. Given the root and any two neighbor labels, an active slot distinct from those two can absorb the constraint. Counting free remaining labels proves the joint law of the root and the two specified neighbors is product uniform. This gives every first and second moment used by the tables.

Their common covariance is positive definite on $\mathcal H_i$. Under $Q_i$ a test array gives the functional $F=\sum_s f_s(X,Y_s)$. If its variance is zero, varying just $Y_s$ shows $f_s(x,y)$ is constant in each row. Varying $x$ shows that the sum of those row constants is constant in $x$. These are exactly the functionals vanishing on $\mathcal H_i$. Thus no nonzero direction in that space has zero variance. Equality of covariance gives the same fact for $P_i$; its support differences consequently span the entire space as well.

The unrestricted support-difference lattice is $\mathcal L_{0i}=\mathcal H_i\cap\mathbb Z^{\text{table entries}}$. Within-row transfers and root changes generate it. Write $\mathcal L_i$ for the lattice generated by legal-star differences. If $r(x)$ is the common root-row sum of a difference array, it satisfies the congruences
\begin{equation}\label{ch:eq:congruence}
 C_x(m)=\sum_s\sum_yL_{i,j(s)}(y)m_s(x,y)-T_i(x)r(x)=0
 \quad\text{in }U_i,\qquad x\in V_i.
\end{equation}
These congruences are complete. To prove this, represent any circle-valued linear functional constant on legal stars by functions $g_s(x,y)\in\mathbb R/\mathbb Z$. Fix $x$ and three active slots, fixing all other labels. In image coordinates their functions satisfy
\[
 f_a(u)+f_b(v)+f_c(t-u-v)=\text{constant}\quad(u,v\in U_i).
\]
Comparison at $(u+h,v)$ and $(u,v+h)$ proves
\[
 f_a(u+h)-f_a(u)=f_b(v+h)-f_b(v)=\chi(h).
\]
The increment is independent of $u,v$ and additive in $h$. All three functions are therefore constants plus the same additive character $\chi$ of $U_i$. Vary any other neighbor label and compensate in an active slot: its function must be a constant plus $\chi\circ L_{ij}$, even if the channel is inactive and not onto. The original active-slot function fixes $\chi$ independently of the initially fixed other labels. Thus for each root $x$ all functions have form
\[
 g_s(x,y)=\beta_s(x)+\chi_x(L_{i,j(s)}y).
\]
Their sum on a legal star is $\sum_s\beta_s(x)+\chi_x(T_i(x))$, which must be constant across roots. Modulo functionals constant on unrestricted stars, these are precisely the characters of \eqref{ch:eq:congruence}.

A within-row transfer in one active slot changes $C_x$ arbitrarily in $U_i$, independently for each $x$. Hence the congruence map $\mathcal L_{0i}\to U_i^{V_i}$ is surjective. Since $\mathcal L_i$ has full real rank and the character calculation leaves no additional character, it is exactly the kernel. A strict sublattice of that kernel would give a separating character of a nontrivial finite quotient, a contradiction. Therefore
\begin{equation}\label{ch:eq:localindex}
 [\mathcal L_{0i}:\mathcal L_i]=C_i^{A_i}.
\end{equation}
The legal affine lattice at size $n_i$ is contained in the unrestricted affine lattice, since a legal star is also unrestricted. No support-saturation assumption at small sizes is made. If $h_i=0$, determinants and covolumes at that type are interpreted as one; then $C_i=1$ as well.

\paragraph{Relative tilted local limit estimate}
We record the analytic argument needed on a fixed neighborhood. For a finite-support lattice law with positive support probabilities, positive covariance on its real difference space, and maximal support-difference lattice $\mathcal L$, the log moment generating function $\Lambda$ is analytic with positive Hessian at zero. The inverse function theorem supplies an analytic mean-to-tilt map on a fixed mean neighborhood. On a smaller compact neighborhood, tilted probabilities stay bounded below, covariance stays positive, and derivatives are uniformly bounded. Exponential tilting and Fourier inversion at the tilted mean give, uniformly for $m\in nz+\mathcal L$ with $m/n$ in that neighborhood,
\begin{equation}\label{ch:eq:llt}
 \Prob(S_n=m)=\operatorname{covol}(\mathcal L)(2\pi n)^{-h/2}
 (\det\Sigma_{m/n})^{-1/2}e^{-n\mathcal I(m/n)}
 [1+O(n^{-1/2})].
\end{equation}
Here all determinants and covolumes use the same Euclidean coordinates. For a direct verification, on a fundamental dual torus zero is the sole modulus-one point, by maximality of the support-difference lattice. Compactness of the tilted laws makes the characteristic modulus at most a fixed $\rho<1$ away from zero. Near zero the centered logarithm is
\[
 -\tfrac12\langle\xi,\Sigma_{m/n}\xi\rangle+O(\|\xi\|^3)
\]
uniformly, with a uniform Gaussian modulus bound. Rescale by $\sqrt n$ and integrate the cubic error against that Gaussian. This gives relative error $O(n^{-1/2})$; the outer torus is exponentially small. The leading Gaussian integral stays bounded above and below after its scaling. This also proves positivity at every indicated affine lattice point for sufficiently large $n$.

Apply \eqref{ch:eq:llt} at each type to $P_i,Q_i$. Their means and covariance agree, so their analytic rate functions agree through quadratic order. Their covolume ratio is \eqref{ch:eq:localindex}, and the determinant ratio stays bounded and tends to one at the mean. Multiplying over $n_i=\alpha_i n$ gives, for some fixed neighborhood,
\begin{align}
 \frac{P_N(M=m)}{Q_N(M=m)}&\le C I_*e^{C\|\Delta\|^3/n^2}
 &&(\|\Delta\|\le\epsilon_0n),\label{ch:eq:cubic}\\
 \frac{P_N(M=m)}{Q_N(M=m)}&=I_*(1+o(1))
 &&(\|\Delta\|\le\sqrt n\log n)\label{ch:eq:ratio}
\end{align}
on the legal affine lattice. Away from it the numerator vanishes. Affine containment and \eqref{ch:eq:llt} ensure a positive denominator wherever the numerator is nonzero in these neighborhoods, for large sizes. All admissible size residues are included.

\paragraph{Weighted concentration and integration}
Under $\nu_N$, generate each internal pairing by a uniform stub permutation paired in consecutive positions, and each cross bijection by a uniform permutation. Every internal pairing has $2^{L/2}(L/2)!$ permutation preimages, so this is exactly uniform. There are $O(n)$ permutation entries and independent root labels in total. A transposition changes at most two paired edges, hence at most four physical incidences. Changing one root label affects at most $2D$ incidences. These bounds also hold for loops and repeated edges.

Expose permutation images and then root labels. Uniform completions after two choices for the next image are coupled by one transposition preserving earlier images; thus the conditional expectations of a scalar table entry have bounded range. Root-label exposures have the same property. The elementary bounded-range martingale exponential estimate, with total squared range $O(n)$, and a fixed-dimensional union bound give
\[
 \nu_N(\|M-\E_{\nu_N}M\|\ge t)\le Ce^{-ct^2/n}.
\]
Cross table means are their independent means. In an internal slot the chance that the partner stub has the same vertex is $(d_{ii}-1)/(n_id_{ii}-1)$. Its expected table entry is consequently
\[
 n_i\left[\frac{1-\rho_i}{|V_i|^2}
       +\frac{\rho_i\one_{x=y}}{|V_i|}\right],
 \qquad\rho_i=\frac{d_{ii}-1}{n_id_{ii}-1}.
\]
This differs from the independent mean by $O(1)$ in each entry. Therefore
\begin{equation}\label{ch:eq:tail}
 \nu_N(\|\Delta\|\ge t)\le Ce^{-ct^2/n}\qquad(t\ge C_0\sqrt n).
\end{equation}
The estimate holds for all larger deviations.

By \eqref{ch:eq:normalizers}--\eqref{ch:eq:exactone}, the central contribution is the $\nu_N$ expectation of the likelihood ratio $P_N/Q_N$. Equation~\eqref{ch:eq:ratio} bounds the logarithmic window by $I_*+o(1)$. Choose fixed $\epsilon>0$ with $C\epsilon<c/2$. On $\|\Delta\|\le\epsilon n$, the cubic factor in \eqref{ch:eq:cubic} is at most $\exp(\eta\|\Delta\|^2/n)$ with $\eta=C\epsilon<c/2$. For $L=\sqrt n\log n$, integration of \eqref{ch:eq:tail} yields
\begin{align*}
 \E_{\nu_N}\left[e^{\eta\|\Delta\|^2/n}\one_{\|\Delta\|\ge L}\right]
 &\le Ce^{-(c-\eta)L^2/n}
 +C\int_L^\infty\frac{2\eta t}{n}e^{-(c-\eta)t^2/n}\,dt\\
 &=O(e^{-c'\log^2n})=o(1).
\end{align*}
Thus the surrounding fixed neighborhood contributes $o(1)$ outside the logarithmic window. We have proved the central bound \eqref{ch:eq:index}, including its finite lattice factor and all prefactors.

\subsection{Proper phases: the selected sparse certificate}
Fix a proper phase satisfying \eqref{ch:eq:activephase}--\eqref{ch:eq:onsitephase}. Quotient roots by $H_i$, outputs by $K_i$, and subtract the type-constant quotient point $c_i+H_i$. The offset equation makes that point exact for every configuration. A solution has homogeneous quotient defect
\begin{equation}\label{ch:eq:quotient}
 z_v=Y_v+H_i-(c_i+H_i),\qquad
 \sum_{vu}\overline L_{ij}(z_u)=\overline T_i(z_v).
\end{equation}
Every active quotient channel is an isomorphism, since both the ambient and restricted channels are isomorphisms. An inactive quotient channel may have a kernel.

For any fixed quotient assignment, all its full lifts satisfying \eqref{ch:eq:system} form an empty affine fiber or a translate of the homogeneous restricted kernel from \eqref{ch:eq:restriction}. Write $Z_H(G)$ for that kernel count and $M_H(n)=\E Z_H(G_n)$. Hence every nonempty fiber has size $Z_H(G)$, and the zero-defect contribution is at most $M_H(n)$. Representatives need not solve the ambient unquotiented equation.

For nonzero support $S=\{v:z_v\ne0\}$ of size $s$, select all active edges incident to $S$, retaining their actual labeled stubs. At an outside zero vertex with exactly one active incidence from $S$, that incidence has nonzero contribution through its active quotient isomorphism. Its equation forces at least one further incidence from $S$ with nonzero contribution, necessarily inactive. Select one such cancellation edge. A fixed stub order can make this choice deterministic. At a vertex with two or more active incidences no extra edge is needed; vertices with no active incidence are not added.

Every selected edge touches $S$. Every selected outside vertex now has at least two selected boundary incidences, while every vertex of $S$ has at least three selected incidences. If $a$ selected edges are internal to $S$, $b$ are boundary edges, $t$ is the outside count, and $E=a+b$, then
\begin{equation}\label{ch:eq:excess}
 2a+b\ge3s,\quad b\ge2t,\quad E\le Ds,\quad
 E-(s+t)\ge\frac{s}{2},\quad t\le\frac{Ds}{2}.
\end{equation}
A loop counts twice toward support degree and once in $a$. Multiplicities and $s=1$ are retained. The inactive support stubs need not be saturated, and there is no absent-edge condition on unselected stubs.

\paragraph{Exact conditional forcing and rank control}
A certificate $C$ prescribes distinct stub pairs consistent with their factors. To force a new cross pair $u\mapsto v$, swap the image of $u$ with the current preimage of $v$. With $m$ unused stubs on each side, each permitted output has exactly $m$ preimages. Earlier prescriptions remain fixed, so uniform conditional input gives uniform conditional output. For an internal pair $\{u,v\}$, if not already paired, replace $\{u,x\},\{v,y\}$ by $\{u,v\},\{x,y\}$. With $m$ free stubs each output has exactly
\[
 1+2(m-2)/2=m-1
\]
preimages: itself and the two inverse orientations for every other output pair. This is also exact uniform conditional forcing, unaffected by coincident vertex owners. Sequential forcing over the independent factors gives a coupled $G^C$ with exactly the configuration law conditional on $C$, using at most $E$ switches.

Only the owners of at most four stubs have changed equation rows at each switch. With $r_{\rm out}=\max_i\dim U_i$, the restricted rectangular operators therefore satisfy
\[
 \rank(\mathcal A_H(G^C)-\mathcal A_H(G))\le4r_{\rm out}E,
 \qquad Z_H(G^C)\le p^{4r_{\rm out}E}Z_H(G).
\]
Taking expectations under the coupling proves
\begin{equation}\label{ch:eq:conditional}
 \E[Z_H\mid C]\le p^{4r_{\rm out}E}M_H(n).
\end{equation}
No independence of the certificate and the restricted kernel is assumed.

Let $\lambda=\min_{d_{ij}>0}\alpha_id_{ij}>0$. If $s\le\epsilon n$ and $\epsilon\le\lambda/(4D)$, then every prescribed-pair probability denominator is at least $\lambda n/2$: cross probabilities are reciprocal falling factorials, and internal probabilities use $m-1,m-3,\ldots$. Since $E\le Ds$,
\begin{equation}\label{ch:eq:weighted}
 \Prob(C)\le(C_0/n)^E,\qquad
 \E[Z_H\one_C]\le M_H(n)(C_0p^{4r_{\rm out}}/n)^E.
\end{equation}

\paragraph{Enumeration and the boundary neighborhood}
For fixed $s,t,E$, choose support and outside vertices in at most $n^{s+t}/(s!t!)$ ways. Quotient labels on the support cost at most $p^{r_{\rm root}s}$, where $r_{\rm root}=\max_i\dim V_i$. Selected stubs at the support and outside vertices cost at most $2^{D(s+t)}$. This includes the necessary choice of unsaturated inactive stubs at $S$; allowing arbitrary subsets only overcounts. The selected $2E$ stubs have at most $(2E)^E\le(2Ds)^E$ pairings. Discard invalid colors, repeated stubs, edges not touching $S$, and degree patterns failing \eqref{ch:eq:excess}.

Associate a quotient solution to its selected certificate. Drop the further quotient equations, the tests choosing a particular cancellation edge, and all restrictions on unselected stubs. By the affine-fiber fact, the sum of $Z_H\one_C$ over the resulting assignments and certificates dominates the original solution count. Applying \eqref{ch:eq:weighted}, every fixed factor is bounded by a constant to the power $s$, since $E\le Ds$ and $t\le Ds/2$. The remaining powers obey
\begin{equation}\label{ch:eq:powers}
 \frac{s^E n^{s+t-E}}{s!t!}
 =\left(\frac{s}{n}\right)^{E-s-t}\frac{s^s}{s!}\frac{s^t}{t!}
 \le\left(\frac{s}{n}\right)^{s/2}\frac{s^s}{s!}\frac{s^t}{t!}.
\end{equation}
Use $s^s/s!\le e^s$, $\sum_{t\ge0}s^t/t!=e^s$, and at most $Ds+1\le(D+1)^s$ choices of $E$. For a fixed $B<\infty$,
\begin{equation}\label{ch:eq:boundary}
 \E\#\{Y\text{ solves \eqref{ch:eq:system}}:\#\{v:z_v\ne0\}=s\}
 \le M_H(n)B^s(s/n)^{s/2},\qquad1\le s\le\epsilon n.
\end{equation}
For $f_s=B^s(s/n)^{s/2}$,
\[
 \frac{f_{s+1}}{f_s}
 =B\sqrt{\frac{s+1}{n}}(1+1/s)^{s/2}
 \le B\sqrt{\frac{e(s+1)}n}.
\]
Shrink $\epsilon$ so this is at most $1/2$ for consecutive terms in the indicated range. The sum starts at $f_1=B n^{-1/2}$, so
\begin{equation}\label{ch:eq:boundarysum}
 \E\#\{Y:1\le\#\{v:z_v\ne0\}\le\epsilon n\}
 \le O(n^{-1/2})M_H(n).
\end{equation}
Together with the zero-defect bound, this controls the entire root-coset neighborhood of the proper phase.

In particular, when the restricted root dimension is zero and the
phase point is an exact old assignment, its neighborhood contributes
$1+O(n^{-1/2})$. More generally, whenever the restricted moment is
bounded, positive-defect leakage is $o(1)$, also inside a narrower
within-coset edge-table neighborhood. These statements supply the
vanishing boundary errors used by the later scalar and spectral
refinements; no stronger generic support exponent is required.


\subsection{Induction, simplicity, and fixed powers}
Induct on $\sum_i\dim V_i$ in the stated channel class. Total dimension zero has exactly one assignment. Every proper phase has a strictly smaller restricted root dimension and remains in the class by \eqref{ch:eq:restriction}. Its expected restricted kernel count is bounded by induction, so \eqref{ch:eq:boundarysum} and the affine-fiber bound control its full root-coset neighborhood. There are finitely many proper phases. The full-uniform phase is controlled by the central table estimate.

These are open neighborhoods in the compact feasible-profile space after using strict inequalities in their definitions. Near a proper phase the total root mass outside its cosets is arbitrarily small; near the full phase its physical pair tables are close to their uniform means. They cover every entropy zero. The complementary closed feasible set has strictly negative maximum of the continuous exponent \eqref{ch:eq:entropy}. Its exact type sum in \eqref{eq:fmp-exact-types} is at most $\operatorname{poly}(n)e^{-cn}$. Therefore the total expectation has finite limsup. Each of the finitely many smaller admissible sizes has finite count. This proves the configuration assertion of Theorem~\ref{ch:thm:main}. The finite map sets permit a uniform maximum on fixed finite spaces.

For completeness, the joint simplicity probability is bounded below by a positive constant along admissible large sizes. In an internal degree-$b$ factor, the loop and unordered double-edge pattern counts have limiting factorial moments of independent Poisson variables with means $(b-1)/2$ and $(b-1)^2/4$. In a cross factor with endpoint degrees $b,c$, the double-edge mean is $(b-1)(c-1)/2$. To verify these assertions directly, count ordered lists of fixed stub patterns: vertex-disjoint patterns give the displayed products using reciprocal falling factorial or pairing denominators; every consistent overlap uses fewer vertex choices relative to prescribed edges and is $O(1/n)$. For each fixed factorial moment the number of overlap types is finite. Bonferroni inequalities for the zero-pattern event, followed by the limit in size and then in truncation order, give the positive zero-count limits $e^{-\lambda}$. No loops and no double-edge patterns is exactly simplicity. Independence of the finitely many factors gives a positive joint limit.

Thus $\E[Z_T\mid\mathrm{simple}]\le\E Z_T/\Prob(\mathrm{simple})$ is bounded. Each simple graph in a factor has the same number of labeled-stub representations, so its conditional projection is exactly uniform; factor independence persists. Finitely many smaller feasible simple sizes again cause no issue. This proves the simple-model assertion directly, without a qualitative contiguity claim about an unbounded observable.

Finally $Z_T(G)^s$ counts $s$ independent solution coordinates on the same graph. This is the channel system on the direct sums of $s$ copies of all spaces and maps, which still satisfies every hypothesis. For intertwiners it is the target $\operatorname{diag}(R,\ldots,R)$. Theorem~\ref{ch:thm:main} therefore proves every fixed-power bound and Corollary~\ref{ch:cor:intertwiner}.

The resulting exponential-nullity bounds can be inserted into a separately certified kernel and enumeration routine whose conditional charged cost is $\widetilde O(n^2p^{s h})$. That composition also requires structural identifiability, the appropriate right kernel, and certified bit-work accounting. This note proves the moment input; it does not construct or audit those separate primitives, assert a growing-field estimate, or establish historical priority.

\section{Sharp scalar moments with positive omission weights}
\label{sec:frame-scalar}
Use the configuration law of Theorem~\ref{thm:frame-bounded-moment}.
In this section require a certificate with \emph{positive} weights on
\emph{every} present factor, $0<q_{ij}<1/2$, and assume connected
nonbipartite old support, counting an internal factor as a self-incidence.
Common total degree $d\ge3$ is a special case. The more general
zero-weight theorem does not imply the scalar equality below.
Fix an odd prime $p$ and
$\lambda\in\Fp$, and put $\ell=\dim_{\Fp}\ker(Q-\lambda I)$.

\begin{theorem}[Exact configuration scalar moment]\label{thm:frame-scalar}
Along every configuration-admissible size sequence,
\[
 \E|\ker_{\Fp}(A-\lambda I)|=2p^\ell+o_{Q,p}(1).
\]
The fixed size-ratio vector is implicit in the error constant.
No degree, size, or pool multiplier is required to be invertible modulo
$p$. The exact equality is in the configuration law; after conditioning
simple we use its upper bound.
\end{theorem}
\begin{proof}
We sharpen the full-uniform sector in the preceding moment proof.
Use its full physical-slot tables, legal law $P$, unrestricted law $P_0$,
and exact normalized weighted reference law $\nu_N=K_NP_0W$.
Let $L_0$ be the unrestricted integer table-difference lattice and let
$L$ be its legal sublattice. For each old type and root label the legal
congruence is
\[
 C_{i,x}(m)=\sum_s\sum_y y\,m_{is}(x,y)-\lambda x\,r_i(x)=0
 \quad\text{in }\Fp,
\]
where $r_i$ is the common row margin. The preceding maximal-lattice
proof gives $[L_0:L]=I=p^{kp}$. Let $K_0\le L_0$ impose pooled
cross reversal and pooled internal symmetry with even diagonal.
The all-zero star array $m_0(N)$ is compatible and legal at every
admissible size, so the compatible affine lattices are
$m_0(N)+K_0$ and $m_0(N)+(K_0\cap L)$. Write
$J_*=[K_0:K_0\cap L]$.

For each fixed $z\in K_0$, uniformly in the logarithmic central window,
\begin{equation}\label{eq:frame-scalar-shift}
 \nu_N(m+z)/\nu_N(m)=1+o(1).
\end{equation}
For $P_0$ this follows from its relative full-table local limit
estimate, since a bounded shift changes its Gaussian exponent by
$O(\log N/\sqrt N)$. For a cross pool the exact factorial ratio has
leading factor $(N_id_{ij}/p^2)^{\sum_{x,y}z_{ij}(x,y)}=1$.
For an internal pool, even diagonal shifts contribute half their
diagonal sum and off-diagonal shifts their full sum. The exponent is
$\frac12\sum_xz_{ii}(x,x)+\sum_{x<y}z_{ii}(x,y)=0$.
All entries are of order $N$ in this window; parity and compatibility
are preserved by $z$. Ratios are therefore uniformly $1+o(1)$,
also for negative shifts.

The actual-reference concentration proved above puts mass $1-o(1)$
in this window and any fixed-width inward shrinkage. Translation by
fixed representatives of $K_0/(K_0\cap L)$ consequently makes its
coset masses equal up to $o(1)$: sum
\eqref{eq:frame-scalar-shift} on an inner window, and use the
subgaussian tails for both complements. There are finitely many
cosets, so each has mass $1/J_*+o(1)$. The legal/unrestricted
likelihood ratio is $I(1+o(1))$ on the legal central lattice.
The fixed-neighborhood annulus has vanishing weighted contribution
by the preceding cubic-ratio and reference-tail integration.
Hence the full-uniform expected contribution is $I/J_*+o(1)$.

We compute this index directly on the integer lattice. An annihilator
of the congruence image $C(K_0)$ has functions $h_i:\Fp\to\Fp$.
Insert a pooled zero-margin two-by-two interaction in one physical
slot on each orientation of a cross factor. Its character equation is
\[
 [h_i(x)-h_i(x')](y-y')+
 [h_j(y)-h_j(y')](x-x')=0.
\]
Anchor labels zero and one give $h_i(x)=s_ix+c_i$, with
$s_i+s_j=0$. For an internal factor, insert twice the symmetric
interaction with diagonal $+1$ at $(x,x),(y,y)$ and off-diagonal
$-1$ at $(x,y),(y,x)$ in one physical slot. It is an actual
integer compatible direction, including even diagonal. It forces
$2[h_i(x)-h_i(y)](x-y)=0$, so $h_i$ is constant.
An odd cross cycle likewise kills every alternating slope.
Connectivity and nonbipartiteness thus give $h_i(x)=c_i$ for all $i$.

For these constants the functional is exactly
\[
 \sum_i[(Q-\lambda I)c]_i\sum_xx\,r_i(x).
\]
The root-margin moments here really vary independently on $K_0$.
Given arbitrary zero-total integer vectors $r_i$, use twice them.
For a cross pool choose an integer table with row $2d_{ij}r_i$
and column $2d_{ji}r_j$, using label zero as an anchor.
For an internal pool use the diagonal table $2d_{ii}r_i(x)$,
which has even diagonal. Split a pooled table into physical slots
by placing the anchor-column table $2r_i(x)$ in the first
$d_{ij}-1$ slots and the residual in the last. Every slot then has
row $2r_i$, total zero, and the pooled constraints hold.
Negative entries are allowed in a direction. Multiplication by two
is invertible modulo $p$, so each $\sum_xx r_i(x)$ is free modulo
$p$. Thus the annihilator is exactly $\ker(Q-\lambda I)$ and has
dimension $\ell$. Conversely every such constant vector annihilates
all compatible directions by the same pooled-margin identity.
Therefore
\[
 J_*=|C(K_0)|=p^{kp-\ell},\qquad I/J_*=p^\ell.
\]
No mod-$p$ saturation of a real compatibility space has been assumed.

For a scalar alphabet the only subspaces are zero and the whole
field. Corollary~\ref{ch:cor:positive-phases} therefore gives precisely
the full-uniform phase and $p^\ell$ point phases $Qc=\lambda c$.
Each point phase is an exact old labeling on every configuration.
Its sparse leakage contribution is $o(1)$ by the preceding boundary
bound with restricted dimension zero; configurations may have
singleton defects, whose contribution is still
$O(N^{-1/2})$. Choose point neighborhoods disjoint, using a radius
smaller than one quarter of the smallest class-size coefficient,
and disjoint from a small uniform-table neighborhood. The compact
entropy complement contributes a polynomial times $e^{-cN}$.
The point and uniform sectors each contribute $p^\ell+o(1)$,
proving the theorem.
\end{proof}

\begin{corollary}[Cubic moment, uniformly prescribed matching]
\label{thm:cubic-raw}
For the two-class $(a,b)=(2,1)$ model with any prescribed cross
matching independent of its internal configurations, every fixed odd
prime satisfies
\[
 \E_{\rm config}|\ker_{\Fp}(A-I)|=2p+o_p(1)
\]
at every $N$, uniformly over that matching. In the simple internal law,
\[
 \limsup_N\sup_B\E|\ker_{\Fp}(A-I)|\le(1+e^{3/2})p.
\]
\end{corollary}
\begin{proof}
Apply Theorem~\ref{thm:frame-scalar} to
$Q=\left(\begin{smallmatrix}2&1\\1&2\end{smallmatrix}\right)$ and
$\lambda=1$. Its old kernel has dimension one for every odd prime.
A relabeling of the minus block gauges any prescribed cross matching
to identity. The independent internal configuration laws are invariant
under that relabeling, so the distribution of the kernel size is
identical for every matching and for an independent uniform matching.
Simplicity of the two internal configurations has probability tending
to $e^{-3/2}$. Subtracting the deterministic planted $p$ vectors before
conditioning gives the displayed simple upper bound.
\end{proof}

The same configuration scalar theorem gives $2p+o_p(1)$ at the planted
target for every fixed two-class $a>b>0$ when $p\nmid b$.
Subtracting the planted line bounds the probability of real nullity
at least two by $1/(p-1)+o(1)$ in the configuration law.
Positive-probability simplicity conditioning and then arbitrarily
large fixed primes give real planted-line simplicity.
The size limit always precedes the prime limit. No polynomial
real-collision rate or exact simple-model scalar limit is inferred.


\section{Certified full kernels and nearly quadratic partition listing}\label{sec:full-kernel}
The next algorithm computes the entire kernel on every fixed symmetric
input, with no corank-one promise. Its classical preconditioning,
recurrence, extension construction, deterministic certificate, retry
cost, and base-field descent are all charged. This stronger algorithm
will be averaged using a bounded raw moment for the actual graph law,
rather than a qualitative high-probability rank assertion.

\subsection{Unconditional per-matrix algorithm}
Fix an odd prime $p$. Let $M=M^T\in\F_p^{n\times n}$, $n\ge1$,
be supplied as a sparse matrix, with $w=\nnz(M)$,
$k=\rank_{\F_p}M$, and $r=n-k$.
All logarithmic constants below may depend on the fixed $p$.

\begin{theorem}\label{fk:thm:kernel}
There is a Las Vegas algorithm, valid for every such fixed $M$,
which returns a basis $C\in\F_p^{n\times r}$ of its entire kernel,
pivot rows $J$ with $C_{J,:}=I_r$, and a deterministically checkable
rank certificate. Its expected work is
\[
 \widetilde O_p\bigl(n(w+n)(r+1)\bigr)
 \quad\text{bit operations}.
\]
Its workspace is $\widetilde O_p(n(r+1)+w)$ bits, apart from a retry
counter. The certificate has $O_p(nr+n\log n)$ bits and can be
verified in $\widetilde O_p(n(w+n))$ bit work.
For bounded-degree graph matrices, $w=O(n)$, giving expected
$\widetilde O_p(n^2(r+1))$ bit work. These assertions include the
zero matrix and nonsingular matrices. No numerical gap or generic
rank profile is assumed.
\end{theorem}

\subsection{Extension field and primary preconditioning inputs}
Choose the least positive $s$ for which
\[
 q=p^s\ge64(n+1)^6,
\]
and construct $K=\F_q=\F_p[z]/(h(z))$ once, reusing it across trials,
using a verified monic irreducible polynomial of degree $s$.
Rank and corank are unchanged
by this scalar extension.
In a fresh trial draw independent
\[
 D=\operatorname{diag}(d_1,\ldots,d_n),\quad d_i\text{ uniform in }K^*,
 \quad u,v\text{ uniform in }K^n,\qquad T=DM.
\]
The following are distinct primary results, both valid in the fixed
characteristic $p$, however small relative to $n$.
Chen et al.\ \cite[Theorem 4.2]{Chen} give
\begin{equation}\label{fk:eq:cyc}
 \Prob(T\text{ not cyclic on its nonzero primary part})
 \le\frac{n(n-1)}{2(q-1)}.
\end{equation}
Their Theorem 4.5, using symmetry of $M$, gives
\begin{equation}\label{fk:eq:nil}
 \Prob(T\text{ has a nilpotent zero block of size }>1)
 \le\frac r{q-1}.
\end{equation}
No independence between these two events is required for their union
bound. The squarefree conclusion of their Theorem 4.3 has a
characteristic restriction and is \emph{not} invoked.

On their common good event,
\[
 K^n=\ker T\oplus\operatorname{im}T,
 \qquad T|_{\operatorname{im}T}\text{ invertible and cyclic of dimension }k.
\]
Repeated nonzero roots of its minimal polynomial cause no problem.

\subsection{A certifying trial, including unlucky choices}
Compute exactly the $2n$ scalars
\begin{equation}\label{fk:eq:seq}
 a_j=u^TT^{j+1}v,\qquad 0\le j<2n.
\end{equation}
Use Berlekamp--Massey to obtain their true shortest recurrence length
$\ell$ and monic polynomial $g$, of degree $\ell$.
If the implementation returns connection coefficients in backward
recurrence order, reverse them before the Horner polynomial evaluation.
The zero sequence has length $0$ and $g=1$.
Reject the trial immediately if $g(0)=0$.

\begin{lemma}[Unconditional lower-rank certificate]\label{fk:lem:lower}
For every $T,u,v$, the finite-sequence complexity in~\eqref{fk:eq:seq}
satisfies $\ell\le\rank T=k$.
\end{lemma}
\begin{proof}
$Tv$ lies in the invariant subspace $\operatorname{im}T$, of
dimension $k$. Cayley--Hamilton on that subspace yields an infinite
recurrence of order at most $k$ and therefore one of that order on
the finite prefix. The minimal prefix length cannot exceed $k$.
This includes singular matrices with nilpotent blocks.
\end{proof}

Set $\widehat r=n-\ell$. Generate at most $n+1$ fresh independent
uniform $b\in K^n$. For each, compute
\[
 y=g(T)b
\]
by Horner's rule. Verify $My=0$ exactly, and retain $y$ only if it
is linearly independent of previously retained verified vectors.
Stop and accept the trial when $\widehat r$ independent verified
vectors have been obtained. If $\widehat r=0$, accept with the empty
basis immediately. If the cap is exhausted, discard the trial and
start another independent one.

Acceptance is sound on \emph{every} random choice. The verified
independent vectors give $\rank M\le n-\widehat r=\ell$;
Lemma~\ref{fk:lem:lower} gives the reverse inequality. Thus acceptance
certifies $\ell=k$ and a complete extension-field kernel basis.
An underestimated rank cannot pass: it asks for more independent
kernel vectors than exist. A bad polynomial action cannot pass
without the explicit kernel and independence checks.

\subsection{Good-trial probability, probes, and expected work}
Conditional on~\eqref{fk:eq:cyc}--\eqref{fk:eq:nil} being good, random
left and right probes retain the minimum polynomial of the invertible
cyclic $k$-dimensional restriction with failure probability at most
$2k/q$. This follows from the nonzero degree-$k$ controllability and
observability determinant polynomials and Schwartz--Zippel.
The restriction's projected scalar minimum polynomial is determined
exactly by the first $2n$ terms, since its degree is at most $n$;
this is the scalar Wiedemann recurrence fact \cite[Section II]{Wiedemann}.
Hence on the combined good event $\mathcal G$,
\[
 \ell=k,\quad g(0)\ne0,\quad
 g(T)=g(0)P_{\ker T\parallel\operatorname{im}T},
\]
and
\begin{equation}\label{fk:eq:bad}
 \Prob(\mathcal G^c)
 \le\frac{n(n-1)/2+r}{q-1}+\frac{2k}{q}
 =O(n^2/q).
\end{equation}
For $k=0$ the statement means $g=1$ and the projection is the identity.

Thus each $g(T)b$ on a good trial is an independent uniform vector
in the $r$-dimensional kernel (up to a nonzero scalar).
The expected number of draws to span it, without a cap, is
\[
 \sum_{j=0}^{r-1}\frac1{1-q^{j-r}}\le r+2,
\]
with zero draws when $r=0$.
After $n+1$ draws, the failure probability is at most
\begin{equation}\label{fk:eq:cap}
 \frac{q^r-1}{q-1}\,q^{-(n+1)}
 \le\frac1{q(q-1)}.
\end{equation}
Indeed a proper span is contained in a hyperplane, and there are
$(q^r-1)/(q-1)$ such hyperplanes, all samples lying in a given one with probability
$q^{-(n+1)}$. Our field size makes each trial succeed with
probability at least $1/2$, uniformly over fixed $M$.

\paragraph{Charging an entire trial.}
The sequence, dot products, and recurrence require
$O(n(w+n))$ field operations, $2n$ applications of $T$, and
$O(n)$ temporary field storage. One polynomial action requires
$O(n(w+n))$ field operations. Kernel verification costs $O(w)$
per candidate. Maintaining independent verified columns costs
$O(nj)$ operations when $j$ columns have already been retained.
On a good trial these costs total in expectation
\[
 O\bigl(n(w+n)(r+1)+nr(r+1)\bigr)
 =O\bigl(n(w+n)(r+1)\bigr).
\]
Even on a bad trial the cap gives the deterministic work bound
$O(n^2(w+n))$. By~\eqref{fk:eq:bad}, its expected contribution is
$O(n^4(w+n)/q)$, absorbed by the displayed good-trial bound.
At all times there are at most $r$ retained vectors, since they
are verified and independent, even on bad trials.

Independent retries satisfy the renewal identity
$\E W=\E W_{\mathrm{trial}}/\Prob(\mathrm{accept})$.
They therefore terminate almost surely with the asserted expected
field work on \emph{every fixed matrix}, including all exceptional
graphs. No rational fallback or high-probability corank-one event
is hidden in this Las Vegas assertion.

\subsection{Descent to the base field and deterministic verification}
Let $B\in K^{n\times r}$ be an accepted full kernel basis.
Find row indices $J$ with $B_{J,:}$ invertible and compute
\[
 C=B(B_{J,:})^{-1},\qquad C_{J,:}=I_r.
\]
This costs $O(nr^2+r^3)$ field operations, bounded by
$O(n(w+n)(r+1))$, and $O(nr)$ additional storage.

\begin{lemma}[Canonical normalization descends]\label{fk:lem:descent}
Every entry of $C$ belongs to the embedded base field $\F_p$.
\end{lemma}
\begin{proof}
The kernel is defined by $M$ with base-field coefficients, so applying
the coefficientwise Frobenius $x\mapsto x^p$ to $C$ yields another
kernel basis with the same identity on $J$. Restriction to these
pivot coordinates is a bijection from the kernel to $K^r$.
Consequently this normalized basis is unique, and $C^{(p)}=C$.
The fixed elements of this Frobenius in $K$ are precisely $\F_p$.
\end{proof}

In the polynomial representation, check that all positive-degree
coefficients of $C$ vanish, read off the base-field entries, and
verify $MC=0$ in $\F_p$. These are exact checks, not an assumption
that arbitrary extension-field vectors can simply be enumerated as
base-field vectors.

The certificate comprises the irreducible polynomial $h$, $D,u,v$,
$\ell$, the base-field matrix $C$, and its pivot rows $J$.
A verifier checks the representation sizes and bounds on $q$, tests
irreducibility, recomputes~\eqref{fk:eq:seq} and its exact BM length,
and checks
\[
 r=n-\ell,\qquad MC=0,\qquad C_{J,:}=I_r.
\]
These give the lower and upper rank bounds deterministically, without
trusting the random preconditioner or the polynomial projection.
Verification needs $O(n(w+n)+wr+r^2)$ field/base operations, within
the theorem's bound. It does not recompute a dense elimination.

\subsection{Field construction, bit work, and storage}
The extension degree is $s=O_p(\log(n+1))$. Draw monic degree-$s$
polynomials over $\F_p$ and test irreducibility deterministically.
The standard irreducible-polynomial count gives success probability
at least $1/(2s)$, so the expected number of draws is $O(s)$.
One may use Rabin's exact criterion \cite{Rabin}, or the elementary
version checking
\[
 \gcd\bigl(h,X^{p^i}-X\bmod h\bigr)=1
 \quad\text{for all }1\le i\le\lfloor s/2\rfloor.
\]
A reducible degree-$s$ polynomial has an irreducible factor of degree
at most $s/2$, so this version is sufficient and necessary.
Successive modular Frobenius powers and classical polynomial gcd
give polynomial work in $s$, hence polylogarithmic bit work in $n$.
Every accepted field construction is correct; randomness only
controls its running time.

Represent elements by $s$ base-field coefficients. Addition,
multiplication modulo $h$, and inversion by extended Euclid each
cost a polynomial in $s$, hence $\widetilde O_p(1)$ bit operations.
Independent uniform field choices use $O(\log q)$ random bits in
expectation per element; rejecting zero for diagonal entries has
bounded expected overhead. This charges all field construction,
randomness, polynomial actions, and basis arithmetic in
Theorem~\ref{fk:thm:kernel}. Workspace is $O(n(r+1)+w)$ field/base
elements; each has $O_p(\log(n+1))$ bits. Retry-counter storage can
be treated separately and has finite expected size.

\subsection{Exact partition listing from a single corank moment}
Let $M_{\mathbb Z}=A-\lambda I$ for a simple graph of maximum degree
$d$ and a known integer $|\lambda|\le d$, and let $M$ be its reduction modulo
the fixed odd prime $p$. Compute the full base-field kernel by
Theorem~\ref{fk:thm:kernel}. Enumerate all $p^r$ coefficient vectors
$a\in\F_p^r$ and their images $Ca$.
Retain only those with last coordinate $1$ and every coordinate in
$\{1,-1\}$. Lift these residues uniquely to integer signs and
check exactly
\[
 \one^T\sigma=0,\qquad M_{\mathbb Z}\sigma=0.
\]
Return every passing sign vector, with the orientation $\sigma_n=1$.
The list is complete on every input: every true balanced integer
sign kernel vector reduces to a vector of this finite-field kernel.
Its integer verification eliminates all modular artifacts.

\paragraph{The enumeration factor is charged without an extra $r$.}
Use a base-$p$ odometer or Gray ordering and maintain the current
image $Ca$. In the full run, the number of digit changes is at most
$p^r\sum_{j\ge0}p^{-j}=O_p(p^r)$.
Each digit change adds a known basis column in $O(n)$ base operations.
All sign, balance, and integer degree checks cost $O_d(n)$ per
candidate. The full enumeration therefore costs
\[
 \widetilde O_{d,p}(n p^r)
\]
bit operations, including the list entries. Recomputing every image
by a fresh $nr$ multiplication would be an invalid expectation
shortcut: $\E p^r<\infty$ need not bound $\E(rp^r)$.

\begin{corollary}[Moment-to-full-list theorem]\label{fk:cor:moment}
Suppose a bounded-degree random-graph model satisfies, for one fixed
odd prime $p$,
\[
 \sup_n\E\bigl[p^{\dim\ker_{\F_p}(A-\lambda I)}\bigr]\le C<\infty.
\]
Then the preceding algorithm lists all balanced integer sign kernels
of $A-\lambda I$, up to global exchange, in unconditional expected
$\widetilde O_{d,p,C}(n^2)$ bit work. It has no algorithmic error
and does not assume real or modular corank one. Its expected working
storage is $\widetilde O_{d,p,C}(n)$ bits and its expected output
size is $O_{d,p,C}(n)$ bits.
\end{corollary}
\begin{proof}
The per-input conditional expected work is
\[
 \widetilde O_{d,p}\bigl(n^2(r+1)+n p^r\bigr).
\]
Since $r+1\le O_p(p^r)$, the stated raw moment bounds its expectation
by $\widetilde O(n^2)$ and gives the storage/output bounds as well.
This averages a valid all-input Las Vegas bound, so exceptional
graphs incur no unjustified dense fallback cost.
\end{proof}


\section{Certified vector intertwiners and fast unknown-base recovery}
\label{sec:vector-algorithm}
The symmetric full-kernel algorithm also handles nonsymmetric quotient
matrices. The reduction computes a doubled kernel, then projects to its
right component before enumeration. Its expected-work application needs
one bounded moment for that right kernel in the actual graph law.

\subsection{A complete right kernel on every input}
Fix an odd prime $p$. Let $A=A^T$ be an integer adjacency matrix of a
simple graph of maximum degree at most $D$, and let $R$ be a fixed
$r\times r$ integer matrix with bounded entries. Stack matrix columns
in the convention $\operatorname{vec}(X)=(X_{:,1},\ldots,X_{:,r})$.
Over $\Fp$, put
\begin{equation}\label{eq:vector-operator}
 \mathcal K_R=I_r\otimes A-R^T\otimes I_n,\qquad
 m=rn,\quad w=\operatorname{nnz}(\mathcal K_R),\quad
 h=\nullity_{\Fp}\mathcal K_R.
\end{equation}
Then $\mathcal K_R\operatorname{vec}(X)=\operatorname{vec}(AX-XR)$;
its transpose acts as $\operatorname{vec}(Y)\mapsto
\operatorname{vec}(AY-YR^T)$. Each product costs
$O(r\operatorname{nnz}(A)+n\operatorname{nnz}(R))$ field operations.
At fixed $r,D$ and bounded quotient entries, $m=O(n)$ and $w=O(n)$.
The structured products or an explicit sparse representation suffice.

\begin{theorem}[Certified right kernels]\label{vec:thm:kernel}
For every fixed input in \eqref{eq:vector-operator}, a Las Vegas
algorithm returns the complete right-kernel basis
$C\in\Fp^{m\times h}$, normalized on pivot rows $J$ by $C_{J,:}=I_h$.
Its conditional expected bit work is
\[
 \widetilde O_p\bigl(m(w+m)(h+1)\bigr),
\]
and its working storage is
$\widetilde O_p(m(h+1)+w)$ bits, apart from a retry counter.
A certificate of $O_p(mh+m\log m)$ bits is deterministically
checkable in $\widetilde O_p(m(w+m))$ bit work.
Every accepted answer is correct, including on inputs with large
nullity. No symmetry assumption on $R$ is imposed.
\end{theorem}
\begin{proof}
Apply Theorem~\ref{fk:thm:kernel} to the symmetric matrix
\begin{equation}\label{eq:vector-doubling}
 \mathcal S_R=
 \begin{pmatrix}0&\mathcal K_R\\\mathcal K_R^T&0\end{pmatrix}
 \quad\text{of order }2m.
\end{equation}
For $(u,v)$ it acts as $(\mathcal K_Rv,\mathcal K_R^Tu)$, so
\[
 \ker\mathcal S_R=\ker\mathcal K_R^T\oplus\ker\mathcal K_R,
 \qquad \rank\mathcal S_R=2\rank\mathcal K_R,
 \qquad \nullity\mathcal S_R=2h.
\]
These are field identities; no Gram-matrix or positivity argument is
used. One product uses one $\mathcal K_R$ and one $\mathcal K_R^T$
product. The per-input theorem, with order $2m$, sparsity $2w$, and
nullity $2h$, has the displayed work and storage bounds. Its verified
extension-field construction, preconditioning retries and base-field
descent are all included.

The bottom $m$ rows of a complete $\mathcal S_R$-kernel basis span
$\ker\mathcal K_R$: every right nullvector $v$ occurs as $(0,v)$.
Select $h$ independent projected columns by exact elimination, choose
$h$ invertible pivot rows, and normalize there. Arbitrary first columns
need not be independent. Selection and normalization cost
$O(mh^2+h^3)\le O(m^2(h+1))$, within the theorem's work bound.
They use $O(mh)$ storage. The symmetric primitive has already returned
a base-field basis. Equivalently, a projected extension-field basis
can be normalized first: Frobenius preserves the base-defined right
kernel and its uniquely pivot-normalized basis, so that basis lies in
$\Fp$. Extract its base coefficients and check $\mathcal K_RC=0$.
The cases $h=0$ and $h=m$ require no exception to these arguments.

For a compact final certificate, retain the symmetric primitive's
verified field definition, nonzero diagonal preconditioner $\mathcal D$,
and scalar probes $u,v$, together with $C,J$. Its checker regenerates
the $4m$ terms
\[
 u^T(\mathcal D\mathcal S_R)^{j+1}v,
 \qquad 0\le j<4m,
\]
and recomputes their true shortest Berlekamp--Massey length $\ell$.
Require $0\le h\le m$, the prescribed dimensions, distinct pivot
rows $J$, and canonical base-field entries of $C$. Require
$\ell=2(m-h)$ and check $\mathcal K_RC=0$ and $C_{J,:}=I_h$.
The unconditional shifted-sequence bound of Lemma~\ref{fk:lem:lower}
gives $\ell\le\rank\mathcal S_R=2\rank\mathcal K_R$ on every draw.
The verified independent columns give $\rank\mathcal K_R\le m-h$.
Thus acceptance forces
\[
 m-h=\ell/2\le\rank\mathcal K_R\le m-h,
\]
certifying completeness with zero verification error. This is an
integer dimension comparison, not division by two in the field.
The verifier checks field sizes and encodings before the primitive's
irreducibility and sequence checks. Its $O(m)$ sparse products,
quadratic recurrence computation, and residual/pivot checks cost
$\widetilde O_p(m(w+m))$. Retaining the right basis and probes gives
the stated certificate size; the unused left basis is unnecessary.
\end{proof}

\subsection{Enumeration under a single actual-law moment}
\begin{theorem}[Moment-to-vector-list theorem]\label{vec:thm:moment}
Fix $K,D$ and a finite family $\mathcal R$ of nonnegative integer
quotient matrices of orders at most $K$, with row sums at most $D$.
Suppose a law on simple graphs of maximum degree at most $D$ satisfies,
for one fixed odd prime $p$ and every $R\in\mathcal R$,
\begin{equation}\label{eq:vector-moment}
 \sup_n\E p^{h_R(A)}\le C_R<\infty,\qquad
 h_R(A)=\nullity_{\Fp}(I_r\otimes A-R^T\otimes I_n).
\end{equation}
There is an always-correct Las Vegas algorithm which lists all
nonempty labeled equitable partitions with quotient in $\mathcal R$.
Its unconditional expected bit work is $\widetilde O(n^2)$, its
expected working storage is $\widetilde O(n)$ bits, and its expected
output size is $O(n)$ bits. Constants may depend on
$p,K,D,\mathcal R$ and the moment bounds. The algorithm is complete
on every input; the moment assumption controls its expected cost.
\end{theorem}
\begin{proof}
For each $R$, compute the certified right basis $C$ from
Theorem~\ref{vec:thm:kernel}. Enumerate only its $p^h$ base-field
coefficient vectors. Maintain the current vector using a base-$p$
odometer or a Gray order. Across all states there are $O_p(p^h)$
digit changes, each adding a scalar multiple of one column of $C$
in $O(m)$ field operations. Total generation costs
$\widetilde O_p(mp^h)$ bit work, with no extra factor $h$.
The doubled kernel is never enumerated.

Reshape each vector into $X\in\Fp^{n\times r}$. Require that every
row has exactly one entry equal to one and all remaining entries
zero, and that every column is nonzero. Lift these rows to integer
zero-one rows and check $AX=XR$ over the integers. Check any class
sizes actually prescribed by the task. A blind unknown-frame search
has no such size promise and imposes only nonemptiness. At fixed
$K,D$, these scans and sparse integer neighbor-count checks cost
$\widetilde O(n)$ per candidate. Every true integer partition reduces
to a field-kernel vector, and the integer check rejects modular
artifacts. Hence the resulting list is exact on every input.

For a fixed input and target the conditional expected work is
\begin{equation}\label{eq:vector-per-input}
 \widetilde O_p\bigl(m(w+m)(h+1)+mp^h\bigr).
\end{equation}
Because $h+1\le p^h$ and $m,w=O_{K,D}(n)$, assumption
\eqref{eq:vector-moment} bounds its graph expectation by
$\widetilde O(n^2)$. This averages the valid per-input Las Vegas
bound, without excluding rare large kernels. Neither
$\E p^{2h}$ nor $\E[h p^h]$ is used.
The same bound gives expected $\widetilde O(n)$ working storage.
Stream the output: at most $p^h$ candidates pass, and a partition
uses $O_K(n)$ bits. Its expected output size is therefore $O(n)$.
Finite summation over $\mathcal R$ preserves all these bounds;
independence of different runs is unnecessary. Reuse their workspace.
\end{proof}

For unknown quotients, use every order $1\le r\le K$ and all
nonnegative integer matrices with row sums at most $D$. There are
at most $\sum_{r=1}^K(D+1)^{r^2}$ candidates. Retain a verified
partition with the maximum number of classes, using only $O(n)$
additional bits. On the rigidity event of
Theorem~\ref{thm:unknown-frame}, every candidate coarsens the plant,
so this maximum is the planted partition, up to class names.
Completeness and zero algorithmic error hold outside that event too;
statistical recovery there is a separate question.

Nonsymmetric quotient candidates must remain in this blind search.
For any accepted one-hot $X$ with class sizes $N_i>0$, symmetry of
$A$ gives the integer balance identities
\[
 X^TAX=\diag(N_i)R=(\diag(N_i)R)^T,
 \qquad N_iR_{ij}=N_jR_{ji}.
\]
Equal positive sizes, when prescribed, force $R=R^T$. Equal sizes
of the true plant do not force its permitted coarsenings to have
equal sizes, so they do not justify discarding nonsymmetric candidates.
There is no division by a class size, including when $p\mid N_i$.
For a fixed $R$, only its allowed label automorphisms preserve that
target. To list unlabeled partitions in the unknown-$R$ search,
canonicalize $(X,R)$ simultaneously under the at most $K!$ label
permutations and output only the canonical representative. This
uses a constant number of scans and no growing deduplication table.

\subsection{Unknown certified frames in the actual simple law}
\begin{corollary}[Nearly quadratic unknown-frame recovery]
\label{vec:cor:frame}
Under the omission certificate of
Theorem~\ref{thm:unknown-frame}, the algorithm given only $A,K,D$
lists all equitable partitions with at most $K$ classes in
unconditional expected $\widetilde O_{K,D}(n^2)$ bit work, expected
$\widetilde O_{K,D}(n)$ working bits, and expected $O_{K,D}(n)$
output bits. Selecting a largest-class partition recovers the plant
and its unknown frame with failure $O_{K,D}(n^{-(d_{\min}-2)})$.
Varying total degrees, zero omission weights, internal degrees zero or
one, unequal planted sizes, and all
admissible sizes at which the simple factors exist are included.
\end{corollary}
\begin{proof}
Connectedness and integer balance fix the planted size ratios.
They can be written $N_i=c_iN$ for a primitive positive integer vector
$(c_i)$; this is only a parametrization of all admissible sizes.
Theorem~\ref{thm:frame-bounded-moment} gives
\eqref{eq:vector-moment} in exactly the independent uniform-simple
factor law, for $p=3$ and every reduction of every candidate $R$.
There are finitely many permitted planted frames, size-ratio vectors,
and candidate targets at fixed $K,D$, so a common moment constant
exists. Apply Theorem~\ref{vec:thm:moment} and the structural failure
bound of Theorem~\ref{thm:unknown-frame}.
\end{proof}

There is also a deterministic polynomial truncation preserving the
structural failure rate. Over $\F_2$, use dense elimination for each
candidate and cap its kernel size at $n^{d_{\min}-1}$, reading the minimum degree from the input. The bounded binary moment and a finite union give
interruption probability $O_{K,D}(n^{-(d_{\min}-1)})$.
If any cap is exceeded, return abstention rather than a partial maximum.
Gray enumeration and
integer checks then use bounded
$\widetilde O_{K,D}(n^3+n^{d_{\min}})$ bit work and $O_K(n^2)$ space.
Adding the structural failure preserves $O(n^{-(d_{\min}-2)})$.
The expected nearly quadratic Las Vegas algorithm has no such cap.

\subsection{The random-lift specialization}
\begin{corollary}[Nearly quadratic unknown-base recovery]
\label{vec:cor:lift}
Let $H$ be a fixed connected simple $d$-regular graph, $d\ge3$,
with $k\le K$ vertices, and let $A_N$ be its random $N$-lift,
with independent uniform matchings above all base edges. Given only
the lifted graph and fixed bounds $K,D\ge d$, the algorithm recovers
the base fibers and $H$ up to isomorphism with probability $1-o(1)$.
Its unconditional expected bit work is
$\widetilde O_{H,K,D}(n^2)$, where $n=kN$, with expected
$\widetilde O_{H,K,D}(n)$ working bits. Listing all equitable
partitions with at most $K$ classes has expected $O_{H,K,D}(n)$
output bits. All positive lift sizes are covered, including bipartite
bases; no real-corank or eigengap promise is needed.
\end{corollary}
\begin{proof}
Choose one fixed odd prime, for example $p=3$. The bounded lift
moment theorem, Theorem~\ref{thm:lift-bounded-moment}, states that
for this actual unconditioned lift law, every fixed $r$ and target
$\overline R\in\Fp^{r\times r}$ satisfy
\[
 \sup_N\E\bigl|\{Y\in\Fp^{n\times r}:A_NY=Y\overline R\}\bigr|
 <\infty.
\]
Apply it to the reductions of the finite integer quotient family
above. Its solution count is exactly $p^{h_R(A_N)}$, so
Theorem~\ref{vec:thm:moment} gives the work, space and output bounds.
A lift of a simple base is automatically simple; there is no
conditioning or change of graph law. Both $H$ and $R$ are analysis
parameters in this moment invocation, not information supplied to
the algorithm. Corollary~\ref{cor:random-lift} supplies the rigidity
event, and maximum-class selection recovers the plant on that event.
\end{proof}

The omission certificate is sufficient; minimum degree three alone is
not substituted for it. The nearly quadratic work is an expected bound,
not a worst-case quadratic guarantee.

\section{Old and new spectra of nonbipartite regular lifts}
\label{sec:lift-spectral}
Let $B$ be the adjacency matrix of a fixed connected simple
nonbipartite $d$-regular base on $k$ vertices, with $d\ge3$.
Let $A_N$ be its random $N$-lift, with independent uniform matchings
above the base edges. The fiber-constant embedding $J$ satisfies
$A_NJ=JB$. Over the reals put $U_N=\im J$ and $W_N=U_N^\perp$.
Both spaces are invariant because $A_N$ is symmetric.

\begin{theorem}[Old/new spectral disjointness]\label{lift:thm:old-new}
For every such fixed base,
\[
 \Prob\{\spec(A_N|_{W_N})\cap\spec(B)\ne\varnothing\}
 \longrightarrow0
\]
as $N\to\infty$ through all positive integers. Thus every old
eigenvalue retains exactly its base multiplicity with high probability,
including irrational old eigenvalues. This statement concerns exact
equality of old and new eigenvalues; it asserts neither simplicity of
the new spectrum nor a numerical gap around the old eigenvalues.
\end{theorem}

The proof sharpens the full-uniform coefficient in the lift moment
calculation for cyclic targets, assembles an exact moment for
irreducible companions, and transfers real collisions through integer
polynomial kernels. All finite-field limits keep the field fixed.

\subsection{The exact uniform coefficient for a cyclic target}
Fix an odd prime $p$ and $r\ge1$, write $V=\Fp^{1\times r}$ and $q=p^r$,
and let $R\in\Fp^{r\times r}$. A cyclic target has a row $v$ for which
$v,vR,\ldots,vR^{r-1}$ form a basis. Labels are rows, so a legal star at old
vertex $i$ obeys $\sum_{j\sim i}b_j=aR$. Define
\[
 h_B(R)=\dim_{\Fp}\{X\in\Fp^{k\times r}:BX=XR\}.
\]
We use the exact star representation, entropy phases, sparse bound,
and finite-support relative local limit estimate established in the
channel moment proof, with positive weights $q_{ij}=1/d$ and
Corollary~\ref{ch:cor:positive-phases}. The following argument
retains the exact compatible-lattice fraction in that representation.

\begin{lemma}[Cyclic uniform coefficient]\label{lift:lem:cyclic}
If $R$ is cyclic, the contribution from a sufficiently small fixed
full-uniform edge-table neighborhood to
$\E|\{Y:A_NY=YR\}|$ is $p^{h_B(R)}+o(1)$.
This holds for every lift-size sequence, without modular
semisimplicity assumptions on $B$ or $A_N$.
\end{lemma}
\begin{proof}
Let $P$ be the uniform legal-star law, with the root and any $d-1$
neighbors independent uniform, and let $Q$ make the root and all
$d$ neighbors independent uniform. Take $N$ independent stars at
each old vertex. Their directed tables $m_{ij}(a,b)$ have common
root margins $n_i(a)$ for all slots $j\sim i$. Write
\[
 \mathcal W(m)=\one_{\{m_{ji}=m_{ij}^T\}}
       \prod_{ij\in E(B)}\frac{\prod_{a,b}m_{ij}(a,b)!}{N!}.
\]
Exact ordered-stub matching enumeration gives
\begin{equation}\label{eq:lift-two-normalizers}
 \E|\{Y:A_NY=YR\}|=q^{dkN}\E_P\mathcal W(M),
 \qquad q^{dkN}\E_Q\mathcal W(M)=1.
\end{equation}
For the second identity, every root assignment and every lift determine
one consistent array of neighbor labels. There are $q^{kN}$ root
assignments, whereas $Q$ has $q^{(d+1)kN}$ equally likely arrays.

Let $L_0$ be the full integer direction lattice of directed tables
with common row margins and total zero in each table. Independent
stars generate $L_0$. Legal stars generate precisely its sublattice
\begin{equation}\label{eq:lift-congruences}
 L=\{m\in L_0:C_i(a;m)=0\text{ in }V\},\quad
 C_i(a;m)=\sum_{j\sim i,b}b\,m_{ij}(a,b)-(aR)n_i(a).
\end{equation}
Here $[L_0:L]=p^{kqr}$. To recall maximality, a circle-valued
character constant on legal stars has slot functions
$\theta_j(a,b)$. At fixed $a$, freely varying two neighbors with a
third compensating forces an additive Cauchy equation, because
$d\ge3$. Thus $\theta_j(a,b)=\alpha_j(a)+\chi_a(b)$ with the
same additive character in every slot. Varying the root leaves,
modulo characters of the common-row lattice, exactly the characters
of $C_i(a;m)$. They are independent: changing one neighbor at one
root label changes that congruence by any element of $V$.
There are $q^q$ choices per old vertex, giving the stated full index.

The two laws have the same positive covariance on the real span
of $L_0$: any two neighbors under $P$ are independent uniform
conditional on the root, as they are under $Q$. The tilted relative
lattice local limit estimate therefore gives, uniformly in the
$\sqrt N\log N$ window and on the legal affine lattice,
\begin{equation}\label{eq:lift-relative-ratio}
 \Prob_P(M=m)=p^{kqr}(1+o(1))\Prob_Q(M=m).
\end{equation}
Indeed one tilts to mean $m/N$, applies Fourier inversion on each
maximal-lattice torus, and undoes the tilt. The common covariance
gives a common quadratic rate; the remaining exponent error is
$O(\log^3N/\sqrt N)$. The lattice covolume ratio is exactly
$[L_0:L]$. At other lattice points the legal probability is zero.

We compute the index after imposing compatibility. Let $\mathcal T$
be the compatible real direction space and $K_0=L_0\cap\mathcal T$.
Its integer directions have a free integral parametrization: choose
arbitrary integer margins $u_i(a)$ of total zero, and arbitrary
entries $m_{ij}(a,b)$ for $a,b\ne0$ on every oriented base edge.
Complete that table by
\begin{align*}
 m(a,0)&=u_i(a)-\sum_{b\ne0}m(a,b),&&a\ne0,\\
 m(0,b)&=u_j(b)-\sum_{a\ne0}m(a,b),&&b\ne0,\\
 m(0,0)&=u_i(0)+u_j(0)+\sum_{a,b\ne0}m(a,b).
\end{align*}
Use its transpose in the reverse slot. Hence reduction of $K_0$
modulo $p$ is the entire compatible field direction space.
If $\rho$ is the field rank of the restricted congruence map, then
\begin{equation}\label{eq:lift-compatible-index}
 K=K_0\cap L,\qquad [K_0:K]=p^\rho.
\end{equation}
No saturation factor or division by $N$, $q$, or $d$ is involved.

A left dependency of the congruences is given by columns $h_i(a)$,
paired with the row vectors $C_i(a;m)$. Testing a zero-margin
two-by-two transfer on edge $ij$ gives
\[
 (b-b')[h_i(a)-h_i(a')]
 +(a-a')[h_j(b)-h_j(b')]=0.
\]
Set $a'=b'=0$ and test coordinate rows $b$. Every coordinate of
$h_i(a)-h_i(0)$ is then linear in $a$. Thus
$h_i(a)=K_i a^T+c_i$, and the same equation gives
$K_j=-K_i^T$. Connectivity and an odd cycle force a common
$K_i=K$ with $K^T=-K$. Substitution in the dependency functional
leaves
\[
 \sum_{i,a}\left[a\left(\sum_{j\sim i}c_j-Rc_i\right)
                         -aRK a^T\right]u_i(a).
\]
Margins vary freely subject only to their own total zero; the bracket
must be constant in $a$ and is zero at $a=0$. Testing $a$ and $-a$
in odd characteristic separates its linear and quadratic parts:
\begin{equation}\label{eq:lift-dependencies}
 \sum_{j\sim i}c_j=Rc_i,\qquad K^T=-K,\qquad RK=KR^T.
\end{equation}
These conditions are also sufficient, so no dependency has been lost.
Writing the rows $c_i^T$ as $X$ gives $BX=XR^T$.
This Sylvester operator is the Frobenius adjoint of $Y\mapsto BY-YR$;
their kernel dimensions agree because $B=B^T$. The linear part
therefore has dimension $h_B(R)$.

The remaining $K$ vanishes for cyclic $R$. Define the alternating
form $\beta(a,b)=aKb^T$ on rows. The last condition in
\eqref{eq:lift-dependencies} makes right multiplication by $R$
self-adjoint for this form. For a cyclic row $v$, all rows have
the form $vf(R)$, and
\[
 \beta(vf(R),vg(R))=\beta(v,v g(R)f(R)).
\]
This is symmetric in $f,g$. The form is both symmetric and
alternating, so it is zero in odd characteristic. This also permits
degenerate forms. Consequently
\begin{equation}\label{eq:lift-compatible-rank}
 \rho=kqr-h_B(R).
\end{equation}

We now retain the exact Gaussian fraction. The all-zero legal star
gives a compatible point $v_0$ in the one-step affine support.
The compatible affine lattice hulls are $Nv_0+K$ and $Nv_0+K_0$,
nonempty for every $N$. For compatible centered scaled tables write
\[
 X_{ij}(a,b)=\frac{u_i(a)+u_j(b)}q+Z_{ij}(a,b),
\]
where $u_i$ has total zero and $Z_{ij}$ has zero row and column sums.
The common covariance density and matching Stirling weight give
the Gaussian exponent
\begin{equation}\label{eq:lift-weighted-precision}
 -\frac q2\sum_i\|u_i\|^2
       -\frac{q^2}2\sum_{ij\in E(B)}\|Z_{ij}\|^2.
\end{equation}
Its negative is positive definite on compatible directions.
The directed and compatible dimensions differ by
$|E(B)|(q^2-1)$, exactly twice the matching factorial power of $N$.
Thus all powers cancel. Translated Gaussian Riemann sums on the
two fixed lattices have ratio $[K_0:K]^{-1}=p^{-\rho}$, uniformly
in the affine translate. The relative local limit estimate also
ensures actual support at every central affine-lattice point for
sufficiently large $N$; equality of finite support and its lattice
hull is not assumed.

The normalized $Q$ central matching mass tends to one. Its weighted
law in \eqref{eq:lift-two-normalizers} is exactly a random lift
with independent uniform root labels. On each edge its matched
label pairs are iid uniform, so Chebyshev and the finite edge union
put all tables in the logarithmic window with probability $1-o(1)$.
Multiplying its legal-lattice fraction $p^{-\rho}$ by
\eqref{eq:lift-relative-ratio} gives the coefficient
$p^{kqr-\rho}=p^{h_B(R)}$.

Finally this window covers the whole sufficiently small fixed
uniform neighborhood. The finite-support rate function is analytic
near its mean. Subtracting the matching factorial Hessian gives
exactly the positive form \eqref{eq:lift-weighted-precision}; shrink
the neighborhood to absorb cubic terms. Outside the logarithmic
window the remaining weighted bound is a fixed polynomial times
$\exp(-c\log^2N)$, so its sum is $o(1)$. This proves the lemma
with every lattice residue and prefactor retained.
\end{proof}

\subsection{Exact irreducible-companion moments}
\begin{theorem}\label{lift:thm:irreducible-moment}
Let $g\in\Fp[t]$ be monic irreducible of degree $r\ge1$ and let
$R_g$ be a companion matrix. For every fixed odd $p$,
\begin{equation}\label{eq:lift-irreducible-moment}
 \E|\{Y:A_NY=YR_g\}|=2p^{h_B(g)}+o_{B,p,g}(1),\qquad
 h_B(g)=\dim_{\Fp}\ker g(B),
\end{equation}
along all $N$. Modular diagonalizability is unnecessary.
\end{theorem}
\begin{proof}
For any matrix $M$, evaluation at a cyclic column $v$ of $R_g$
gives an isomorphism
\begin{equation}\label{eq:lift-hom-kernel}
 \{Y:MY=YR_g\}\longrightarrow\ker g(M),\qquad Y\mapsto Yv.
\end{equation}
The intertwiner must send $R_g^jv$ to $M^jYv$; conversely these
relations define it from any $w\in\ker g(M)$, since the last cyclic
relation is $g(M)w=0$. This proves the dimension identity in the
theorem without semisimplicity.

The positive-weight phase classification of
Corollary~\ref{ch:cor:positive-phases} gives common $R_g$-invariant
subspaces and quotient coset point maps on a connected nonbipartite
base. An irreducible companion has only the invariant subspaces
$0$ and $V$. The former gives precisely the $p^{h_B(g)}$ old
point intertwiners; the latter gives one full-uniform phase.
Every point map lifts deterministically. The channel boundary estimate
with zero restricted dimension makes each small support neighborhood
contribute $1+o(1)$.

Different point maps differ on at least one full fiber. Choose
support radii below $N/2$ and then small enough to separate all
point neighborhoods from the uniform edge-table neighborhood.
They are disjoint and contribute $p^{h_B(g)}+o(1)$ altogether.
Lemma~\ref{lift:lem:cyclic} gives the same contribution at full
uniformity, since every companion is cyclic.
For completeness, each exact legal-star type contributes
\[
 \prod_i\frac{N!}{\prod_s n_i(s)!}
 \prod_{ij\in E(B)}\frac{\prod_{a,b}m_{ij}(a,b)!}{N!}.
\]
Its logarithm is $N\Phi+O(\log(N+1))$, uniformly at zero entries,
where
\[
 \Phi=\sum_iH(P_i)-\sum_{ij\in E(B)}H(\nu_{ij})
\]
for the star laws and directed-edge laws.
The entropy is nonpositive and has only the phases just listed.
The complementary feasible type region is compact with strictly
negative maximum. Polynomially many types and the uniform factorial
bounds make its contribution exponentially small. Combining the
three regions proves \eqref{eq:lift-irreducible-moment}.
\end{proof}

The space in \eqref{eq:lift-hom-kernel} always has dimension divisible
by $r$, because $\ker g(M)$ is a vector space over $\Fp[t]/(g)$.
Also $J$ remains injective over every field and intertwines $B$ with
$A_N$, even when $p\mid N$. Hence its dimension $h_N(g)$ in the
lift is always at least $h_B(g)$. Subtracting this deterministic
baseline before Markov gives
\[
 \Prob\{h_N(g)\ge h_B(g)+r\}
 \le\frac{\E p^{h_N(g)}-p^{h_B(g)}}
              {p^{h_B(g)}(p^r-1)}.
\]
Consequently \eqref{eq:lift-irreducible-moment} yields
\begin{equation}\label{eq:lift-shifted-markov}
 \limsup_N\Prob\{h_N(g)\ge h_B(g)+r\}
 \le\frac1{p^r-1}.
\end{equation}
This is a bound on total modular kernel dimension; a modular old/new
direct-sum decomposition is not required.

\subsection{Good primes and transfer of real collisions}
Factor the fixed base characteristic polynomial over the integers:
\[
 \chi_B(t)=\prod_\alpha f_\alpha(t)^{t_\alpha},\qquad
 F(t)=\prod_\alpha f_\alpha(t),
\]
with distinct monic irreducible factors. Real symmetry gives
$F(B)=0$ and makes $F$ the squarefree minimal polynomial.
Its integer discriminant is nonzero. Call $p$ good if it is odd
and $p\nmid\operatorname{disc}(F)$; only finitely many primes are
excluded. The reduction of $F$ is squarefree, so its distinct
factors remain coprime and $B\bmod p$ is semisimple.

Fix $f=f_\alpha$ of degree $\ell$ and old exponent $t=t_\alpha$.
At any good prime write
\[
 \overline f=\prod_j g_j,\qquad \deg g_j=r_j,\qquad
 \sum_jr_j=\ell.
\]
Good-base semisimplicity and the characteristic polynomial identity
give $\dim\ker g_j(B)=r_jt$. No prescribed splitting pattern of
the primes is needed. The discriminant used here is that of $F$;
the characteristic polynomial itself can have repeated old roots.

Let $S$ be the fiber-sum map. Over $\Q$ the exact identities
$A_NJ=JB$, $SA_N=BS$ and $SJ=NI_k$ give the invariant splitting
$\Q^{kN}=\im J\oplus\ker S$. Its second summand is the rational
new space. Its characteristic polynomial is
$\chi_{A_N}/\chi_B\in\Z[t]$, by monic polynomial division.
If its real spectrum contains one root of $f$, rationality and
irreducibility force $f$ to divide that new characteristic polynomial.
Real symmetry then implies
\begin{equation}\label{eq:lift-real-polynomial-nullity}
 \dim_\Q\ker f(A_N)\ge\ell(t+1).
\end{equation}

Reduction of the integer matrix $f(A_N)$ cannot increase rank,
so the same lower bound holds for its modular nullity at every
prime. At a good prime, polynomial CRT gives
\begin{equation}\label{eq:lift-crt}
 \ker\overline f(A_N\bmod p)=
       \bigoplus_j\ker g_j(A_N\bmod p).
\end{equation}
This holds for any modular matrix: the usual polynomial projectors
sum to identity on the kernel annihilated by $\overline f$ and
project into the indicated kernels. Jordan blocks elsewhere cause
no problem. The old dimensions on the right sum to $\ell t$.
Thus \eqref{eq:lift-real-polynomial-nullity} forces some index $j$
to have dimension greater than $r_jt$, hence at least $r_j(t+1)$
by its divisibility by $r_j$. Equations
\eqref{eq:lift-hom-kernel} and \eqref{eq:lift-shifted-markov} yield
\[
 \limsup_N\Prob\{\spec(A_N|_{W_N})\cap\{f=0\}\ne\varnothing\}
 \le\sum_j\frac1{p^{r_j}-1}\le\frac\ell{p-1}.
\]
Every field, factor, and error limit here is fixed before taking
$N\to\infty$. Summing over the finitely many old factors gives
a collision bound $\deg(F)/(p-1)$. Arbitrarily large fixed good
primes may then tend to infinity, proving
Theorem~\ref{lift:thm:old-new}.

This transfer uses geometric kernels of the integer matrices
$f(A_N)$. Modular characteristic-polynomial multiplicities alone
would not suffice, since a larger Jordan block can preserve geometric
kernel dimension. Nor does the proof split old and new modules
modulo $p$: $J$ injects the old baseline even when $p\mid N$.
The nonbipartite hypothesis is used in the cyclic congruence calculation
and the phase assembly; no bipartite spectral extension is claimed.

\section{Bipartite completion of old/new spectral disjointness}
\label{sec:bipartite-spectral}
The remaining bipartite case requires additional mixed entropy phases
and a two-dimensional excess threshold at zero. We prove the exact
companion moments needed for this case before completing the transfer.

\begin{theorem}[All simple regular bases]\label{lift:thm:all-base}
Let $B$ be the adjacency matrix of any fixed connected simple
$d$-regular base, $d\ge3$, and let $A_N$ be its independent uniform
random $N$-lift. With $U_N$ the real fiber-constant space and
$W_N=U_N^\perp$,
\[
 \Prob\{\spec(A_N|_{W_N})\cap\spec(B)\ne\varnothing\}
 \longrightarrow0
\]
through all positive lift sizes. Every old eigenvalue retains exactly
its base multiplicity, including irrational and repeated old modes.
\end{theorem}

The nonbipartite case is Theorem~\ref{lift:thm:old-new}.
For the rest of this section assume the base is bipartite. Regularity
gives equally many base vertices on its two sides, say $m$ each, and
\begin{equation}\label{eq:bip-square-blocks}
 B=\begin{pmatrix}0&T\\T^T&0\end{pmatrix},\qquad
 A_N=\begin{pmatrix}0&X\\X^T&0\end{pmatrix},
 \qquad T\in\Z^{m\times m},\quad X\in\Z^{mN\times mN}.
\end{equation}
These square blocks are essential to the zero-target argument.
Write the two old parts as $\mathcal L,\mathcal R$.
Fix an odd prime $p$ and a monic irreducible $g\in\Fp[t]$ of
degree $r$. For a companion $R_g$, the evaluation isomorphism
\eqref{eq:lift-hom-kernel} gives
\[
 Z_g(M):=|\{Y:MY=YR_g\}|=|\ker g(M)|,
 \qquad h=\dim_{\Fp}\ker g(B).
\]
This is geometric kernel dimension, even when a modular matrix has
Jordan blocks. All limits below fix $B,p,g$ before $N\to\infty$.
We use the phase, relative local limit and weighted sparse estimates
of Section~\ref{sec:frame-moment-proof}, using the positive-weight
phase specialization in Corollary~\ref{ch:cor:positive-phases}. The new compatible indices are proved below.

\begin{theorem}[Bipartite companion moments]\label{bip:thm:moments}
For every positive lift-size sequence,
\begin{equation}\label{eq:bip-moments}
 \E Z_g(A_N)=c_g p^h+o(1),\qquad
 c_g=\begin{cases}
 p+3,&g=t,\\
 2,&g\ne t\text{ and }g\text{ is not even},\\
 2(p^s+1),&g\text{ is even of degree }2s.
 \end{cases}
\end{equation}
Here even means $g(-t)=g(t)$. No condition on $p\mid N$ or
$p\mid d$ is imposed.
\end{theorem}

\subsection{The full-uniform coefficient and its extra forms}
Use row labels $a\in V=\Fp^{1\times r}$, let $q=p^r$, and
temporarily allow any target $R$. The maximal legal-star congruences
are \eqref{eq:lift-congruences}; before matching compatibility their
index is $p^{krq}$, where $k=2m$. The compatible integer lattice
has the free root-margin and edge-interior parametrization used in
the proof of Lemma~\ref{lift:lem:cyclic}. Its reduction modulo $p$
is the whole compatible direction space, with no saturation factor.

A dependency of the congruences has columns $h_i(a)$.
The same rectangular edge transfers force
\[
 h_i(a)=K_i a^T+c_i,\qquad K_j=-K_i^T\quad(i\sim j).
\]
Connectivity now leaves $K_i=M$ on the left and $K_i=-M^T$ on
the right, with no odd-cycle skew constraint. The independent
root-margin tests at $a$ and $-a$ give
\[
 \sum_{j\sim i}c_j=Rc_i,\qquad
 RM\text{ and }RM^T\text{ are skew-symmetric}.
\]
Conversely these conditions annihilate every compatible direction.
The constant columns form the Frobenius-adjoint old intertwiner
space and have dimension $h_B(R)=\dim\{Y:BY=YR\}$. Put
\[
 b(R)=\dim_{\Fp}\{M:RM\text{ and }RM^T
                                      \text{ are skew-symmetric}\}.
\]
The restricted congruence rank is therefore
$\rho=krq-h_B(R)-b(R)$.

This rank gives an exact coefficient. The relative legal/independent
star density ratio on the logarithmic central window is
$p^{krq}(1+o(1))$. The compatible legal sublattice has index $p^\rho$.
Its Gaussian sum has fraction $p^{-\rho}$ of the independent one,
whose exact weighted normalizer is one by
\eqref{eq:lift-two-normalizers}. The mean, covariance and matching
Stirling factors are the same as in the preceding lift calculation;
the positive precision is \eqref{eq:lift-weighted-precision} and
all powers of $N$ cancel. The all-zero legal stars provide a
compatible affine point for every $N$. Translated lattice sums and
the relative local limit estimate therefore retain this fraction
at every residue, and the fixed-neighborhood tails are $o(1)$.
Thus the full-uniform contribution is exactly
\begin{equation}\label{eq:bip-uniform}
 p^{h_B(R)+b(R)}+o(1).
\end{equation}

If $R$ is invertible, set $L=RM$. The two skew conditions are
equivalent to
\begin{equation}\label{eq:bip-form}
 L^T=-L,\qquad R^2L=L(R^T)^2.
\end{equation}
Indeed $M=R^{-1}L$, and skewness of $RM^T=-RL R^{-T}$ is
equivalent, after multiplication by $R$ and $R^T$, to the second
identity. Hence $b(R)$ counts alternating forms for which right
multiplication by $R^2$ is self-adjoint, including degenerate forms.

For an irreducible $g\ne t$, identify this row module with
$E=\Fp[\lambda]$ and put $K=\Fp[\lambda^2]$.
The extension $E/K$ has degree at most two. Its degree is two
exactly when $g$ is even: a nontrivial automorphism sends
$\lambda$ to $-\lambda$, while an even minimal polynomial supplies
that automorphism. The nonzero constant term forces the degree of
an even $g$ to be $2s$, with $|K|=p^s$.
If $g$ is not even, $K=E$, so $R^2$ is cyclic; the cyclic
alternating-form argument forces $b(R)=0$.

In the even case $E$ has dimension two over $K$.
The forms in \eqref{eq:bip-form} are precisely the $K$-balanced
alternating $\Fp$-forms $\beta(a,b)=aLb^T$, where balance means
$\beta(cx,y)=\beta(x,cy)$ for $c\in K$.
They correspond to alternating $K$-bilinear
forms $\Omega$ through $\beta(x,y)=\operatorname{Tr}_{K/\Fp}\Omega(x,y)$.
To justify this correspondence, the trace polynomial
$z+z^p+\cdots+z^{p^{s-1}}$ is nonzero and has degree below $|K|$;
therefore the pairing $\operatorname{Tr}(cz)$ is nondegenerate.
Define $\Omega(x,y)$ uniquely by
\[
 \operatorname{Tr}(c\Omega(x,y))=\beta(cx,y)\quad(c\in K).
\]
Uniqueness and balance give $K$-linearity in each argument, and
skewness plus odd characteristic gives alternation. Conversely the
trace of any alternating $K$-form has these properties. The space
of alternating forms on a two-dimensional $K$-space has dimension
one over $K$, so $b(R)=s$. This proof also covers $p\mid s$ and
never assumes $\operatorname{Tr}(1)\ne0$.
For $g=t$, the scalar target is zero, both skew conditions are
vacuous, and $b(0)=1$.

\subsection{Nonzero targets and every mixed line phase}
The positive-weight bipartite phase classification
gives subspaces $H_L,H_R$ with
$\lambda H_L\subseteq H_R$ and $\lambda H_R\subseteq H_L$.
For nonzero irreducible targets, invertibility implies
\[
 H_R=\lambda H_L,\qquad \lambda^2H_L=H_L.
\]
If $g$ is not even, irreducibility of multiplication by $\lambda^2$
leaves only the point and full-uniform phases. The point family
has $p^h$ members, each contributing $1+o(1)$ by the sparse bound.
Equation~\eqref{eq:bip-uniform} contributes another $p^h+o(1)$.
This proves the non-even case of \eqref{eq:bip-moments}.

If $g$ is even of degree $2s$, the proper nonzero possibilities are
the $K$-lines
\[
 H_L=uK,\qquad H_R=\lambda uK,
 \qquad \#\{H_L\}=p^s+1.
\]
Coordinates $xu$ on the left and $y\lambda u$ on the right make
the two root maps $I$ and $C$, where $C$ is multiplication by
$c=\lambda^2$ on $K$. This $C$ is cyclic over $\Fp$.
Choose any complementary $K$-line $u'K$ and its right partner
$\lambda u'K$. Multiplication by $\lambda$ preserves the two
pairs of summands. For the base and every lift, the full equation
operator is a direct sum of two identical restricted $(I,C)$
operators. If their old solution dimension is $h_H$, then
\begin{equation}\label{eq:bip-line-splitting}
 h=2h_H.
\end{equation}
The quotient is isomorphic to the complementary restriction. It has
$p^{h_H}$ old point maps, each with an exact old lift through this
fixed complement. Translating that lift removes all affine forcing
at every $N$; an arbitrary affine forcing is not being assumed soluble.

The restricted full-uniform coefficient is exactly $p^{h_H}$.
Indeed its interaction dependencies have slopes $M,-M^T$.
The left map $I$ forces $M^T=-M$; the right map $C$ then forces
$CM=MC^T$. Cyclicity of $C$ kills $M$. The remaining constant
dependencies have dimension $h_H$, by the adjoint of the square
old two-map operator. With $Q_0=|K|$, the legal lattice index
$p^{kQ_0s}$ and compatible index $p^{kQ_0s-h_H}$ give coefficient
$p^{h_H}$ by the same exact normalization and Gaussian fraction.
All affine translates here come from literal exact old sections.

For each quotient point map, count only the restricted full-uniform
edge neighborhood inside its cosets. Its contribution is
$p^{h_H}+o(1)$, so one line pair contributes $p^{2h_H}+o(1)=p^h+o(1)$.
Using the entire restricted moment here would also count its point
subphases. Leakage outside these cosets contributes $o(1)$: the
weighted sparse-boundary bound is $O(N^{-1/2})$ times the bounded
restricted two-space moment. Its matching-forcing rank argument
controls the correlation with a large bulk kernel.
Summing all line pairs, point phases and original full uniformity
gives
\[
 (p^s+1)p^h+p^h+p^{h+s}=2(p^s+1)p^h,
\]
as required.

\subsection{Zero target and singleton-side coefficients}
Put $s_0=\dim\ker T=\dim\ker T^T$, so $h=2s_0$.
For scalar zero there are four subspace pairs:
$(0,0)$, $(\Fp,\Fp)$, $(\Fp,0)$ and $(0,\Fp)$.
The point family contributes $p^h+o(1)$ and full uniformity
contributes $p^{h+1}+o(1)$ by \eqref{eq:bip-uniform}.
For $(\Fp,0)$ the quotient point maps are
$a_R\in\ker T$, numbering $p^{s_0}$. Each has an exact section
with left labels zero and right labels fiber-constant. Its homogeneous
restriction is $X^Tz_L=0$.

We retain the full singleton-side prefactor. Left roots range over
$\Fp$ and right roots are zero. Before compatibility only right
old types have a nontrivial neighbor-sum congruence. These $m$
scalar congruences are independent and are all maximal periods,
by the same $d\ge3$ additive-character argument. Their full index
is $p^m$. Compatible integer directions are arbitrary left root
margins $n_i(a)$ of total zero; all edge tables are those margins
or their transposes. The moments $w_i=\sum_a a n_i(a)$ range
independently over $\Fp$ using label transfers between zero and
one. The restricted congruence map is $T^Tw$, of rank $m-s_0$.
Thus the compatible index is $p^{m-s_0}$.

There are $p^{dmN}$ legal star arrays in the iid normalization,
and the independent-star weighted normalizer after multiplication
by this factor is one. The full and compatible dimensions are
$m(d+1)(p-1)$ and $m(p-1)$; their difference cancels twice the
matching Stirling power. The remaining Gaussian exponent is
$-(p/2)\sum_{i\in\mathcal L}\|u_i\|^2$, with positive precision in every
nonzero direction. The relative density ratio $p^m$, compatible
fraction $p^{-(m-s_0)}$, and exact normalizer therefore give the
restricted uniform coefficient $p^{s_0}$.
The all-zero stars supply compatible affine origins at every $N$,
and the exact old section transports this calculation to each
quotient point map. This mixed family contributes
$p^{s_0}p^{s_0}+o(1)=p^h+o(1)$, as does the opposite family.

In each case choose narrow edge-table neighborhoods of the distinct
uniform-coset phase centers, together with small outside-coset radii.
There are finitely many centers. They can be chosen disjoint from
one another and from point centers. The sparse estimate makes
positive-defect leakage $o(1)$ using the bounded smaller-dimensional
two-space moment. The remaining compact feasible type region has
a strict entropy gap, so its exact factorial sum is exponentially
small. The same argument closes the nonzero cases above.
Adding the four zero families gives $(p+3)p^h+o(1)$, completing
the proof of Theorem~\ref{bip:thm:moments}.

\subsection{Vanishing excess probabilities and real transfer}
Fiber-constant pullback injects every old companion space into the
lift space over every field, even when $p\mid N$. For nonzero
irreducible $g$, both dimensions are multiples of $r=\deg g$.
Subtracting the deterministic baseline $p^h$ in Markov and using
\eqref{eq:bip-moments} gives
\begin{equation}\label{eq:bip-excess}
 \limsup_N\Prob\{\dim\ker g(A_N)>h\}\le
 \begin{cases}
 (p^r-1)^{-1},&g\ne t\text{ not even},\\
 (2p^s+1)/(p^{2s}-1),&g\text{ even of degree }2s,\\
 (p+2)/(p^2-1),&g=t.
 \end{cases}
\end{equation}
For the last denominator, the square blocks in
\eqref{eq:bip-square-blocks} give
$\dim\ker A_N=2\dim\ker X$ over every field. The old nullity
$h=2s_0$ is also even, so every excess is at least two. A
one-dimensional threshold would not give a vanishing zero-target
bound. For $p\ge3$ every bound in \eqref{eq:bip-excess} is at
most $2/(p-1)$.

To complete Theorem~\ref{lift:thm:all-base}, let $F$ be the squarefree
minimal polynomial of the real symmetric integer base $B$, and
fix an odd prime not dividing its nonzero discriminant. Then
$B\bmod p$ is semisimple. For an old irreducible integer factor
$f$ of degree $\ell$ and old exponent $t$, factor
$\overline f=\prod_jg_j$ into distinct irreducibles of degrees $r_j$.
Its old geometric dimensions are $\dim\ker g_j(B)=r_jt$.

A new real root of $f$ forces
$\dim_\Q\ker f(A_N)\ge\ell(t+1)$: the rational new characteristic
polynomial is divisible by $f$, and $A_N$ is real symmetric.
Reduction of the integer matrix $f(A_N)$ cannot increase rank.
Polynomial CRT gives
\[
 \ker\overline f(A_N\bmod p)
       =\bigoplus_j\ker g_j(A_N\bmod p)
\]
for any modular matrix, without semisimplicity of the lift.
Some factor must exceed its injected old baseline. Apply
\eqref{eq:bip-excess}, including its even increment at zero, and
sum over at most $\ell$ factors. This gives limiting collision
probability at most $2\ell/(p-1)$ for $f$, and at most
$2\deg F/(p-1)$ for all old roots. Arbitrarily large fixed good
primes tend to infinity only after the lift-size limit superior.
This proves the bipartite case and hence the all-base theorem.

This extension concerns simple regular bases with equal-size lift
fibers. It does not extend spectral disjointness to arbitrary
reversible degree frames: unequal numbers of old types on the two
bipartite sides can force new zero modes, as in Section~\ref{sec:scope}.
No simplicity of the new spectrum, numerical eigengap, or
growing-parameter uniformity is asserted.

\section{Sharp lift ambiguity and the degree-two boundary}\label{sec:lift-boundary}
The bounds in this section use independent uniform matchings above each
base edge. Statistical claims hide the planted fibers by a uniform
relabeling before revealing the graph.

\subsection{Complete bases attain the structural failure rate}
\begin{theorem}[Complete-base leading constant]\label{thm:complete-lift-sharp}
Fix $d\ge3$, let $k=d+1$, and let $G$ be the random $N$-lift of
$K_k$. For any fixed $K\ge k$, let $\mathcal A_K$ be the event that
an equitable partition into at most $K$ classes does not coarsen the plant.
Then, with $\beta=d-2$ and $c_d=\binom{d+1}{2}$,
\[
 \Prob(\mathcal A_K)=c_dN^{-\beta}+o(N^{-\beta}).
\]
If the known base is $K_k$ and the balanced plant is hidden uniformly,
the optimal recovery Bayes error is
$(c_d/2)N^{-\beta}+o(N^{-\beta})$. Both leading asymptotics persist
after conditioning the total lift connected.
\end{theorem}
\begin{proof}
For $u\in V_i,v\in V_j$, $i\ne j$, let $I_{uv}$ assert adjacency
and equality of their neighbors in every other old fiber. Independence
of the base-edge matchings gives $\Prob(I_{uv})=N^{-d}$.
On this event the closed neighborhoods of $u,v$ agree, so transposing
them is a graph automorphism. Swapping them between their fibers gives
a different balanced cover partition for $N\ge2$. If $X$ counts these
unordered cross-fiber twin pairs, then exactly
\[
 \E X=\binom k2N^{-(d-2)}.
\]
One pair corresponds to a colored motif $F$ consisting of its two
centers, their edge, and $d-1$ common neighbors, one in each remaining
fiber. Thus $v(F)=d+1$, $e(F)=2d-1$, and its excess is $\beta$.

For two distinct feasible motifs with common graph $J$, include all
common vertices, even isolated ones, and put $m=v(J),h=e(J)$.
A fixed overlap pattern has $O_d(N^{2k-m})$ placements and probability
$O_d(N^{-2(2d-1)+h})$, by the exact bounded partial-permutation
probabilities $\prod_{ij}(N)_{r_{ij}}^{-1}$.
Moreover $h\le m$. If the centers coincide, matching injectivity
forces the whole motif to coincide, which is excluded. Otherwise
one center of the first motif is absent or has degree at most two in
the second. Removing it from the first motif leaves a star. Hence
the common graph is a forest with at most one added cycle, proving
$h\le m$. Summing the finitely many overlap types yields
\[
 \E[X(X-1)]=O_d(N^{-2\beta}),\qquad
 \Prob(X\ge1)=c_dN^{-\beta}+O_d(N^{-2\beta}),\quad
 \Prob(X\ge2)=O_d(N^{-2\beta}).
\]

We refine the rigidity proof of Section~\ref{sec:unknown-frame}.
Its macroscopic alternative event has exponentially small probability.
For small bad supports, \eqref{eq:unknown-support} now sums from
size three, giving $O_{d,K}(N^{-3\beta/2})=o(N^{-\beta})$.
A singleton is impossible. A bad support of size two is exactly one
of the twin pairs above. Two vertices of the same old fiber cannot
share a neighbor, by matching injectivity. For vertices in different
fibers that are not adjacent, the neighbor of the first in the second
fiber is exterior and has exactly one support neighbor. If they are
adjacent, badness says precisely that all other neighbors coincide.

Outside the macroscopic event, every candidate $Y$ is within $\eta n$
rows of a point map $JT$. Take $\eta<1/[2k(d+1)]$. Every old fiber
then contains an unchanged closed neighborhood, proving $QT=TR$.
Its nonzero difference $Y-JT$ has a bad row support. Outside the
size-at-least-three event it is a twin pair. Conversely every twin
pair supplies a noncoarsening candidate. These inclusions prove the
ambiguity asymptotic, with error $O_{d,K}(N^{-3\beta/2})$.

For the Bayes assertion, every ordered balanced $K_k$-cover partition
has constant likelihood $(N!)^{-\binom k2}$ on compatible labeled
graphs. Each unordered one has $k!$ orderings, so the posterior is
uniform on its $Z(G)$ compatible balanced cover partitions.
Near a point map, a balanced $k$-class candidate must use an onto
map: an omitted label would force at least $N$ changes. Thus that map
is a permutation. After naming it identically, a size-two difference
preserving all class sizes can only swap the two twin vertices.
Outside the previous exceptional events and $\{X\ge2\}$, $Z=1+X$.
Consequently $\Prob(Z\ge3)=o(N^{-\beta})$ and optimal error
$\E(1-1/Z)\sim \frac12c_dN^{-\beta}$.

Finally every component of a cover meets each fiber in the same
number $m$. A disconnected lift has one with $1\le m\le N/2$.
For specified $m$-subsets in every fiber, all matchings preserve them
with probability $\binom Nm^{-\binom k2}$. Union bounding gives
\[
 \Prob(G\text{ disconnected})
 \le\sum_{m=1}^{\lfloor N/2\rfloor}\binom Nm^{-r_0}
 =O_d(N^{-r_0}),\qquad r_0=\binom k2-k=k(d-2)/2>\beta.
\]
For the displayed order, separate $m=1$ and use
$\binom Nm\ge\binom N2$ for $2\le m\le N/2$.
Conditioning connected therefore preserves both leading asymptotics,
and does not change the uniform posterior on a fixed connected graph.
\end{proof}

The Bayes constant here assumes the base is supplied. It also gives an
information lower bound for any unknown-base class containing this
ensemble. A posterior over unknown bases would require its own prior.
The complete base shows that the universal structural exponent $d-2$
cannot generally be increased.

\subsection{Exact failure of recovery for cycle bases}
\begin{theorem}[Known cycle, hidden fibers]\label{thm:cycle-boundary}
Fix $k\ge3$ and consider the unconditioned random $N$-lift of the
known simple cycle $C_k$, followed by a uniform relabeling. If $R(G)$
is its component count, the posterior on compatible unordered fiber
partitions is uniform on $(2k)^{R(G)-1}$ possibilities. Optimal Bayes
success, and optimal randomized minimax success, are exactly
\begin{equation}\label{eq:cycle-success}
 P_{k,N}=\frac{2k\,\Gamma(N+1/(2k))}
 {\Gamma(1/(2k))\Gamma(N+1)}
 \sim\frac{2k}{\Gamma(1/(2k))}N^{-1+1/(2k)}.
\end{equation}
Thus success tends to zero without any computational restriction.
Connectedness has probability $1/N$; conditional on connectedness,
the supplied-base unordered fiber partition is uniquely recoverable.
\end{theorem}
\begin{proof}
Transport names along a spanning path of the base, making its first
$k-1$ matchings identities. The last matching is a uniform permutation
$\pi\in S_N$. Its cycles of lengths $\ell_1,\ldots,\ell_R$ give graph
components $C_{k\ell_1},\ldots,C_{k\ell_R}$.

A component $C_{k\ell}$ has exactly $2k$ locally bijective maps to
the named base: choose one vertex's image in $k$ ways and a direction
in two ways. Local bijectivity forces that direction thereafter.
Every such map puts $\ell$ vertices in each fiber. Thus all components
give $(2k)^R$ balanced named maps. Two maps have the same unordered
fibers exactly when they differ by a single global automorphism of
$C_k$. Indeed the induced permutation of fiber names must preserve all
base edges; conversely any base automorphism works. This action is
free and has order $2k$. Repeated component lengths create no extra
quotient, since observed vertices are labeled. The partition count is
therefore $(2k)^{R-1}$.

There are $n!/(N!)^k$ balanced named maps. The prior probability of one
is $(N!)^k/n!$; its likelihood for a compatible graph is $(N!)^{-k}$.
Every compatible map thus has joint probability $1/n!$. Summing over
its $2k$ global-name variants gives the same joint probability to each
compatible unordered partition, proving the posterior claim.

For the cycle count $R_N$ of a uniform permutation, inserting the last
element as a singleton or after one of the other $N-1$ elements gives
\[
 \E t^{R_N}=\frac{t(t+1)\cdots(t+N-1)}{N!}
 =\frac{\Gamma(N+t)}{\Gamma(t)\Gamma(N+1)}.
\]
Take $t=1/(2k)$ and multiply by $2k$ for \eqref{eq:cycle-success};
Stirling gives its asymptotic. A permutation is one cycle with
probability $1/N$, and at $R=1$ the partition count is one.

Finally the uniform hidden prior bounds any worst-case success from
above by $P_{k,N}$. Sampling uniformly from the compatible partitions
attains that success for every fixed hidden named map: fix a canonical
phase and direction on one component, and choose the $2k$ possibilities
independently on each other component. Thus it is randomized minimax
optimal. The hidden-prior or minimax formulation is essential; a fixed
publicly named plant can be hardcoded by an estimator.
\end{proof}

Connected conditioning is a different degree-two model. If the cycle
base size is unknown, a connected cycle can cover other cycle bases;
no identification among those bases is asserted by the last theorem.


\section{Two-class and prescribed-cross consequences}\label{sec:results}
Fix integers $a\ge0$, $b\ge1$, and $D=a+b\ge3$, and put
$\delta=a-b$, $n=2N$. There are two planted classes of size $N$;
their independent internal factors are uniform simple $a$-regular,
and the independent cross factor is uniform simple $b$-regular
bipartite. Degree-zero factors are absent. All limits are through
admissible sizes, with $aN$ even and the simple factors existing.
Before hiding the plant write $\sigma=(\one_N,-\one_N)$.
Deterministically,
\begin{equation}\label{eq:planted}
 A\one_n=D\one_n,\qquad A\sigma=\delta\sigma.
\end{equation}
A sign vector $\tau$ satisfies $A\tau=\delta\tau$ exactly when
each vertex has $a$ neighbors of its own sign and $b$ of the other.
It is automatically balanced, since $(D-\delta)\one^T\tau=2b\one^T\tau=0$.

\begin{theorem}[All two-class parameters]\label{thm:main-recovery}
For every fixed $a\ge0,b\ge1,D\ge3$, all valid unordered $(a,b)$
bisections can be listed by an always-correct Las Vegas algorithm in
unconditional expected $\widetilde O_{a,b}(n^2)$ bit work and
expected $\widetilde O_{a,b}(n)$ workspace. Returning any member
recovers the plant with probability $1-O_{a,b}(N^{-\nu})$, where
\[
 \nu=(D-2)m,\qquad
 m=\begin{cases}a-b+1,&a\ge b,\\b-a,&a<b.\end{cases}
\]
Under a uniformly hidden balanced plant this selection attains the
optimal Bayes risk exactly. A deterministic polynomial estimator with
bounded $O_{a,b}(n^3+n^{\nu+2})$ elementary work and $O(n^2)$ storage
attains its leading constant. All admissible sizes are included.
No assortativity, real-corank, or adjacency outlier assumption is needed.
The unknown pair $a,b$ can also be recovered with high probability
in expected nearly quadratic work.
\end{theorem}
\begin{proof}[Algorithmic implications]
Theorem~\ref{thm:frame-bounded-moment} with scalar target $\delta$
bounds the actual-law finite-field kernel moment. Apply
Corollary~\ref{fk:cor:moment}, check the original integer sign
equation, and identify opposite signs. The sharp ambiguity theorem
below gives the recovery rate, uniform posterior and Bayes assertion;
Section~\ref{sec:enumeration} supplies the deterministic cap.
For unknown parameters, read $D$ from the graph and enumerate the
finite family $a\ge0,b\ge1,a+b=D$. Alternatively apply the unknown-frame
algorithm with class bound two. On its rigidity event, every nonempty
two-class candidate is the plant and its integer neighbor counts fix
$a,b$. The finite-family moment theorem preserves the expected work.
\end{proof}

\begin{theorem}[Bounded scalar moment]\label{thm:general-moment}
For every fixed $a\ge0,b\ge1,D\ge3$ and prime $p$, in both the
configuration and independent uniform-simple laws,
\begin{equation}\label{eq:poly-moment}
 \sup_N\E|\ker_{\Fp}(A-\delta I)|\le C_{a,b,p}<\infty.
\end{equation}
The simple-law supremum ranges over sizes where its factors exist.
\end{theorem}
This is the two-equal-class specialization of
Theorem~\ref{thm:frame-bounded-moment}. The Boolean rigidity theorem
below proves failure $O(N^{-\nu})$ and its sharp form computes
the exact leading constant. Neither statement asserts real
eigenspace simplicity.

For the separate prescribed-cross spectral theorem, fix a connected
symmetric $Q=(d_{ij})\in\Z_{\ge0}^{k\times k}$ with $d_{ii}\ge3$.
Its off-diagonal support is connected, with the one-type convention.
There are $k$ equal classes of size $N$, with $d_{ii}N$ even.
Internal blocks are independent uniform simple $d_{ii}$-regular.
Every cross block is an arbitrary prescribed simple $d_{ij}$-regular
bipartite graph independent of all internal randomness. Set
\begin{equation}\label{eq:J}
 Jc=(c_1\one_N,\ldots,c_k\one_N),\qquad W=\im J.
\end{equation}
Then $AJ=JQ$; the total graph is regular precisely when $Q$ has
constant row sum.

\begin{theorem}[Uniform fixed-target exclusion]\label{thm:frame}
For every fixed real $\lambda$, with
$\ell=\dim_\R\ker(Q-\lambda I)$,
\begin{equation}\label{eq:frame-main}
 \sup_{(B_{ij})}\Prob\{
 \ker_\R(A-\lambda I)\ne J\ker_\R(Q-\lambda I)\}\longrightarrow0
\end{equation}
along every admissible sequence. All quotient eigenspaces thus have
exactly their old multiplicities simultaneously with high probability.
Cross blocks may instead have any joint law independent of the
internal blocks.
\end{theorem}

\begin{corollary}[Known-frame rational recovery]\label{cor:frame-recovery}
Given $A,Q$ in Theorem~\ref{thm:frame}, exact rational computation
of $\ker_\Q m_Q(A)$ and grouping its equal rows recovers the planted
classes with high probability, uniformly over the cross blocks.
Here $m_Q$ is the squarefree minimal polynomial of $Q$.
Repeated and irrational quotient eigenvalues are allowed.
\end{corollary}

\begin{corollary}[Two-class prescribed cross graph]\label{cor:two-frame}
For fixed $a\ge3,b\ge1$, and arbitrary prescribed simple $b$-regular
cross block independent of its internal graphs,
\[
 \sup_B\Prob\{\ker_\R(A-(a-b)I)\ne\R\sigma\}\longrightarrow0.
\]
This also holds for $(a,b)=(2,1)$, hence for every $a>b>0$.
Every admissible size is covered. Rational nullspace computation
recovers the plant with high probability. The $a\ge3$ statement
itself requires no assortativity.
\end{corollary}

The prescribed-cross proof uses direct all-target internal configuration
mixing and positive-probability simplicity conditioning, including odd
admissible $N$. The cubic case uses the sharp regular-frame scalar
moment and the cross-matching gauge. This robust quantifier is distinct
from the all-parameter random-cross theorem; it is not asserted for
every prescribed cross graph when $a=0$ or $a=1$.


\section{Exact two-class ambiguity for all parameters}\label{sec:sharp-ambiguity}
\label{sec:rigidity}
Fix integers $a\ge0$, $b\ge1$, and $D=a+b\ge3$. The planted classes
$V_+,V_-$ have size $N$. Their internal factors are independent uniform
simple $a$-regular graphs, and the independent cross factor is uniform
simple $b$-regular bipartite. Degree-zero factors are absent. Limits
run through every admissible size, with $aN$ even and the simple factors
existing. Put $n=2N$, $\delta=a-b$, and
$\sigma=(\one_N,-\one_N)$. Let $K(G)$ count valid unordered
$(a,b)$ bisections, including the plant. Define
\begin{equation}\label{eq:ambiguity-exponent}
 \begin{split}
 m&=\begin{cases}\delta+1,&\delta\ge0,\\-\delta,&\delta<0,\end{cases}
 \qquad \nu=(D-2)m,\\
 C_{a,b}&=\begin{cases}
 \displaystyle\frac{((bt)!)^2(a!)^{2t}}
 {(t!)^2(b!)^{2t}a^{t\delta}},&\delta\ge0,\quad t=\delta+1,\\[6pt]
 \displaystyle\frac{((as)!)^2(b!)^{2s}}
 {(s!)^2(a!)^{2s}b^{s^2}},&\delta<0,\quad s=-\delta.
 \end{cases}
 \end{split}
\end{equation}
The first expression has $a\ge b\ge1$; at $\delta=0$ it has $a=b\ge2$.
The second is well defined at $a=0$ using $0!=1$.

\begin{theorem}[Boolean rigidity]\label{thm:boolean}
For every such fixed parameter pair,
\begin{equation}\label{eq:boolean-rate}
 \Prob\{\exists\tau\ne\pm\sigma:\tau\in\{-1,1\}^{2N},\quad
                         A\tau=\delta\tau\}=O_{a,b}(N^{-\nu}).
\end{equation}
\end{theorem}

\begin{theorem}[Exact ambiguity and Bayes error]\label{thm:sharp-ambiguity}
Along every admissible size sequence,
\begin{equation}\label{eq:sharp-ambiguity}
 \Prob\{K(G)\ge2\}\sim C_{a,b}N^{-\nu},\qquad
 \Prob\{K(G)\ge3\}=O_{a,b}(N^{-\nu-1}).
\end{equation}
Under a uniformly hidden balanced unordered plant, optimal Bayes error
is $R_N\sim(C_{a,b}/2)N^{-\nu}$. An always-correct Las Vegas complete
listing, followed by selection of any member, attains $R_N$ exactly
in unconditional expected $\widetilde O_{a,b}(n^2)$ bit work and
expected $\widetilde O_{a,b}(n)$ workspace. A deterministic estimator
with bounded $O_{a,b}(n^3+n^{\nu+2})$ elementary work and $O(n^2)$
storage attains the same leading Bayes constant.
\end{theorem}

These constants use block size $N$. For example, $(1,2)$ gives
$(\nu,C)=(1,2)$, $(1,3)$ gives $(4,16)$, $a=b\ge2$ gives
$(\nu,C)=(2a-2,(a!)^2)$, and $a=0,b\ge3$ gives
\[
 \nu=b(b-2),\qquad C_{0,b}=\frac{(b!)^{2b-2}}{b^{b^2}}.
\]
The assortative cubic pair $(2,1)$ has $(\nu,C)=(2,4)$.

\subsection{Minimum balanced flips and all larger supports}
Since $A\one=D\one$ and $D-\delta=2b>0$, every integer sign
eigenvector at $\delta$ is balanced. An alternative therefore flips
equally many $j$ vertices in each old class. Identifying global sign
lets us take $j\le N/2$. Write
\[
 w=(\sigma-\tau)/2=\one_{S_+}-\one_{S_-},\qquad
 |S_+|=|S_-|=j,\qquad Aw=\delta w.
\]
At a supported vertex, $\delta=i-c$, where $i$ is its number of
same-old-side support neighbors and $c$ its number of opposite-side
support neighbors. If $\delta\ge0$, simplicity gives
$\delta\le j-1$. Thus $j\ge t=\delta+1$, and equality forces two
$K_t$ internal cores and no support-cross edge. Every support vertex
then has $b$ internal and $b$ cross exterior edges. This includes
$\delta=0$, when both cores are singletons. If $\delta=-s<0$,
then $-s=i-c\ge-j$, so $j\ge s$. Equality forces no internal
support edge and the complete support core $K_{s,s}$; every support
vertex has $a$ internal and $a$ cross exterior edges.
These minimum assertions are used only in the simple graph law.

We bound all $j\ge m+1$ uniformly. At a zero coordinate of $w$,
exactly one incident edge from its nonzero support is impossible.
For a fixed support $S$ of size $h$, explore all incident pairings
in the unconditioned typed configuration model. A collision means
that the partner lies in $S$ or at a previously touched exterior
vertex. If $e$ edges lie in $S$, $B$ leave it, and $z$ exterior
vertices are touched, then
\[
 Dh=2e+B,\qquad z\le B/2,\qquad
 \#\text{collisions}=e+B-z\ge Dh/2.
\]
For $h\le\varepsilon N$, all present pools retain $\Theta_{a,b}(N)$
stubs and only $O_D(h)$ eligible collision stubs have appeared.
Each conditional collision probability is at most $C_{a,b}h/N$.
Absent internal pools are omitted, so no division by $a$ is used at
$a=0$. Degree-one matchings obey the same conditional exposure law.
Padding to $Dh$ steps with dummy noncollisions and union bounding
the collision positions gives
\[
 \Prob\{S\text{ has the necessary exterior condition}\}
 \le2^{Dh}(C_{a,b}h/N)^{Dh/2}.
\]
After summing over $\binom{2N}{h}$ supports this is at most
\begin{equation}\label{ap:eq:support}
 f_h=\left[L_{a,b}(h/N)^\eta\right]^h,
 \qquad \eta=(D-2)/2.
\end{equation}
Its consecutive ratio is at most
$L_{a,b}e^\eta((h+1)/N)^\eta$. Choose a small fixed
$\varepsilon$ making this at most one-half. The sum starting at
$h=2m+2$ is $O(N^{-(D-2)(m+1)})=O(N^{-\nu-1})$.
Conditioning the actual factors simple costs a bounded positive
factor, as justified below.

For macroscopic flips use real one-hot labels, not a sign-variance
argument requiring an internal slot. The candidate quotient is
$\left(\begin{smallmatrix}a&b\\b&a\end{smallmatrix}\right)$.
The exact count \eqref{eq:fmp-exact-types} applies to these integer
local equations. Symmetrization only within each old neighbor type
preserves root and pooled edge laws and increases the exponent.
Let $\Phi$ be the symmetrized exponent per physical vertex, with
old weights $\alpha_+=\alpha_-=1/2$, internal weights $a/4$, and
cross weight $b/2$. The regular frame has the symmetric omission
certificate $q_{ij}=1/D$ on every present factor. The oriented flip
fraction $r=j/N\le1/2$ is exactly the defect from the closest
old-block point map. Thus \eqref{eq:unknown-gap}, with
$c=1-2/D$, gives
\begin{equation}\label{ap:eq:macro-gap}
 2N\Phi\le-c_DNr^2,\qquad
 c_D=\frac{9(D-2)^3}{2D^7}>0.
\end{equation}
This gap is valid for $a=0,1$ and singular candidate quotients,
including $\delta=0$. Polynomially many fixed-alphabet types and
the uniform factorial remainders make every range $r\ge\varepsilon/2$
exponentially small. Combining it with \eqref{ap:eq:support} proves
\begin{equation}\label{ap:eq:large}
 \Prob\{\text{an alternative with least-oriented }j\ge m+1\}
                         =O_{a,b}(N^{-\nu-1}).
\end{equation}
The minimum $2m$ support size here comes from balanced flips; no
corresponding assertion about arbitrary real eigenvectors is used.

\subsection{All minimum patterns and their exact multiplicity}
Put $c=b$ in the nonnegative-$\delta$ case and $c=a$ otherwise.
At an exterior old-plus vertex, $Aw=\delta w$ requires equal
numbers of internal neighbors in $S_+$ and cross neighbors in $S_-$;
denote this common number by $m_v$. The same holds on the other
side, and on each side $\sum_v m_v=cm$. Hence each touched exterior
set has at most $cm$ vertices. A minimum flip is leading when both
have $cm$ vertices, so every positive $m_v=1$.

Its complete support-incident motif $H$ has
\[
 v_0=2m+2cm,\qquad e_0=2e_{\rm int}+e_{\rm cross},
 \qquad e_0-v_0=\nu,
\]
where, per internal factor,
\[
 \begin{array}{c|cc}
 &e_{\rm int}&e_{\rm cross}\\ \hline
 \delta\ge0&t\delta/2+bt&2bt\\
 \delta<0&as&s^2+2as=sD
 \end{array}
\]
A nonleading minimum pattern merges at least one exterior vertex
while preserving all $e_0$ distinct incident edges. There are finitely
many feasible simple typed patterns, each with at most $v_0-1$
vertices. Their placements and prescribed-stub probabilities give
$O(N^{-\nu-1})$, even after bounded-cost simplicity conditioning.
At $c=0$ the exterior is empty and all minimum patterns are leading.

Choose $S_+,S_-$ of size $m$ and disjoint exterior sets $W_+,W_-$
of size $cm$ on their old sides. Four maps specify the boundary:
\[
 W_+\to S_+\text{ internally},\quad W_+\to S_-\text{ across},\quad
 W_-\to S_-\text{ internally},\quad W_-\to S_+\text{ across}.
\]
Each has exactly $c$ preimages at every support vertex. Thus the
exact number of leading data choices is
\begin{equation}\label{ap:eq:data}
 A_N=\binom Nm^2\binom{N-m}{cm}^2
       \left[\frac{(cm)!}{(c!)^m}\right]^4
 \sim N^{v_0}\frac{((cm)!)^2}{(m!)^2(c!)^{4m}}.
\end{equation}
Empty maps count one. A realized leading flip determines its exterior
sets and all maps uniquely. Conversely these edges saturate every
support vertex and balance each exterior contribution, so they give
a valid flip. Other exterior edges do not matter. For $N>2m$ only
one orientation of an alternative has $m$ flips, so there is no
extra orientation or drawing-automorphism factor.

For $c>0$ the number of physical stub assignments for one data choice is
\[
 c_H=(a!b!)^{2m}(ab)^{2cm}.
\]
All support stubs are assigned to distinct prescribed neighbors, so
these assignments are disjoint events. One has exact configuration
probability
\begin{equation}\label{ap:eq:prescribed}
 P_N^*=\prod_{j=0}^{e_{\rm int}-1}(aN-2j-1)^{-2}
                         (bN)_{e_{\rm cross}}^{-1}.
\end{equation}
There is no orientation factor for the ordinary internal pairings.
When $a=0$, omit the internal products and exterior stub factor
directly: $m=b$, $c_H=(b!)^{2b}$ and
$P_N^*=(bN)_{b^2}^{-1}$. No $0^0$ convention is needed.

\subsection{Residual simplicity and the leading constants}
After prescribing one assignment, all remaining factors are independent
uniform configurations. Support residual degrees are zero; exterior
vertices, if present, have internal degree $a-1$ and cross degree
$b-1$; all bulk vertices retain $a,b$. Every prescribed edge touches
a saturated support vertex. Hence a residual edge cannot duplicate
a prescribed edge, and full simplicity is exactly residual simplicity.
All leading data and stub assignments give the same residual lists
up to labels, so their residual simplicity probability $s_N^*$ is
the same finite-$N$ number.

For a residual list let $\gamma$ denote its second factorial moment
sum divided by its stub total. If $\delta\ge0$, those totals and
parameters are
\begin{align*}
 L_a&=aN-tD,&\gamma_a&=(a-1)(1-bt/L_a),\\
 L_b&=bN-2bt,&\gamma_b&=(b-1)(1-bt/L_b).
\end{align*}
If $\delta<0$ and $a\ge1$, they are
\begin{align*}
 L_a&=aN-2as,&\gamma_a&=(a-1)(1-as/L_a),\\
 L_b&=bN-sD,&\gamma_b&=(b-1)(1-as/L_b).
\end{align*}
The internal totals are even: $tD=\delta(\delta+1)+2bt$ in the
first case, and $2as$ is even in the second. The cross totals agree
on both sides and need no evenness assumption. At $a=1$ the internal
residual graph is a perfect matching on $N-2s$ vertices and is always
simple. At $a=0$ the internal factors are absent and the cross
residual list is $b$-regular on $N-b$ vertices per side.

For completeness, bounded-degree configuration defect counts give
internal loop and double-edge means $\gamma/2$ and $\gamma^2/4$,
and cross double-edge mean $\gamma_+\gamma_-/2$. At each fixed
factorial-moment order, disjoint defect prescriptions give the product
limits; overlapping distinct prescriptions have edge excess and lose
a power of $N$. Finite Bonferroni bounds for their zero event give
the simplicity limit. The displayed residual parameters tend to
$a-1,b-1$, just as in the original lists. Thus original and residual
product probabilities have the same positive limit
\begin{equation}\label{ap:eq:simplicity}
 s_N,s_N^*\longrightarrow
 \exp\!\left[-\mathbf1_{a\ge1}\frac{a^2-1}{2}
                         -\frac{(b-1)^2}{2}\right],
 \qquad s_N^*/s_N=1+o(1).
\end{equation}
Absent internal factors contribute one, including at $a=0$.
The degree-one internal factors also contribute one. The ratio is
uniform over all leading data because their residual lists coincide
up to labels. This proves eventual simple completability at all
sufficiently large admissible sizes as well. Conditioning produces
exactly the independent uniform simple factors, since their
distinguished-stub multiplicities are constant.

For one leading data pattern its simple-model probability is exactly
$c_HP_N^*s_N^*/s_N$. When $\delta\ge0$, boundary stub powers cancel
to give
\[
 (1+o(1))(a!b!)^{2t}a^{-t\delta}N^{-e_0}.
\]
When $\delta<0$ and $a\ge1$, they instead give
\[
 (1+o(1))(a!b!)^{2s}b^{-s^2}N^{-e_0}.
\]
At $a=0$ the latter coefficient follows by direct omission of the
empty factors in \eqref{ap:eq:prescribed}. Multiplying by
\eqref{ap:eq:data}, for $X$ counting leading flip sets, proves
\begin{equation}\label{ap:eq:first}
 \mathbb E X\sim C_{a,b}N^{-\nu}
\end{equation}
with precisely the two constants in \eqref{eq:ambiguity-exponent}.
The ratio tending to one, rather than a mere simplicity lower bound,
is required for these coefficients.

\subsection{All overlaps and the multiple-state error}
Every possible leading motif is connected and has minimum degree
at least two. For $\delta>0$ its two clique cores are joined by
exterior two-edge paths. For $\delta=0$ it is $K_{2,2a}$. For
$\delta<0$ its $K_{s,s}$ core is connected, with all exterior vertices
attached. At $a=0$ the whole motif is an isolated $K_{b,b}$.
Support vertices have motif degree $D>2$ and exterior vertices degree
two; at $a=0$ every motif vertex is support.

If $H$ is connected with minimum degree at least two, every proper
subgraph $F$ with nonempty vertex set satisfies
\[
 e(F)-v(F)\le e(H)-v(H)-1.
\]
A spanning proper subgraph omits an edge. Otherwise remove $r>0$
vertices. If $E_R$ edges lie among them and $B$ leave them, then
$B>0$ and $2E_R+B\ge2r$, so $E_R+B>r$. At least $r+1$ edges
are removed; further deletion only lowers excess. This includes
retained isolated vertices.

For two distinct leading occurrences include every common vertex
and edge in their intersection $F$. If nonempty it is proper: an
identical full image would have the same motif-degree-$D$ support
and the same uniquely determined boundary data. Distinct data on one
support cannot co-occur because support stubs are saturated. Their
union therefore has excess at least $\nu+1$. Each of the finitely
many typed overlap patterns contributes $O(N^{v-e})$ by its finite
stub count and bounded-cost simplicity conditioning. Incompatible
patterns contribute zero. Disjoint occurrences contribute
$O(N^{-2\nu})=O(N^{-\nu-1})$, since $\nu\ge1$. In particular the
exponent-one case $(1,2)$ requires no stronger disjoint bound.
Consequently
\begin{equation}\label{ap:eq:second}
 \mathbb E[X(X-1)]=O(N^{-\nu-1}),\qquad
 \Prob(X\ge1)=\mathbb E X+O(N^{-\nu-1}),\qquad
 \Prob(X\ge2)=O(N^{-\nu-1}).
\end{equation}
Outside the larger-flip and nonleading-minimum exceptional events,
every alternative is uniquely a leading minimum flip, so $K(G)=1+X$.
Equations \eqref{ap:eq:large}--\eqref{ap:eq:second} prove
\eqref{eq:sharp-ambiguity}, and hence Theorem~\ref{thm:boolean}.
The multiple-state error does not require a rate for the $o(1)$
simplicity ratio in \eqref{ap:eq:simplicity}.

\subsection{Posterior, exact listing, and a charged polynomial cap}\label{sec:enumeration}
For a fixed valid plant every compatible labeled graph has likelihood
\[
 [G_a(N)^2B_b(N)]^{-1},
\]
independent of the partition, and incompatible graphs have likelihood
zero. The three-factor decomposition is unique once the partition is
specified. Here $G_0(N)=1$ and $G_1(N)$ is the admissible matching
count. A uniform hidden balanced unordered plant therefore has an
exactly uniform posterior on its $K(G)$ valid partitions. Relabeling
invariance identifies its ambiguity probability with the fixed-plant
one. Its optimal conditional error is $1-1/K(G)$, attained by any
verified valid partition, and
\[
 0\le R_N-\tfrac12\Prob(K\ge2)
                  \le\tfrac12\Prob(K\ge3)=O(N^{-\nu-1}).
\]
This proves the Bayes constant $C_{a,b}/2$ without a size-bias or
orientation correction. A publicly supplied deterministic plant
would not support this statistical claim.

Theorem~\ref{thm:frame-bounded-moment} applies to the actual two-class
frame for every parameter pair here. At one fixed odd prime it bounds
the scalar kernel moment at $\delta$, so
Corollary~\ref{fk:cor:moment} gives complete sign listing in
unconditional expected nearly quadratic bit work and nearly linear
workspace. The integer equation rejects modular artifacts, also when
the prime divides $b$; every passing vector is automatically balanced.
Select one representative of each global-sign pair. Returning any
listed partition attains the exact conditional Bayes optimum on every
graph. This proves the Las Vegas assertion of
Theorem~\ref{thm:sharp-ambiguity}.

For the deterministic estimator first check that the input is a
simple $D$-regular graph, returning failure otherwise. Use the bounded
binary moment from the same core theorem and encode signs by
$\tau_i=(-1)^{z_i}$.
Every valid assignment solves
\[
 (A+DI)z=(b\bmod2)\one\quad\text{over }\mathbb F_2,
\]
whose homogeneous matrix is $A-\delta I$ modulo two because
$D+\delta=2a$. Solve this affine system by dense binary elimination.
Cap its cardinality at $n^{\nu+1}$; the core moment and Markov give
interruption probability $O(N^{-\nu-1})$. Otherwise enumerate in
Gray order, checking exact integer neighbor counts, and choose any
valid output. On completion the plant is present in the search.
Elimination costs $O(n^3)$, and generation and checks cost $O_{a,b}(n)$
per candidate, giving bounded $O_{a,b}(n^3+n^{\nu+2})$ elementary
work and $O(n^2)$ storage. On interruption use an arbitrary output.
This raises Bayes risk by at most $O(N^{-\nu-1})=o(N^{-\nu})$ and
therefore preserves its leading constant. It completes the proof of
Theorem~\ref{thm:sharp-ambiguity}.

The sharp constants concern independent random cross factors and the
unconditioned simple-factor law. At $a=0$ the leading ambiguity is an
isolated $K_{b,b}$, so connected conditioning can remove it. No
prescribed-cross analogue of these constants, real-corank asymptotic,
growing-parameter or noisy-recovery conclusion is asserted.

\section{Exact quotient eigenspaces for prescribed cross graphs}\label{sec:frame}

This section proves Theorem~\ref{thm:frame} and
Corollary~\ref{cor:frame-recovery}. Its hypotheses and conclusion
differ from those of the main recovery theorem: all internal degrees
are at least three, all admissible block sizes are allowed, and cross graphs can be
arbitrary prescribed simple biregular graphs. The conclusion identifies
whole real eigenspaces, rather than only their sign vectors.

\subsection{Strong mixing directly in the configuration law}
Initially let each $C_i$ be an independent degree-$d_i=d_{ii}$ uniform
configuration adjacency matrix on $N$ vertices, with loops counted
twice. The only parity conditions are $d_iN$ even. Fix an odd prime
$p$, put $s(x)=\one_N^Tx$, and define
\begin{equation}\label{eq:q-reference}
 q_i(x,r)=\begin{cases}
 \one_{\{r=d_ic\one_N\}},&x=c\one_N,\\
 p^{-(N-1)}\one_{\{s(r)=d_is(x)\}},&x\text{ nonconstant}.
 \end{cases}
\end{equation}
Let $P_i(x,r)=\Prob(C_ix=r)$ and
\begin{equation}\label{eq:mix-error}
 e_i(x)=\sup_r|P_i(x,r)-q_i(x,r)|,\qquad
 \eps_N=\max_i\sum_xe_i(x).
\end{equation}
Theorem~\ref{thm:config-mixing}, proved in
Appendix~\ref{app:configuration-mixing}, gives $\eps_N\to0$ at
fixed $p$ and fixed frame along every admissible size sequence.
The sum of maximum-atom errors is unscaled. Neither $p\nmid d_i$
nor $p\nmid N$ is required. No matching-sum moment is transferred
through contiguity here; we shall condition this actual configuration
moment directly on simplicity.

\begin{lemma}[Replacement for coupled targets]\label{lem:replacement}
For arbitrary target functions $r_i(x_1,\ldots,x_k)$, replacing all
$P_i(x_i,r_i)$ by $q_i(x_i,r_i)$ changes their summed product by at most
\begin{equation}\label{eq:replacement}
 k\eps_N(2p+\eps_N)^{k-1}=o_{p,Q}(1).
\end{equation}
The bound is uniform over every choice of target functions.
\end{lemma}
\begin{proof}
The summed maximum atoms satisfy
\[
 \sum_x\sup_r q_i(x,r)
 =p+(p^N-p)p^{-(N-1)}\le2p,
 \qquad
 \sum_x\sup_r P_i(x,r)\le2p+\eps_N.
\]
Telescope the product. In a term replacing factor $i$, bound its
absolute difference by $e_i(x_i)$ and each other factor by its
maximum atom, which depends only on its own $x_j$. The unrestricted
sum then factorizes. Its bound is
$\eps_N(2p+\eps_N)^{k-1}$ per replacement.
No reference kernel is being treated as a random linear map.
\end{proof}

\subsection{The reference first moment}
For $\theta\in\Fp$ put
$\ell_p=\dim_{\Fp}\ker(Q-\theta I)$. Independence gives
\[
 \E|\ker_{\Fp}(A-\theta I)|
 =\sum_{x_1,\ldots,x_k}
    \prod_iP_i\!\left(x_i,\theta x_i-\sum_{j\ne i}B_{ij}x_j\right).
\]
Use Lemma~\ref{lem:replacement} and split the reference sum according
to which $x_i$ are constant.

If every block is nonconstant, the $i$th reference condition is
\[
 d_{ii}s(x_i)=\theta s(x_i)-\sum_{j\ne i}d_{ij}s(x_j),
 \quad\text{or equivalently}\quad (Q-\theta I)s=0.
\]
For each prescribed sum, there are exactly $p^{N-1}$ vectors, at
most $p$ of which are constant. This remains true when $p\mid N$.
Thus the all-nonconstant contribution is between
$p^{\ell_p}(1-p^{2-N})^k$ and $p^{\ell_p}$, and is
$p^{\ell_p}+o(1)$. If every block is constant, write
$x_i=c_i\one_N$; the equations are exactly $(Q-\theta I)c=0$.
This contribution is exactly $p^{\ell_p}$.

For mixed patterns we use a rank bound which holds in every field.
\begin{lemma}[Disjoint row supports]\label{lem:cross-rank}
If $B$ is a zero-one square matrix of order $N$ with all row and
column sums $d\ge1$, then over every field
\begin{equation}\label{eq:cross-rank}
 \rank B\ge\frac{N}{1+d(d-1)}.
\end{equation}
\end{lemma}
\begin{proof}
Each row support meets at most $d(d-1)$ other row supports. Greedily
select at least $N/[1+d(d-1)]$ rows with pairwise disjoint nonempty
supports. They are linearly independent over every field, since each
has a coordinate equal to one absent from the other selected rows.
\end{proof}

Let $S$ be the nonempty proper set of nonconstant blocks and $T$ its
complement. Connectedness supplies an edge $ij$ with $i\in T$,
$j\in S$, and $d_{ij}>0$. The deterministic equation at constant
block $i$ forces
\begin{equation}\label{eq:mixed-equation}
 \sum_{r\in S}B_{ir}x_r
 =\left((\theta-d_{ii})c_i-
           \sum_{r\in T\setminus\{i\}}d_{ir}c_r\right)\one_N.
\end{equation}
The concatenated coefficient matrix contains $B_{ij}$ and hence has
rank at least $\gamma N$, with
\[
 \gamma=\min_{i<j:d_{ij}>0}[1+d_{ij}(d_{ij}-1)]^{-1}>0.
\]
For fixed constants, \eqref{eq:mixed-equation} permits at most
$p^{|S|N-\gamma N}$ old-vector tuples. Each nonconstant reference
factor is at most $p^{-(N-1)}$ and each constant factor at most one.
The whole pattern contributes at most $p^{k-\gamma N}$, after summing
the $p^{|T|}$ constants. Ignore all other equations, which can only
reduce this count. There are $2^k-2$ mixed patterns; when $k=1$ there
are none.

Combining the three cases proves the sharp configuration-model moment
\begin{equation}\label{eq:frame-moment}
 \E|\ker_{\Fp}(A-\theta I)|
 =2p^{\ell_p}+o_{p,Q}(1),
\end{equation}
uniformly over all permitted cross matrices and over $\theta\in\Fp$.
More explicitly, for $k>1$ the absolute error is at most
\[
 kp^{\ell_p}p^{2-N}+(2^k-2)p^{k-\gamma N}
       +k\eps_N(2p+\eps_N)^{k-1}.
\]
This is a first moment in the configuration model. Below we condition
the nonnegative excess over the deterministic old kernel directly
on simplicity; a qualitative contiguity transfer is unnecessary.

\subsection{Reduction of every fixed real target}
For integer $\lambda$, put $\ell=\nullity_\Q(Q-\lambda I)$.
Outside finitely many primes, quotient rank is preserved: all larger
integer minors vanish, and one nonzero minor of order $k-\ell$
remains nonzero modulo $p$. The old modular kernel has dimension
$\ell$. Real nullity at least $\ell+1$ makes all minors of order
$kN-\ell$ of $A-\lambda I$ vanish, forcing modular nullity at least
$\ell+1$ for every prime. If $K_p$ is the modular nullity, the
deterministic lower bound $K_p\ge\ell$ gives
\begin{equation}\label{eq:subtracted-markov}
 \Prob(K_p\ge\ell+1)
 \le\frac{\E p^{K_p}-p^\ell}{p^{\ell+1}-p^\ell}.
\end{equation}
Equation~\eqref{eq:frame-moment} yields a limiting upper bound
$1/(p-1)$ on real failure, uniformly in the cross matrices.
First let $N\to\infty$ at fixed good odd $p$, then take arbitrarily
large such primes. This proves the integer-target assertion.

We give the elementary algebraic extension to avoid an unnecessary
integer-eigenvalue restriction.
\begin{lemma}[Root primes]\label{lem:root-primes}
For every nonconstant $f\in\Z[t]$ there are arbitrarily large primes
at which $f$ has a root, avoiding any prescribed finite set.
\end{lemma}
\begin{proof}
Choose an integer $c=f(b)\ne0$ and enlarge the forbidden set to
include two and the prime divisors of $c$. Let $M$ be their product.
Expansion gives $f(b+cMt)=c h(t)$, where $h\in\Z[t]$ is nonconstant
and is congruent to one at every forbidden prime. Choose an integer
$t$ with $|h(t)|>1$. A prime divisor of $h(t)$ lies outside that set
and makes $b+cMt$ a residue root of $f$. This is the elementary
Schur prime-divisor argument; it asserts no density or effective rate.
\end{proof}

\begin{lemma}[Good reductions of the quotient]\label{lem:good-roots}
Let $\lambda$ be a real algebraic integer, $f$ its monic irreducible
minimal polynomial, and $\ell=\nullity_\R(Q-\lambda I)$. There is a
nonzero integer $m$ such that for every prime $p\nmid m$ and every
root $\theta\in\Fp$ of $f$, one has
$\nullity_{\Fp}(Q-\theta I)=\ell$.
\end{lemma}
\begin{proof}
Each minor of $Q-tI$ has integer coefficients. Minors of order
$k-\ell+1$ vanish at $\lambda$, hence are divisible by $f$ in
$\Z[t]$. For $\ell=0$ that assertion is vacuous. Choose a minor
$g(t)$ of order $k-\ell$ nonzero at $\lambda$; use the empty minor
$g=1$ when $\ell=k$. Irreducibility gives $\gcd(f,g)=1$ over
$\Q[t]$. Clearing denominators in B\'ezout's identity yields
$Uf+Vg=m\ne0$ with $U,V\in\Z[t]$. At every residue root of $f$
outside the prime divisors of $m$, the chosen minor is nonzero and
all larger minors vanish, proving the claim.
\end{proof}

If $\nullity_\R(A-\lambda I)\ge\ell+1$, every minor of order
$kN-\ell$ of $A-tI$ vanishes at $\lambda$. Its polynomial is
divisible by the monic $f$ in $\Z[t]$, so at every residue root
$\theta$ it vanishes modulo $p$. Hence modular nullity is at least
$\ell+1$. Lemma~\ref{lem:good-roots} identifies the old modular
dimension at the good primes; \eqref{eq:subtracted-markov} applies
with \eqref{eq:frame-moment}. Lemma~\ref{lem:root-primes} supplies
arbitrarily large good odd root primes, completing the same order of
limits. Finally, a real number which is not an algebraic integer
cannot be an eigenvalue of any integer matrix, including $A$ or $Q$.
This proves the fixed-real-target assertion in the configuration model.

\subsection{Direct simple-model moments and uniform quantifiers}
Let $\mathcal S$ be simplicity of all independent internal factors.
Its probability satisfies
\[
 s_N\longrightarrow s_Q:=
 \exp\!\left[-\frac14\sum_i(d_{ii}^2-1)\right]>0.
\]
Conditioning gives exactly independent uniform simple internal
factors. For fixed $p,\theta$, write
$K=|\ker_{\Fp}(A-\theta I)|$ and
$\ell_p=\nullity_{\Fp}(Q-\theta I)$. The deterministic old kernel
has $p^{\ell_p}$ elements. Hence
\begin{equation}\label{eq:frame-simple-moment}
 \E_{\rm simple}K
 \le p^{\ell_p}+
       \frac{\E_{\rm config}K-p^{\ell_p}}{s_N},\qquad
 \limsup_N\sup_B\E_{\rm simple}K
 \le(1+s_Q^{-1})p^{\ell_p}.
\end{equation}
The last inequality follows from \eqref{eq:frame-moment} and is a
raw moment bound in the actual simple law. It does not assert that
the configuration coefficient two remains the simple-model limit.
For the good primes and residue roots in the real-target reduction,
subtract the old contribution in Markov as in
\eqref{eq:subtracted-markov}. The resulting limiting real failure
bound is $s_Q^{-1}/(p-1)$, which tends to zero as arbitrarily large
fixed good primes follow the size limit. This proves
Theorem~\ref{thm:frame} along every admissible sequence.

All configuration estimates were uniform in the prescribed cross
matrices and the conditioning probability is independent of them,
so their supremum is controlled directly. Averaging allows any joint
random cross law independent of the internal factors. This does not
mean simultaneous success for every cross choice on one realization,
or allow cross matrices chosen after observing the internals.

For the cubic pair $(2,1)$, Theorem~\ref{thm:cubic-raw} supplies
\eqref{eq:raw-simple-moment}. Real nullity at least two forces
modular nullity at least two for every prime; subtracting the planted
$p$ vectors gives limiting failure at most $e^{3/2}/(p-1)$.
Large fixed primes after the size limit give the planted real line
at every $N$. Gauge invariance makes uniformity over the prescribed
cross matching exact. Together with the $a\ge3$ case above, this
proves planted-line simplicity for every fixed $a>b>0$, every
admissible size, and every prescribed simple cross matrix independent
of the internal randomness. The argument gives no effective rate.

\subsection{Rational row clustering recovers the whole frame}
Let $\chi_Q$ be the characteristic polynomial and put
\begin{equation}\label{eq:mQ}
 m_Q(t)=\frac{\chi_Q(t)}{\gcd(\chi_Q(t),\chi'_Q(t))},
\end{equation}
with monic gcd. Because $Q$ is symmetric and integral, this is its
minimal polynomial, is squarefree, and belongs to $\Z[t]$.
Its roots are all the distinct real quotient eigenvalues and its
degree is at most $k$.
Always $m_Q(A)J=Jm_Q(Q)=0$. Apply Theorem~\ref{thm:frame} to the
finitely many roots of $m_Q$. Real diagonalizability of $A$ gives,
with probability $1-o(1)$ uniformly in the cross matrices,
\begin{equation}\label{eq:old-polynomial}
 \ker_\R m_Q(A)
 =\bigoplus_{\lambda\in\spec Q}\ker_\R(A-\lambda I)=W.
\end{equation}
Compute a rational basis matrix $V$ for this kernel. On this event
$V=JT$ for an invertible $k\times k$ matrix $T$.
Its rows are equal exactly within each planted class: distinct rows
of $T$ cannot coincide because $T$ is invertible. Grouping equal
rational rows recovers the classes without choosing individual
irrational eigenvectors or splitting repeated eigenspaces.

All coefficients and degrees of $m_Q$ are fixed. If the maximum
total degree of $A$ is $\Delta_Q$ and $m_Q(t)=\sum c_rt^r$, every
absolute row sum of $m_Q(A)$ is at most
$\sum_r|c_r|\Delta_Q^r=O_Q(1)$. Its minors are therefore bounded by
$\exp(O_Q(kN))$ by Hadamard's inequality. Forming this fixed-degree
matrix polynomial and using fraction-free rational elimination have
polynomial bit complexity. This proves
Corollary~\ref{cor:frame-recovery}.

Within the promised frame model, finding exact nullity
$\nullity_\Q m_Q(A)=k$ certifies \eqref{eq:old-polynomial}, because
the old $k$-dimensional space is always included. The algorithm assumes
the frame is known; unknown-frame identifiability is a different problem.

For the homogeneous frame $Q=(a-b)I_k+b\one_k\one_k^T$, with
$a\ge3$, $b\ge1$, and $k>1$, its two distinct eigenvalues are
$a+(k-1)b$ and $a-b$, the latter of multiplicity $k-1$ when $k>1$.
Then $m_Q(t)=(t-a-(k-1)b)(t-a+b)$ for $k>1$. This supplies recovery
for every fixed number of equal classes at all admissible sizes, without an assortativity
or outlier condition. In particular $k=2$ gives the $a\ge3$ part
of Corollary~\ref{cor:two-frame}, without an assortativity condition.

\section{Nearly quadratic two-class listing}\label{sec:two-class-fast}
\subsection{A bounded raw moment in the actual simple graph law}
Fix $a\ge3,b\ge1$ or $(a,b)=(2,1)$, put $D=a+b$,
$\delta=a-b$, and choose one fixed odd prime
$p_0\nmid b$. All limits in the next input keep $p_0$ fixed.
For independent degree-$a$ internal configurations and any prescribed
simple $b$-biregular cross matrix $B$,
Appendix~\ref{app:configuration-mixing} and
Corollary~\ref{thm:cubic-raw} give
\begin{equation}\label{eq:raw-config-moment}
 \E_{\rm config}|\ker_{\F_{p_0}}(A-\delta I)|=2p_0+o(1)
\end{equation}
uniformly over $B$, at every admissible $N$.
The appendix treats $a\ge3$ without requiring assortativity. For the
remaining pair
$(2,1)$, the sharp general-frame scalar moment and invariance under
relabeling give exactly the same law for every prescribed matching.
The internal factors remain actual uniform configurations.

Let $K=|\ker_{\F_{p_0}}(A-\delta I)|$. Each configuration contains
the planted $p_0$ vectors. Simplicity of both internal configurations
has probability $s_N\to\exp[-(a^2-1)/2]>0$. Since $K-p_0\ge0$,
\begin{equation}\label{eq:raw-simple-transfer}
 \E_{\rm simple}K
 =p_0+\frac{\E_{\rm config}[(K-p_0)\one_{\mathcal S}]}{s_N}
 \le p_0+\frac{\E_{\rm config}K-p_0}{s_N}.
\end{equation}
In particular
\begin{equation}\label{eq:raw-simple-moment}
 \limsup_N\sup_B\E_{\rm simple}K
 \le\bigl[1+e^{(a^2-1)/2}\bigr]p_0.
\end{equation}
This is an upper bound, not an asserted simple-model limiting value.
Finitely many initial admissible sizes give a finite uniform supremum
over all sizes. Averaging permits any independent random cross law.
No unbounded expectation is transferred through qualitative contiguity.

\begin{theorem}[Full two-class listing in expected nearly quadratic work]\label{thm:fast-full-list}
For every fixed known $a\ge3,b\ge1$, and also for $(a,b)=(2,1)$,
there is an always-correct Las Vegas
algorithm which lists every equitable $(a,b)$ bisection, up to exchange,
on every simple input graph. In the model with independent uniform simple
internal graphs and any prescribed simple cross graph independent
of the internals, its unconditional expected bit work is
$\widetilde O_{a,b}(n^2)$, its expected workspace is
$\widetilde O_{a,b}(n)$ bits, and its expected output size is
$O_{a,b}(n)$ bits, uniformly over the prescribed cross graph.
All admissible block sizes are covered. No real or modular corank-one
promise and no numerical eigengap is used.
\end{theorem}
\begin{proof}
The input domain is simple graphs. Check $D$-regularity; a graph
failing this check has no equitable $(a,b)$ partition, so return the empty list.
Apply Corollary~\ref{fk:cor:moment} to the actual-law bound
\eqref{eq:raw-simple-moment}. On every input, the full base-field
kernel is certified, every coefficient vector is enumerated, and
integer neighbor counts verify every candidate. Every planted input
has a nonempty list. The field extension used by preconditioning
grows with $n$ only for computational success; the probabilistic
kernel moment is at the fixed base prime $p_0$.
\end{proof}

For a random cross factor, Theorem~\ref{thm:boolean} makes this list
unique with probability $1-O_{a,b}(N^{-\nu})$, where
$\nu$ is defined in \eqref{eq:ambiguity-exponent}.
There is no algorithmic error, so this rate has no separate cutoff or
real-corank failure term. Section~\ref{sec:sharp-ambiguity} proves
the exact leading ambiguity and optimal statistical-error constants.
For arbitrary prescribed cross graphs, the real planted-line theorem
proved above gives uniqueness with qualitative probability $1-o(1)$ uniformly
over $B$, without claiming the same quantitative random-cross rate.

\subsection{Unknown assortative degree pair}
If only the total degree $D$ is supplied, the same expected-work
guarantee can be obtained with high-probability recovery of $a,b$.
At $D=3$ the only positive assortative pair is $(2,1)$ and its target
is $\delta=1$. At $D\ge4$, the actual internal degree is at least
three. Enumerate the finitely many positive targets
$\lambda'=D-2b'$ with $1\le b'<D/2$, and choose a fixed odd prime
$p_0>2D$ for every trial. The generic fixed-target configuration
count in Section~\ref{sec:frame} gives
\[
 \E_{\rm config}|\ker_{\F_{p_0}}(A-\lambda'I)|
 =2p_0^{\ell_{p_0}}+o(1),\qquad
 \ell_{p_0}=\nullity_{\F_{p_0}}(Q-\lambda'I).
\]
Here $Q=\left(\begin{smallmatrix}a&b\\b&a\end{smallmatrix}\right)$
has quotient eigenvalues $D,\delta$. The prime choice preserves all
these integer ranks: $\ell_{p_0}=1$ for $\lambda'=\delta$ and zero
for any other candidate. Conditioning gives a bounded raw moment
at every candidate target. Thus the finite collection of full-kernel
listings still costs expected $\widetilde O_D(n^2)$ bit work.
The fixed-target exclusion theorem, by a finite union, leaves only
the planted target with probability $1-o(1)$ uniformly over $B$.
Exact neighbor counts then recover $a,b$. No claim about unproved
cubic moments at other integer shifts is needed.

\section{Scope and verification}\label{sec:scope}
The model and parameter bounds are part of each theorem. The following
distinctions are necessary for using the results.
\begin{center}
\small
\begin{tabular}{p{0.25\textwidth}p{0.30\textwidth}p{0.34\textwidth}}
Model & Recovery or structure & Work guarantee\\ \hline
Connected reversible frame with symmetric $0\le q_{ij}<1/2$,
$\sum_jd_{ij}q_{ij}=1$
& Unknown frame and finest classes, unequal sizes, varying degrees,
all admissible sizes
& Expected $\widetilde O(n^2)$ complete listing, expected
$\widetilde O(n)$ space; deterministic polynomial cap\\[3pt]
Fixed connected simple regular base lifts, $d\ge3$
& All old eigenvalues retain their exact multiplicities; all $N$
& No numerical gap or new-spectrum simplicity asserted\\[3pt]
Two-class $a\ge0,b\ge1,a+b\ge3$, independent random cross factor
& Exact ambiguity $C_{a,b}N^{-\nu}(1+o(1))$; hidden-prior Bayes risk half
& Expected $\widetilde O(n^2)$ complete listing; deterministic polynomial
estimator attains the leading Bayes constant\\[3pt]
Equal-size known symmetric frame, internal degrees at least three,
prescribed cross factors
& Exact old eigenspaces and rational row clustering, every admissible $N$
& Deterministic polynomial rational computation; no quantitative
real-collision rate asserted\\[3pt]
Known cycle lift, hidden fibers
& Exact optimal success \eqref{eq:cycle-success}, tending to zero
& Information obstruction, irrespective of computation
\end{tabular}
\end{center}

A bound on $\E p^{h_R}$ is a raw moment in the actual graph law.
A high-probability nullity statement or qualitative contiguity does
not transfer such an unbounded quantity. Our fast algorithms average
their per-input cost against the actual moment; they include graphs
with large kernels and do not send exceptional inputs to an uncharged
dense fallback. Their expected bounds are not worst-case quadratic
bounds. Complete partition listing can be long on a particular input.

The unknown-frame theorem searches all nonempty equitable partitions
with at most the supplied constant $K$ classes, without prescribing
their sizes or quotient. It identifies the unique finest one on its
rigidity event. The theorem does not assert uniqueness among all
partitions with an unbounded class count. Unlabeled recovery is the
target; class names and coordinates inside a fiber are not observed
parameters.

Boolean partition uniqueness, real quotient-eigenspace multiplicity,
and simplicity of the new spectrum are different conclusions.
Recovery can even coexist with unavoidable old/new collisions.
For a bipartite reversible regular frame with $k_L,k_R$ old types,
the physical sides have the same size $m$. If its scaled old
off-diagonal block has rank $s$, old left and right nullities are
$k_L-s,k_R-s$. The full off-diagonal block is square and has the
same nullity $z$ on either side, with
$z\ge\max(k_L-s,k_R-s)$. Thus at least
$2z-(k_L-s)-(k_R-s)\ge|k_L-k_R|$ zero modes lie outside the old
space, for every realization. For example,
\[
 Q=\begin{pmatrix}0&3&0\\1&0&2\\0&3&0\end{pmatrix},\qquad
 (N_1,N_2,N_3)=(N,3N,2N)
\]
has old spectrum $\{3,-3,0\}$ and a deterministic new zero mode.
Its random simple factors still satisfy the unknown-frame recovery
theorem. A simple regular base graph has equally many base vertices
on its two bipartite sides, so this particular obstruction does not
apply there.

The exact scalar moment $2p^\ell+o(1)$ is proved only for positive
omission weights and nonbipartite support in the configuration law.
It is not inferred for every zero-weight frame or asserted unchanged
after simplicity conditioning.

The sharp two-class ambiguity asymptotic does not give the same upper
asymptotic for real eigenvalue collisions. An exact algebraic proof
requires no numerical eigengap, but it does not establish numerical
stability. Perturbed or noisy edges need not preserve the equations
$AY=YR$, and none of the recovery claims here concern noisy data.
Likewise the fixed-base results have constants depending on the base,
field, degree and target dimension; no growing-parameter uniformity is
implicit.

The source package supplies a reproducible build, a two-page reading
guide, a source-binding ledger, and exact finite configuration/fiber
checks. Such checks supplement the proofs and do not replace their
asymptotic estimates. The manuscript is a research draft with internal
mathematical audits; it is not a formal verification, external
peer-review verdict, or historical-priority certificate.


% Integration-ready fragment; no existing manuscript source was edited.
% Derived from nonregular-fiber-counting-obstruction.tex, frozen SHA256
% 8b9a8bdfcd9392eb05c4e15e7f3a856e8cab4e95c37b82d2d517af6d2ef30f39.
\subsection{A nonregular base with exponentially ambiguous fibers}
\label{sec:nonregular-obstruction}
Minimum degree three alone cannot replace the omission certificate.
The following obstruction concerns the same known base and equal fiber
sizes, even after revealing all fibers on one of its sides.

\begin{theorem}[Nonregular blind-fiber obstruction]
\label{thm:nonregular-obstruction}
Fix $b>r\ge3$, and let the base be $K_{r,b}$. Reveal its $r$ left
fibers $L_1,\ldots,L_r$, each of size $N$, and hide the unordered
partition of the $bN$ labeled right vertices into its $b$ size-$N$
fibers. Choose that partition uniformly and all base-edge matchings
independently and uniformly. Every estimator, with unlimited work,
has success at most
\begin{equation}\label{eq:nonregular-success}
 s_N:=\frac{b![(bN)!]^{r-1}}
 {(b!)^{rN}(N!)^{(r-1)b}}
 =b!\,b^{(r-1)/2}(2\pi N)^{-(r-1)(b-1)/2}
       e^{-\gamma_{r,b}N}(1+O_{r,b}(N^{-1})),
\end{equation}
where $\gamma_{r,b}=r\log(b!)-(r-1)b\log b$.
If $\gamma_{r,b}>0$, a typical lift has exponentially many compatible
unordered right-fiber partitions, including for every fixed planted
partition. They extend to distinct equal-size equitable partitions
with the same quotient $K_{r,b}$ and with every $L_i$ kept fixed.
For $K_{3,12}$, success is $O(N^{-11}e^{-\gamma_{3,12}N})$, where
$\gamma_{3,12}=3\log(12!)-24\log12>0$.
The asymptotic bounds and high-probability conclusions persist after
conditioning the lift connected; the finite success bound is then
divided by the connection probability.
\end{theorem}
\begin{proof}
Use an ordered hidden right partition first. Its number is
$(bN)!/(N!)^b$. Given one partition, the $rb$ matchings yield distinct
graphs, since each matching is recovered by restriction to its two
fibers. Thus all marked pairs of graph and ordered partition are
equiprobable, and their exact number is
\[
 T_N=(bN)!(N!)^{(r-1)b}.
\]
An observed graph is determined by $r$ maps $f_i:V\to L_i$, each
having $b$ preimages at every left vertex. Consequently its observation
support has size at most
\[
 U_N=\left[\frac{(bN)!}{(b!)^N}\right]^r.
\]
Some such observations need not admit a fiber partition, which only
improves the bound.

For a supported graph let $Z(G)$ count compatible ordered right
partitions and put $W(G)=Z(G)/b!$. Its posterior on the $W(G)$
unordered partitions is uniform. Since $\Prob(G)=Z(G)/T_N$,
\[
 \Prob(\text{successful recovery})
 \le\E\frac1{W(G)}
 =\frac{b!\,|\operatorname{supp}(G)|}{T_N}
 \le\frac{b!U_N}{T_N}=s_N.
\]
Stirling's formula gives \eqref{eq:nonregular-success}. Moreover,
$\Prob(W\le t)\le t s_N$ for $t\ge1$, by applying Markov to $1/W$.
For every $0<c<\gamma_{r,b}$, this yields
\[
 \Prob(W\le e^{cN})
 \le O_{r,b}\!\left(N^{-(r-1)(b-1)/2}
                         e^{-(\gamma_{r,b}-c)N}\right)=o(1).
\]
Relabeling the right vertices acts transitively on their ordered
balanced planted partitions and preserves $W$. Hence this structural
multiplicity bound holds for each fixed plant as well as its uniform
mixture. Distinct right partitions give distinct full partitions with
the same number of nonempty classes, so neither coarsens the other.
The integer inequality $(12!)^3>12^{24}$ certifies the displayed
positive rate for $K_{3,12}$.

Finally, a spanning-tree change of fiber coordinates leaves
$\beta=|E|-|V|+1=(r-1)(b-1)$ independent uniform permutations of
$[N]$. Disconnectedness requires a common invariant subset of size
$1\le j\le N/2$, so
\[
 \Prob(G\text{ disconnected})
 \le\sum_{j=1}^{\lfloor N/2\rfloor}\binom Nj^{1-\beta}
 =O_{r,b}(N^{1-\beta})=o(1).
\]
Conditioning the recovery or multiplicity bound on connectedness
therefore costs only $1+o(1)$.
\end{proof}

The hidden-prior qualification is essential for estimator impossibility;
a supplied canonical plant could simply be printed. The structural
multiplicity conclusion does hold for a fixed plant. This example does
not classify nonregular bases or make $\gamma=0$ a recovery threshold.

There is a close factorization interpretation: each right vertex is an
$r$-partite hyperedge on $L_1,\ldots,L_r$, and a hidden right partition
is a decomposition into $b$ perfect matchings. Frieze and
P\'erez-Gim\'enez~\cite{FriezePerez} prove a recent contiguity theorem
for the corresponding \emph{nonpartite} hypergraph ensembles at the
same factorial entropy inequality, with their rate equal to
$\gamma_{r,b}/r$. The finite partite posterior count above uses no
contiguity theorem. The factorization connection and its standard
planted change of measure are credited; historical priority is not
certified.

% Proposed item for references.tex, whose current version has no such key:
% \bibitem{FriezePerez}
% A.~Frieze and X.~P\'erez-Gim\'enez,
% \emph{Contiguity of unions of perfect matchings and random regular
% uniform hypergraphs}, arXiv:2611.16160v1 (November 12, 2026),
% Theorem 1 and Section 2. \url{https://arxiv.org/html/2611.16160v1}.

\appendix
\section{Strong configuration mixing at every admissible size}\label{app:configuration-mixing}
We supply the configuration-model counterpart needed by the spectral
and raw-moment results. M\'esz\'aros's established strong-mixing
theorem concerns sums of independent perfect matchings~\cite{Meszaros}.
The proof here works directly with uniform stub pairings, extending
the prescribed-output counting interface underlying Huang's
zero-target argument~\cite[Proposition 2.2]{Huang}. Arbitrary targets,
their uniformity, and the central lattice factor are proved below.

Fix $d\ge3$ and an odd prime $p$. Pair $dn$ distinguished stubs
uniformly, $d$ at each vertex, and let $C_n$ be the adjacency matrix
with loops counted twice. Assume only $dn$ even. For $x,r\in\Fp^n$
put $s(x)=\one^Tx$ and
\[
 q_x(r)=\begin{cases}
 \one_{\{r=dc\one\}},&x=c\one,\\
 p^{-n+1}\one_{\{s(r)=ds(x)\}},&x\text{ nonconstant}.
 \end{cases}
\]
\begin{theorem}[All-target, all-parity mixing]\label{thm:config-mixing}
For every fixed $d\ge3$ and fixed odd prime $p$,
\begin{equation}\label{eq:config-mixing}
 \sum_{x\in\Fp^n}\sup_r|\Prob(C_nx=r)-q_x(r)|\longrightarrow0
 \qquad(n\to\infty,\ dn\text{ even}).
\end{equation}
No assumption that $p$ avoids $d$ or $n$ is needed.
\end{theorem}

\subsection{Exact arbitrary-output count}
For a fixed input let $n_a=|\{i:x_i=a\}|$, $\nu_a=n_a/n$.
Let $M_{ab}$ count stubs rooted at label $a$ whose partner has label
$b$. Then $M$ is symmetric with even diagonal and row sums $dn_a$.
Set
\[
 W(M)=\prod_{a<b}M_{ab}!\prod_a(M_{aa}-1)!!,
 \qquad T=(dn-1)!!.
\]
For $t\in\Fp$, let $X^{(t)}$ count the colors in a uniform ordered
$d$-tuple whose field sum is $t$. For each $a$, independently sum
$X^{(r_i)}$ over the vertices with $x_i=a$, obtaining $Z_a$.
Distinguished-stub counting gives exactly
\begin{equation}\label{eq:config-target-count}
 \Prob(C_nx=r)=\frac{p^{(d-1)n}}{T}
       \sum_M W(M)\prod_a\Prob(Z_a=M_a).
\end{equation}
Indeed each vertex has $p^{d-1}$ legal ordered partner-color tuples;
after assigning them to its stubs, compatible pairs number $W(M)$.
Every configuration and output assignment are counted once, including
loops and repeated edges.

For comparison let $Y_a$ count $dn_a$ independent uniform field
labels and define
\[
 B(M)=W(M)\prod_a\frac{(dn_a)!}{\prod_bM_{ab}!}.
\]
Assigning unrestricted partner colors to the fixed labeled stubs and
pairing them gives
\begin{equation}\label{eq:config-baseline}
 \sum_M B(M)=T,\qquad
 W(M)\prod_a\Prob(Y_a=M_a)=p^{-dn}B(M).
\end{equation}
Thus $B(M)/T$ is the exact pairing-type law, rather than a surrogate
independent-table approximation.

\subsection{Uniform nearconstant atoms}
Write $x=c\one+z$ with support $S$ of size $k$; this translation
only shifts the output. Fix a target $r$ for $C_nz$ and put
$U=\{v\notin S:r_v\ne0\}$. If $|U|>dk$ its probability is zero.
Otherwise expose partners of the support stubs. Call a step constrained
when it hits $S\cup U$ or a previously hit exterior vertex.
For $k\le n/4$, each conditional constrained probability is at most
$K_dk/n$, $K_d=2(2d+1)$, since there are at most
$d(2d+1)k$ eligible stubs and at least $dn/2$ available partners.

Let $I$ count support-support edges, $J$ hits in $U$, and $F,R$
first and repeated hits outside $S\cup U$. On the required output,
every touched zero-output exterior vertex has at least two support
edges, so $F\le R$. Hence
$dk=2I+J+F+R$ and $I+J+R\ge dk/2$.
There are fewer than $2^{dk+1}$ possible exposure words. Thus,
uniformly over the entire prescribed target,
\begin{equation}\label{eq:config-near-atom}
 \sup_r\Prob(C_nz=r)
 \le2^{dk+1}(K_dk/n)^{dk/2}
\end{equation}
when $K_dk/n\le1$. No factor for choosing $U$ is needed: the target
is fixed before bounding probability, and the same bound holds for
every target. Sum over at most $p\binom nk(p-1)^k$ inputs. For a
sufficiently small fixed $\eta>0$, the consecutive-term ratio gives
\begin{equation}\label{eq:config-near-total}
 \sum_{\substack{x\text{ nonconstant}\\\max_a\nu_a\ge1-\eta}}
     \sup_r\Prob(C_nx=r)=O_{d,p}(n^{1-d/2})=o(1).
\end{equation}
Loops and arbitrary field cancellations were permitted in this count.

\subsection{All-target entropy tails}
For a probability law $\pi$ on $\Fp$ define
\[
 F_t(\pi)=\sum_{a_1+\cdots+a_d=t}
               \prod_j\pi(a_j)^{(d-1)/d},\qquad F_* =\max_tF_t.
\]
AM--GM bounds each summand by
$d^{-1}\sum_j\prod_{i\ne j}\pi(a_i)$. Summing the constrained
tuples, every omitted-coordinate term sums to one. Therefore
$p^{-(d-1)}\le F_*(\pi)\le1$; the lower bound uses a tuple at a
largest atom. Equality $F_t=1$ requires uniform $\pi$, or a point
mass at $a$ with $t=da$. To see completeness, the support must be
closed under completing every supported $(d-1)$-tuple. Fixing
$d-2$ entries gives reflections preserving the support. Composing
two such reflections gives translation by any difference of support
points. In a prime field two support points force full support;
AM--GM equality then forces all atom probabilities equal.

Write $u=(1/p,\ldots,1/p)$. Pairwise independence of positions in a
uniform tuple of prescribed sum, valid for $d\ge3$, gives uniformly
in $t$
\begin{equation}\label{eq:config-row-Taylor}
 \log F_t(u+z)=-\frac{(d-1)p}{2d}\|z\|_2^2+O_{d,p}(\|z\|_2^3)
 \quad(\one^Tz=0).
\end{equation}
The linear and quadratic cross terms vanish, and bounded third
derivatives in $\|z\|_\infty\le1/(2p)$ make the remainder uniform.

Put $m=M/(dn)$, $\pi_a=m_a/\nu_a$, omitting zero rows. Evaluate
the nonnegative row coefficient generating polynomial at
$z_b=\pi_a(b)^{(d-1)/d}$. With uniform zero-safe factorial errors,
\eqref{eq:config-target-count} yields
\begin{align}
 \Prob(C_nx=r,\text{type }M)
 &\le(n+1)^{C_{d,p}}
             e^{n[-H(\nu)+\Psi(m)]},\label{eq:config-rate-bound}\\
 \Psi(m)&=-\frac{d-2}{2}\KL(m\Vert\nu\nu^T)
             +\sum_a\nu_a\log F_*(\pi_a).\label{eq:config-rate}
\end{align}
Indeed the coefficient bound contributes
$\prod_{a,b}\pi_a(b)^{-(d-1)M_{ab}/d}
\prod_aF_*(\pi_a)^{n_a}$, and
$\log(W(M)/T)=\tfrac{dn}{2}\sum m_{ab}\log m_{ab}+O(\log n)$.
Combining their exponents gives the displayed expression.
Zero coordinates follow by limits. Multiplying by the number of
inputs of empirical law $\nu$, at most $e^{nH(\nu)}$, cancels
the first entropy in \eqref{eq:config-rate-bound}.

The continuous nonpositive rate $\Psi$ has exactly the zeros
$u u^T$ and the $p$ point matrices $e_ae_a^T$.
Its first loss forces $m=\nu\nu^T$ at equality, and the $F_*$
equality cases then classify $\nu$. Near $u u^T$, put
$Q_*=\sum_a\nu_a\|\pi_a-u\|_2^2$.
Symmetry gives $\sum_a\nu_a\pi_a=\nu$, so
\[
 \|\nu-u\|^2\le Q_*,\qquad
 \|m-u u^T\|^2\le(1+1/p)Q_*.
\]
The second bound uses orthogonality of zero-row-sum deviations and
row-constant deviations. Equation~\eqref{eq:config-row-Taylor}
therefore supplies a negative quadratic bound. Away from a fixed
point-mass neighborhood, compactness extends it to
\[
 \Psi(m)\le-c_{d,p,\eta}
       (\|m-u u^T\|^2+\|\nu-u\|^2)
 \quad\text{when }\max_a\nu_a\le1-\eta.
\]
There are polynomially many input profiles and matrices $M$.
Move the output supremum inside their sum, which only enlarges it.
For each $A>0$, choosing a sufficiently large fixed $B$ now gives
\begin{equation}\label{eq:config-middle}
 \sum_{\substack{x:\max\nu_a\le1-\eta\\
          \|\nu-u\|^2>B\log(n+1)/n}}
       \sup_r\Prob(C_nx=r)=O(n^{-A}).
\end{equation}
No enumeration of the exponentially many target vectors occurs.

\subsection{A relative row local limit, uniform over target arrays}
For every prescribed sum $t$, the count vector $X^{(t)}$ of
\eqref{eq:config-target-count} has common mean and covariance
\begin{equation}\label{eq:config-row-cov}
 \mu=(d/p)\one,\qquad\Sigma=(d/p)I-(d/p^2)\one\one^T.
\end{equation}
These equal those of an unconditioned $d$-tuple, since any two
positions are independent uniform labels. Delete count coordinate
zero to obtain positive definite covariance.

Its maximal difference lattice has index $p$ in $\Z^{p-1}$,
given by weighted sum zero modulo $p$. For completeness, a phase
$\sum_j\theta_{a_j}$ constant on all tuples of sum $t$ obeys
$\theta_{a+v}-\theta_a=\theta_b-\theta_{b-v}$: change two entries
and complete their sum with the remaining positions. Hence every
increment depends only on $v$, and
$\theta_b=\theta_0+2\pi\ell b/p$ for some $\ell\in\Fp$.
Characters separate finite-index lattices, proving that there are
no further restrictions. The unconditioned count lattice has index one.

Let $Z=\sum_{i=1}^mX_i^{(t_i)}$ for an arbitrary target sequence.
Uniformly at retained integer coordinates $z$ with
$z-m\mu=O(\sqrt m\log m)$,
\begin{equation}\label{eq:config-row-LLT}
 \Prob(Z=z)=(1+o(1))p\Prob(Y=z)
\end{equation}
on the affine weighted-sum congruence
$\sum_bbz_b=\sum_it_i$, and it is zero off it. Here $Y$ counts
$dm$ independent uniform labels.

We justify the relative error, including the target-array uniformity.
The log moment generating function of the independent array has
derivatives $O(m)$ in a fixed neighborhood, uniformly, because only
the fixed finite family of $p$ laws occurs. The inverse function
theorem supplies the mean-matching tilt
$h=m^{-1}\Sigma^{-1}(z-m\mu)+O(\|z-m\mu\|^2/m^2)$,
of size $O(\log m/\sqrt m)$. Undoing it gives quadratic rate
$-(z-m\mu)^T\Sigma^{-1}(z-m\mu)/(2m)$ with error
$O(\log^3m/\sqrt m)=o(1)$; the unconditioned law has the same
quadratic term and error.

At the tilted mean, Fourier inversion has Gaussian determinant
$(2\pi m)^{-(p-1)/2}(\det\Sigma)^{-1/2}(1+o(1))$.
Every tilted support weight stays uniformly positive. Outside fixed
neighborhoods of the common maximal characters, all characteristic
functions have modulus at most a common $\rho<1$; inside them,
positive covariance gives uniform Gaussian domination. Scaling
frequency by $\sqrt m$ on a logarithmic window makes the third-order
error vanish. The $p$ maximal-character phases agree exactly on
the stated affine congruence and otherwise the probability is zero
by support. The unconditioned law has one character maximum.
This proves \eqref{eq:config-row-LLT} with a uniform relative error,
not merely an absolute-error central limit estimate. It also rules
out central semigroup holes without a normality assumption.

\subsection{Symmetric compatibility cosets and the net factor}
Suppose $n_a=n/p+O(\sqrt{n\log n})$, uniformly in $a$.
Let $\mathcal R_n$ be the type window
$\max_{a,b}|M_{ab}-dn\nu_a\nu_b|\le\sqrt n\log n$.
Dropping $F_*\le1$ from \eqref{eq:config-rate-bound} gives, uniformly
in the target, a bound
$(n+1)^C\exp[-nH(\nu)-(d-2)n\KL(m\Vert\nu\nu^T)/2]$.
Since $nH(\nu)=n\log p-O(\log n)$, Pinsker bounds the total target
mass outside $\mathcal R_n$ by $o(p^{-n})$.
The exact baseline $B(M)/T$ similarly has tail bounded by a polynomial
times $\exp[-dn\KL(m\Vert\nu\nu^T)/2]$, so its mass in this
window is $1-o(1)$. The same holds for windows enlarged or shrunk
by any fixed amount.

The row congruences are
$\Phi_a(M)=\sum_b bM_{ab}=S_a:=\sum_{i:x_i=a}r_i$.
Symmetry and fixed row sums force only
$\sum_a\Phi_a(M)=d\sum_b bn_b$.
To prove central cosets have no additional restriction, use
\[
 Z_j=2(e_j-e_0)(e_j-e_0)^T,\qquad
 \Phi(Z_j)=2j(e_j-e_0),\quad j\ne0.
\]
These bounded integral shifts preserve row sums, symmetry, and
even diagonals. Since $p$ is odd, their images span the total-zero
hyperplane, of size $p^{p-1}$. Bounded combinations realize every
allowed coset difference and preserve nonnegativity in the center.
The exact factorial expression
\[
 B(M)=\frac{\prod_a(dn_a)!}
 {\prod_{a<b}M_{ab}!\prod_a2^{M_{aa}/2}(M_{aa}/2)!}
\]
gives $B(M+Z)/B(M)=1+o(1)$ uniformly for every such shift.
Its leading logarithmic change is
$-\tfrac12\sum_{ab}Z_{ab}\log(dn/p^2)=0$;
the remainder is $O(\log n/\sqrt n)$.
The enlarged/shrunken-window tail estimate handles boundaries.
Thus all $p^{p-1}$ allowed congruence cosets have baseline mass
$p^{-(p-1)}+o(1)$.

Every central row lies in the relative window of
\eqref{eq:config-row-LLT}. In \eqref{eq:config-target-count}, the
$p$ rows give factor $p^p$. Converting their unconditioned masses
by \eqref{eq:config-baseline}, the central contribution at a globally
compatible target is
\[
 (1+o(1))p^{-n+p}
 \Prob_{\rm baseline}(M\in\mathcal R_n,\Phi(M)=S)
 =(1+o(1))p^{-n+1}.
\]
The tail is $o(p^{-n})$. An incompatible global sum is impossible
because $\one^TC_n=d\one^T$. This proves the balanced arbitrary-target
local law uniformly in all targets, with the exact net factor
$p^p/p^{p-1}=p$.

\subsection{Assembly of the strong-mixing sum}
Constant inputs have zero error. Choose the nearconstant radius
$\eta$ so small that
$h(\eta)+\eta\log(p-1)<\log p$, and take the central window
$\|\nu-u\|^2\le B\log(n+1)/n$, with $B$ sufficiently large
for \eqref{eq:config-middle}. These two regions are eventually
disjoint. Near constants the actual maximum atoms sum to $o(1)$
by \eqref{eq:config-near-total}; the reference sum is $o(1)$ because
there are exponentially fewer than $p^n$ such inputs.
In the intermediate region the actual sum is $O(n^{-2})$, and
the reference sum is at most $p$ times the probability that iid
uniform labels miss the central empirical window, which tends to
zero by Hoeffding and a fixed-label union bound.
At the center the per-input error is $o(p^{-n+1})$ uniformly, and
there are at most $p^n$ inputs. Their total error is $o(1)$ with
$p$ fixed. This proves Theorem~\ref{thm:config-mixing}.

All steps used only evenness of the number of stubs $dn$.
Every shift preserved the exact parity constraints. No target-zero
theorem, simple-model mixing assertion, or growing-prime estimate
was substituted for the conclusion just proved.

\begin{thebibliography}{99}
\bibitem{BDGHT}
G.~Brito, I.~Dumitriu, S.~Ganguly, C.~Hoffman, and L.~V.~Tran,
\emph{Recovery and rigidity in a regular stochastic block model},
Proceedings of SODA (2016), 1589--1601.
Version consulted: arXiv:1507.00930v4, Theorems 1--3 and Section 4.
\url{https://arxiv.org/pdf/1507.00930v4}.

\bibitem{Barucca}
P.~Barucca, \emph{Spectral partitioning in equitable graphs},
Physical Review E \textbf{95} (2017), 062310.
\url{https://doi.org/10.1103/PhysRevE.95.062310}.
Author manuscript: \url{https://arxiv.org/abs/1610.02668}.

\bibitem{ZZ}
X.~Zhu and Y.~Zhu, \emph{Non-backtracking eigenvalues and eigenvectors
of random regular graphs and hypergraphs}, arXiv:2312.03300v3,
September 25, 2026. Introduction 1.3, Theorem 3.10, and Lemma 9.1.
\url{https://arxiv.org/html/2312.03300v3}.

\bibitem{Meszaros}
A.~M\'esz\'aros, \emph{The distribution of sandpile groups of random
regular graphs}, Transactions of the American Mathematical Society
\textbf{373} (2020), 6529--6594.
Version consulted: arXiv:1806.03736v4, Theorem 1.6 and preceding definitions.
\url{https://arxiv.org/pdf/1806.03736v4}.
\url{https://doi.org/10.1090/tran/8127}.

\bibitem{NW}
H.~H.~Nguyen and M.~M.~Wood,
\emph{Cokernels of adjacency matrices of random $r$-regular graphs},
arXiv:1806.10068v2, Section 4.
\url{https://arxiv.org/pdf/1806.10068v2}.

\bibitem{JansonSimple}
S.~Janson, \emph{The probability that a random multigraph is simple, II},
arXiv:1307.6344v1, Theorems 1.1, 5.1, and 5.3.
Theorem 5.1 is attributed there to J.~Blanchet and A.~Stauffer.
\url{https://arxiv.org/pdf/1307.6344v1}.

\bibitem{Wiedemann}
D.~H.~Wiedemann, \emph{Solving sparse linear equations over finite fields},
IEEE Transactions on Information Theory \textbf{32} (1986), 54--62.
Section II.
\url{https://www.csd.uwo.ca/~mmorenom/CS424/Ressources/WIEDEMANN-IEEE-1986.pdf}.

\bibitem{Chen}
L.~Chen, W.~Eberly, E.~Kaltofen, B.~D.~Saunders, W.~J.~Turner,
and G.~Villard, \emph{Efficient matrix preconditioners for black box
linear algebra}, Linear Algebra and its Applications \textbf{343--344}
(2002), 119--146. Theorems 4.2 and 4.5.
\url{https://kaltofen.math.ncsu.edu/bibliography/02/CEKSTV02.pdf}.

\bibitem{KP}
S.~Karmakar and M.~Podder, \emph{Uniqueness of communities in regular
stochastic block models}, arXiv:2002.05577v2 (2021), Theorem 3.1.
\url{https://arxiv.org/pdf/2002.05577v2}.

\bibitem{BDH}
G.~Brito, I.~Dumitriu, and K.~D.~Harris,
\emph{Spectral gap in random bipartite biregular graphs and applications},
Combinatorics, Probability and Computing \textbf{31} (2022), 229--267.
Theorem 6.2. \url{https://doi.org/10.1017/S0963548321000249}.

\bibitem{KY}
D.~Kunisky and X.~Yu, \emph{Computational hardness of detecting graph
lifts and certifying lift-monotone properties of random regular graphs},
arXiv:2404.17012v1 (2024), Introduction and Section 2.3.
\url{https://arxiv.org/pdf/2404.17012v1}.

\bibitem{Vontobel}
P.~O.~Vontobel, \emph{Counting in graph covers: a combinatorial
characterization of the Bethe entropy function},
IEEE Transactions on Information Theory \textbf{59} (2013), 6018--6048.
Lemma 29 and Theorems 30, 33.
\url{https://arxiv.org/html/1012.0065v2}.

\bibitem{MoriTanaka}
R.~Mori and T.~Tanaka, \emph{New generalizations of the Bethe approximation
via asymptotic expansion}, arXiv:1210.2592v2 (2012), Section VI-A,
Lemmas 15--16. \url{https://arxiv.org/html/1210.2592v2}.

\bibitem{MoriTanakaCentral}
R.~Mori and T.~Tanaka, \emph{Central approximation in statistical physics
and information theory}, arXiv:1202.0655v1 (2012), Section IV.
\url{https://arxiv.org/html/1202.0655v1}.

\bibitem{DM}
A.~Dembo and T.~McKenzie, \emph{The Gaussian wave for graphs of finite
cone type}, arXiv:2603.04386v1 (2026), Definition 1.6,
Proposition 2.3 and Section 3.
\url{https://arxiv.org/html/2603.04386v1}.

\bibitem{Huang}
J.~Huang, \emph{Invertibility of adjacency matrices for random
$d$-regular graphs}, Duke Mathematical Journal \textbf{170} (2021),
3977--4032. Configuration enumeration and local-limit framework.
\url{https://arxiv.org/pdf/1807.06465v2}.
\url{https://doi.org/10.1215/00127094-2021-0006}.

\bibitem{Rabin}
M.~O.~Rabin, \emph{Probabilistic algorithms in finite fields},
SIAM Journal on Computing \textbf{9} (1980), 273--280.
\url{https://doi.org/10.1137/0209024}.


\bibitem{KidKrypto}
M.~R.~Fellows and N.~Koblitz, \emph{Kid Krypto},
Advances in Cryptology---CRYPTO '92, LNCS \textbf{740} (1993),
371--389, Section 4.5, p.~385.
\url{https://doi.org/10.1007/3-540-48071-4_27}.

\bibitem{BCL}
D.~J.~Bernstein, J.~Cottaar, and T.~Lange,
\emph{IPCC---Improved Perfect Code Cryptosystems}, Chapter 3 in
J.~Cottaar et al., \emph{Report on evaluation of KpqC candidates},
Eindhoven University of Technology, October 31, 2023.
Section 3.1, details of the second attack, p.~17.
\url{https://math.deweger.net/reports/[2023]KpqC.pdf}.

\bibitem{HV}
Y.~Huang and P.~O.~Vontobel,
\emph{Graph-cover-based characterization of the Bethe partition function
of double-edge factor graphs}, arXiv:2506.16250v2 (2026).
Theorem 46, condition (40), and Appendix G, Lemma 57.
\url{https://arxiv.org/html/2506.16250v2}.

\bibitem{HuangThesis}
Y.~Huang, \emph{Finite-graph-cover-based analysis of factor graphs in
classical and quantum information processing systems}, Ph.D. thesis,
The Chinese University of Hong Kong (2024), Theorem 91 and Appendix E.
\url{https://arxiv.org/pdf/2412.05942v1}.

\bibitem{FKMotifs}
J.~Friedman and D.-E.~Kohler,
\emph{The relativized second eigenvalue conjecture of Alon},
arXiv:1403.3462v1 (2014), Theorem 2.4.7.
\url{https://arxiv.org/pdf/1403.3462v1}.

\bibitem{FriezePerez}
A.~Frieze and X.~P\'erez-Gim\'enez,
\emph{Contiguity of unions of perfect matchings and random regular
uniform hypergraphs}, arXiv:2611.16160v1 (November 12, 2026),
Theorem 1 and Section 2.
\url{https://arxiv.org/abs/2611.16160v1}.

\end{thebibliography}

\end{document}
